% VI : rédiger un Curriculum Vitae — Français Tle Tech — Cours signé OnBuch+
% Programme officiel MINESEC. Compiler avec : tectonic N27-vi-rediger-curriculum-vitae.tex
\newcommand{\DOCMATIERE}{Français}
\newcommand{\DOCNIVEAU}{Tle Tech}
\newcommand{\DOCMODULE}{Français – classe de Terminale technique}
\newcommand{\DOCLECON}{N27}
\newcommand{\DOCTITRE}{VI : rédiger un Curriculum Vitae}
% OnBuch+ — préambule « cours signés » ENRICHI (compilé avec tectonic / XeLaTeX)
% Système visuel « Pop » d'OnBuch+ : fond crème (jamais blanc), bordures encre
% épaisses, ombres « dures » décalées, titres Archivo Black en capitales.
% Version MATHÉMATIQUES : graphes (pgfplots/tikz), boîtes pédagogiques
% (définition, propriété, méthode, attention, le savais-tu, exercice résolu,
% exercice, TP, bilan), page de garde et signature OnBuch+ ancrée en bas.
%
% Macros attendues dans main.tex AVANT \input{preamble} :
%   \DOCMATIERE  \DOCNIVEAU  \DOCMODULE  \DOCLECON (numéro)  \DOCTITRE
\documentclass[11pt]{article}

\usepackage{fontspec}
\usepackage[french]{babel}
\usepackage[a4paper,margin=2cm,top=3.4cm,bottom=2.8cm,headheight=1.4cm,footskip=1.3cm]{geometry}
\usepackage{amsmath,amssymb,amsfonts}
\usepackage{mathtools} % \xmapsto, \coloneqq... souvent utilisés par les modèles
\usepackage{cancel} % simplifications barrées (bilans de forces)
% \abs{z} (module, valeur absolue) et \norm{u} (norme), utilisés par les modèles
\providecommand{\abs}{}\let\abs\relax\DeclarePairedDelimiter{\abs}{\lvert}{\rvert}
\providecommand{\norm}{}\let\norm\relax\DeclarePairedDelimiter{\norm}{\lVert}{\rVert}
\usepackage[table]{xcolor} % [table] : fournit aussi \rowcolors (alternance de lignes)
\usepackage{pagecolor}
\usepackage{tikz}
\usetikzlibrary{arrows.meta,positioning,calc,shapes.geometric,shapes.misc,shapes.arrows,decorations.pathmorphing,decorations.pathreplacing,decorations.markings,patterns,shadows,mindmap,trees,fit,backgrounds,matrix,intersections,angles,quotes}
\usepackage{pgfplots}
\usepackage[version=4]{mhchem} % équations nucléaires \ce{^{235}_{92}U}
\usepackage[european,siunitx]{circuitikz} % schémas électriques
\pgfplotsset{compat=1.18}
\usepgfplotslibrary{fillbetween,groupplots} % aires entre courbes (calcul intégral)
\usepackage{enumitem}
\usepackage{fancyhdr}
% [most] charge toutes les bibliothèques tcolorbox (listings, minted, raster,
% documentation...) et gonfle inutilement la mémoire TeX sur un document long
% (~300 boîtes) : on ne charge que ce qui est réellement utilisé.
\usepackage[skins,breakable]{tcolorbox}
\usepackage{graphicx}
\usepackage{colortbl}
\usepackage{booktabs}
\usepackage{tabularx}
\usepackage{multirow}
\usepackage{diagbox} % case d'en-tête diagonale des tableaux à double entrée
% Colonnes extensibles centrée / gauche / droite (C, L, R), souvent utilisées par les modèles
\newcolumntype{C}{>{\centering\arraybackslash}X}
\newcolumntype{L}{>{\raggedright\arraybackslash}X}
\newcolumntype{R}{>{\raggedleft\arraybackslash}X}
\usepackage{array}
\usepackage{multicol}
\usepackage{siunitx}
% parse-numbers=false : accepte aussi \SI{2,5 \times 10^{-3}}{...}
\sisetup{locale=FR,output-decimal-marker={,},group-minimum-digits=4,parse-numbers=false}
\DeclareSIUnit\ohm{\ensuremath{\symup{\Omega}}} % U+2126 absent de la police
\DeclareSIUnit\division{div} % réglages d'oscilloscope (ms/div, V/div)
\DeclareSIUnit\tr{tr} % tours (vitesse de rotation, ex. \SI{1500}{\tr\per\minute})
\DeclareSIUnit\molar{M} % concentration molaire (ex. \SI{3}{\molar}, \SI{1}{\micro\molar})
\DeclareSIUnit\an{an} % année (le modèle confond parfois avec \year, un compteur TeX) — ex. \SI{5}{\kilo\meter\per\an}
\DeclareSIUnit\FCFA{FCFA} % franc CFA (ex. \SI{500}{\FCFA\per\kilo\gram})
\DeclareSIUnit\degreeBrix{\textdegree Brix} % teneur en sucre d'un sirop/jus (réfractométrie)
\DeclareSIUnit\month{mois} % le modèle utilise parfois \month, un compteur TeX, comme unité de durée
\DeclareSIUnit\microliter{\micro\liter} % le modèle écrit parfois \microliter en un seul token
\DeclareSIUnit\million{millions} % dénombrement (ex. \SI{15}{\million\per\mL} spermatozoïdes)
\usepackage{qrcode}
% \degree : commande fréquemment utilisée par les modèles pour un angle ou une
% température en dehors de \SI{}{} (non fournie par les paquets chargés).
\providecommand{\degree}{^{\circ}}
% \qed (fin de démonstration), utilisé par les modèles sans amsthm
\providecommand{\qed}{\hfill$\square$}
% \barre : double barre des valeurs interdites dans un tableau de variations (tkz-tab)
\providecommand{\barre}{\ensuremath{\|}}
% environnement proof (amsthm non chargé) utilisé par les modèles
\ifdefined\proof\else\newenvironment{proof}[1][Démonstration]{\par\noindent\textit{#1.}\ }{\qed\par}\fi
\usepackage{lastpage}
\usepackage{hyperref}
\hypersetup{hidelinks}

% ============================================================
% Palette « Pop » OnBuch+ — mêmes hex que la classe P (plus_ui.dart)
% ============================================================
\definecolor{poporange}{HTML}{F59321}
\definecolor{popdark}{HTML}{E07A0C}
\definecolor{popcream}{HTML}{FFF6E8}
\definecolor{popink}{HTML}{1C1714}
\definecolor{poppurple}{HTML}{7A5AE0}
\definecolor{popgreen}{HTML}{2E9E6A}
\definecolor{poppink}{HTML}{E0457B}
\definecolor{popblue}{HTML}{2F7DE1}
\definecolor{popgold}{HTML}{C98A12}
\definecolor{popmuted}{HTML}{8A7F76}
\definecolor{popline}{HTML}{EADFD2}
\definecolor{popgray}{HTML}{8A7F76}
\colorlet{popgrey}{popgray}
% Alias tolérants (le modèle nomme parfois les couleurs différemment)
\colorlet{popwhite}{white}
\colorlet{popblack}{popink}
\colorlet{popred}{poppink}
\colorlet{popyellow}{popgold}
% Teintes claires (fonds de tableaux/encadrés) — le modèle nomme parfois
% les couleurs avec un suffixe L (ex. popblueL) sans les avoir définies.
\colorlet{poporangeL}{poporange!20}
\colorlet{popdarkL}{popdark!20}
\colorlet{poppurpleL}{poppurple!20}
\colorlet{popgreenL}{popgreen!20}
\colorlet{poppinkL}{poppink!20}
\colorlet{popblueL}{popblue!20}
\colorlet{popgoldL}{popgold!20}
\colorlet{popredL}{popred!20}
\colorlet{popgrayL}{popgray!20}
% Teintes claires pour fonds de tableaux / zones
\colorlet{popblueL}{popblue!10!white}
\colorlet{poporangeL}{poporange!14!white}
\colorlet{popgreenL}{popgreen!12!white}
\colorlet{poppurpleL}{poppurple!12!white}
\colorlet{poppinkL}{poppink!12!white}
\colorlet{popgoldL}{popgold!14!white}

\pagecolor{popcream}

% ============================================================
% Typographie
% ============================================================
\defaultfontfeatures{Ligatures=TeX}
\setmainfont{PlusJakartaSans-Medium.ttf}[Path=../../../fonts/,BoldFont=PlusJakartaSans-Bold.ttf]
% Maths en sans-serif (Fira Math) avec chiffres et lettres droites de Plus
% Jakarta, pour que les formules se fondent dans le texte.
\usepackage[math-style=ISO,bold-style=ISO]{unicode-math}
\setmathfont{FiraMath-Regular.otf}[Path=../../../fonts/]
\setmathfont{PlusJakartaSans-Medium.ttf}[Path=../../../fonts/,range={up/num,up/{Latin,latin},\mathup}]
\usepackage{icomma} % virgule décimale sans espace en mode math
% Caractères mathématiques Unicode absents des polices de texte (Plus Jakarta,
% Archivo Black) : rendus par leur équivalent mathématique, en texte comme en
% formule — sinon ils s'affichent en carré vide (ex. « Arithmétique dans ℤ »).
\usepackage{newunicodechar}
\newunicodechar{ℝ}{\ensuremath{\mathbb{R}}}
\newunicodechar{ℤ}{\ensuremath{\mathbb{Z}}}
\newunicodechar{ℂ}{\ensuremath{\mathbb{C}}}
\newunicodechar{ℕ}{\ensuremath{\mathbb{N}}}
\newunicodechar{ℚ}{\ensuremath{\mathbb{Q}}}
\newunicodechar{𝔻}{\ensuremath{\mathbb{D}}}
\newunicodechar{α}{\ensuremath{\alpha}}
\newunicodechar{β}{\ensuremath{\beta}}
\newunicodechar{γ}{\ensuremath{\gamma}}
\newunicodechar{δ}{\ensuremath{\delta}}
\newunicodechar{ε}{\ensuremath{\varepsilon}}
\newunicodechar{θ}{\ensuremath{\theta}}
\newunicodechar{λ}{\ensuremath{\lambda}}
\newunicodechar{μ}{\ensuremath{\mu}}
\newunicodechar{σ}{\ensuremath{\sigma}}
\newunicodechar{τ}{\ensuremath{\tau}}
\newunicodechar{φ}{\ensuremath{\varphi}}
\newunicodechar{ψ}{\ensuremath{\psi}}
\newunicodechar{ω}{\ensuremath{\omega}}
\newunicodechar{Σ}{\ensuremath{\Sigma}}
% majuscules produites par \MakeUppercase dans les titres (α-aminés -> Α)
\newunicodechar{Α}{\ensuremath{\alpha}}
\newunicodechar{Β}{\ensuremath{\beta}}
\newunicodechar{Γ}{\ensuremath{\Gamma}}
\newunicodechar{Δ}{\ensuremath{\Delta}}
\newunicodechar{Ε}{\ensuremath{\varepsilon}}
\newunicodechar{Θ}{\ensuremath{\Theta}}
\newunicodechar{Λ}{\ensuremath{\Lambda}}
\newunicodechar{Μ}{\ensuremath{\mu}}
\newunicodechar{Τ}{\ensuremath{\tau}}
\newunicodechar{Φ}{\ensuremath{\Phi}}
\newunicodechar{Ψ}{\ensuremath{\Psi}}
\newunicodechar{Ω}{\ensuremath{\Omega}}
\newunicodechar{↦}{\ensuremath{\mapsto}}
\newunicodechar{∈}{\ensuremath{\in}}
\newunicodechar{∉}{\ensuremath{\notin}}
\newunicodechar{∧}{\ensuremath{\wedge}}
\newunicodechar{⇔}{\ensuremath{\Leftrightarrow}}
\newunicodechar{⟺}{\ensuremath{\Longleftrightarrow}}
\newunicodechar{⇒}{\ensuremath{\Rightarrow}}
\newunicodechar{⟹}{\ensuremath{\Longrightarrow}}
\newunicodechar{∘}{\ensuremath{\circ}}
\newunicodechar{∪}{\ensuremath{\cup}}
\newunicodechar{∩}{\ensuremath{\cap}}
\newunicodechar{≡}{\ensuremath{\equiv}}
\newunicodechar{⊂}{\ensuremath{\subset}}
\newunicodechar{⊥}{\ensuremath{\perp}}
\newunicodechar{∃}{\ensuremath{\exists}}
\newunicodechar{∀}{\ensuremath{\forall}}
\newunicodechar{∛}{\ensuremath{\sqrt[3]{\,}}}
\newunicodechar{∜}{\ensuremath{\sqrt[4]{\,}}}
\newunicodechar{⃗}{\ensuremath{{}^{\to}}}
\newunicodechar{ⁿ}{\ifmmode^{n}\else\textsuperscript{\ensuremath{n}}\fi}
\newunicodechar{ˣ}{\ifmmode^{x}\else\textsuperscript{\ensuremath{x}}\fi}
\newunicodechar{ᵇ}{\ifmmode^{b}\else\textsuperscript{\ensuremath{b}}\fi}
\newunicodechar{ʳ}{\ifmmode^{r}\else\textsuperscript{\ensuremath{r}}\fi}
\newunicodechar{ᵘ}{\ifmmode^{u}\else\textsuperscript{\ensuremath{u}}\fi}
\newunicodechar{ʸ}{\ifmmode^{y}\else\textsuperscript{\ensuremath{y}}\fi}
\newunicodechar{ᵃ}{\ifmmode^{a}\else\textsuperscript{\ensuremath{a}}\fi}
\newunicodechar{ᵉ}{\ifmmode^{e}\else\textsuperscript{\ensuremath{e}}\fi}
\newunicodechar{ᵏ}{\ifmmode^{k}\else\textsuperscript{\ensuremath{k}}\fi}
\newunicodechar{ᵖ}{\ifmmode^{p}\else\textsuperscript{\ensuremath{p}}\fi}
\newunicodechar{ᴺ}{\ifmmode^{N}\else\textsuperscript{\ensuremath{N}}\fi}
\newunicodechar{ᶜ}{\ifmmode^{c}\else\textsuperscript{\ensuremath{c}}\fi}
\newunicodechar{ʰ}{\ifmmode^{h}\else\textsuperscript{\ensuremath{h}}\fi}
\newunicodechar{ᵠ}{\ifmmode^{q}\else\textsuperscript{\ensuremath{q}}\fi}
\newunicodechar{ₙ}{\ifmmode_{n}\else\textsubscript{\ensuremath{n}}\fi}
\newunicodechar{ₐ}{\ifmmode_{a}\else\textsubscript{\ensuremath{a}}\fi}
\newunicodechar{ᵢ}{\ifmmode_{i}\else\textsubscript{\ensuremath{i}}\fi}
\newunicodechar{ᵣ}{\ifmmode_{r}\else\textsubscript{\ensuremath{r}}\fi}
\newunicodechar{ₖ}{\ifmmode_{k}\else\textsubscript{\ensuremath{k}}\fi}
\newunicodechar{ᵦ}{\ifmmode_{\beta}\else\textsubscript{\ensuremath{\beta}}\fi}
\newunicodechar{ₓ}{\ifmmode_{x}\else\textsubscript{\ensuremath{x}}\fi}
\newfontfamily\popdisplay{ArchivoBlack-Regular.ttf}[Path=../../../fonts/]
\newfontfamily\poplogo{SpaceGrotesk-Bold.ttf}[Path=../../../fonts/]
\newfontfamily\popsemi{PlusJakartaSans-SemiBold.ttf}[Path=../../../fonts/]
\newfontfamily\popxbold{PlusJakartaSans-ExtraBold.ttf}[Path=../../../fonts/]
\newfontfamily\popxboldls{PlusJakartaSans-ExtraBold.ttf}[Path=../../../fonts/,LetterSpace=3.0]

\setlist[enumerate]{leftmargin=1.7em,itemsep=5pt,topsep=4pt}
\setlist[itemize]{leftmargin=1.7em,itemsep=4pt,topsep=4pt,label={\color{poporange}\textbullet}}
\setlength{\parindent}{0pt}
\setlength{\parskip}{5pt}
\renewcommand{\arraystretch}{1.25}

% pgfplots : style Pop par défaut
\pgfplotsset{
  every axis/.append style={axis line style={popink,line width=1pt},tick style={popink},
    grid style={popline,line width=0.5pt},label style={font=\small},tick label style={font=\footnotesize}},
  popcurve/.style={poporange,line width=1.8pt,no markers,smooth},
}
% Même style disponible dans un \draw TikZ simple (hors pgfplots)
\tikzset{popcurve/.style={poporange,line width=1.8pt}}
\tikzset{popgrid/.style={popline,very thin}}
\colorlet{popcurve}{poporange} % le nom du style est parfois utilisé comme couleur

% ============================================================
% Pill / squiggle
% ============================================================
\newcommand{\poppill}[3]{%
  \tikz[baseline=-0.55ex]\node[draw=popink,line width=1.2pt,rounded corners=2.6pt,%
    fill=#2,inner xsep=6.5pt,inner ysep=3pt]{%
    \popxboldls\fontsize{7}{7}\selectfont\color{#3}\MakeUppercase{#1}};%
}
\newcommand{\popsquiggle}[1]{%
  \tikz[baseline=-0.4ex]{
    \draw[#1,line width=2.6pt,line cap=round]
      plot [smooth] coordinates {(0,0.09) (0.34,0.01) (0.68,0.21) (1.02,0.01) (1.36,0.10)};
  }%
}
% Mot-clé surligné (marqueur orange sous le texte)
\newcommand{\cle}[1]{{\popxbold\color{popdark}#1}}

% ============================================================
% En-tête / pied
% ============================================================
\pagestyle{fancy}
\fancyhf{}
\renewcommand{\headrulewidth}{0pt}
\renewcommand{\footrulewidth}{0pt}
\lhead{\raisebox{-0.15ex}{\poplogo\fontsize{13}{13}\selectfont\color{popink}OnBuch\color{poporange}+}}
\rhead{\poppill{\DOCMATIERE\ \textbullet\ \DOCNIVEAU\ \textbullet\ Leçon \DOCLECON}{white}{popink}}
\lfoot{\popxbold\fontsize{7.6}{7.6}\selectfont\color{popmuted}\MakeUppercase{Cours signé OnBuch+}\ \textbullet\ {\popsemi\DOCTITRE}}
\rfoot{\tikz[baseline=-0.6ex]\node[rounded corners=8.5pt,fill=popink,minimum height=17pt,minimum width=17pt,inner xsep=5pt,inner sep=0]{\popxbold\fontsize{8}{8}\selectfont\color{popcream}\thepage\,/\,\pageref*{LastPage}};}

% ============================================================
% Titres
% ============================================================
\newcommand{\chaptitre}[2]{%
  \begin{tcolorbox}[enhanced,colback=poporange,colframe=popink,boxrule=2.6pt,arc=11pt,
      drop shadow=popink,left=15pt,right=15pt,top=12pt,bottom=11pt,before skip=3pt,after skip=18pt]
    \poppill{#2}{popink}{popcream}\par\vspace{9pt}
    {\raggedright\hyphenpenalty=10000\exhyphenpenalty=10000\popdisplay\fontsize{22}{24}\selectfont\color{popcream}\MakeUppercase{#1}\par}
  \end{tcolorbox}
}
% Section numérotée : pastille orange avec le numéro + titre Archivo
\newcounter{popsec}
\newcommand{\coursec}[1]{%
  \stepcounter{popsec}\setcounter{popsub}{0}%
  \par\vspace{16pt}\noindent
  \begin{minipage}{\linewidth}
  \tikz[baseline=-0.7ex]\node[draw=popink,line width=1.6pt,fill=poporange,rounded corners=4pt,minimum size=22pt,inner sep=0]{\popdisplay\fontsize{12}{12}\selectfont\color{popcream}\thepopsec};%
  \hspace{8pt}{\popdisplay\fontsize{15}{17}\selectfont\color{popink}\MakeUppercase{#1}}\par
  \vspace{4pt}\noindent{\color{poporange}\rule{36pt}{3.6pt}}
  \end{minipage}\par\nopagebreak\vspace{8pt}
}
\newcounter{popsub}
\newcommand{\courssub}[1]{%
  \stepcounter{popsub}%
  \par\vspace{9pt}\noindent{\popxbold\fontsize{11.5}{13}\selectfont\color{popink}{\color{poporange}\thepopsec.\thepopsub}\hspace{6pt}#1}\par\nopagebreak\vspace{4pt}
}

% ============================================================
% Boîtes pédagogiques — même recette : fond blanc, bordure épaisse colorée,
% ombre dure encre, badge plein. Titre optionnel [#1].
% ============================================================
\tcbset{popbase/.style={breakable,enhanced,colback=white,boxrule=2.3pt,arc=9pt,
  shadow={3pt}{-3pt}{0pt}{popink},left=14pt,right=14pt,top=11pt,bottom=11pt,before skip=11pt,after skip=15pt,
  fontupper=\popsemi\fontsize{10.6}{14.5}\selectfont\color{popink},
  fontlower=\popsemi\fontsize{10.6}{14.5}\selectfont\color{popink}}}
% Coche dessinée (le glyphe ✓ n'existe pas dans les polices du document)
\newcommand{\popcheck}[1]{\tikz[baseline=-0.5ex]\draw[#1,line width=1.6pt,line cap=round,line join=round](-0.13,0.0)--(-0.04,-0.09)--(0.13,0.1);}
\providecommand{\coche}{\popcheck{popgreen}}
\providecommand{\resultat}[1]{\par\smallskip\noindent\fcolorbox{popgreen}{popgreenL}{\parbox{\dimexpr\linewidth-2\fboxsep-2\fboxrule\relax}{\textbf{Résultat.} #1}}\par\smallskip}
\newcommand{\popboxhead}[3]{\poppill{#1}{#2}{white}\ifx\relax#3\relax\else\hspace{6pt}{\popxbold\fontsize{10.6}{12}\selectfont\color{popink}#3}\fi\par\vspace{7pt}}

\newtcolorbox{aretenir}[1][]{popbase,colframe=popgreen,before upper={\popboxhead{À retenir}{popgreen}{#1}}}
\newtcolorbox{exemplebox}[1][]{popbase,colframe=poporange,before upper={\popboxhead{Exemple}{poporange}{#1}}}
\newtcolorbox{definition}[1][]{popbase,colframe=popblue,before upper={\popboxhead{Définition}{popblue}{#1}}}
\newtcolorbox{propriete}[1][]{popbase,colframe=poppurple,before upper={\popboxhead{Propriété}{poppurple}{#1}}}
\newtcolorbox{methode}[1][]{popbase,colframe=popgold,before upper={\popboxhead{Méthode}{popgold}{#1}}}
\newtcolorbox{attention}[1][]{popbase,colframe=poppink,before upper={\popboxhead{Attention}{poppink}{#1}}}
\newtcolorbox{experience}[1][]{popbase,colframe=popblue,colback=popblueL,before upper={\popboxhead{Expérience}{popblue}{#1}}}
% « Le savais-tu ? » : carte violette pleine (contexte, culture, Cameroun)
\newtcolorbox{savaistu}[1][]{popbase,colback=poppurple,colframe=popink,
  fontupper=\popsemi\fontsize{10.4}{14.2}\selectfont\color{white},
  before upper={\poppill{Le savais-tu\,?}{white}{poppurple}\ifx\relax#1\relax\else\hspace{6pt}{\popxbold\color{white}#1}\fi\par\vspace{7pt}}}
% Exercice résolu : énoncé en haut, \tcblower puis la solution sur fond crème
\newcounter{popexo}\newcounter{popexr}
\newtcolorbox{exoresolu}[1][]{popbase,colframe=poporange,colbacklower=popcream,
  segmentation style={popink,line width=1.2pt,dashed},
  before upper={\stepcounter{popexr}\popboxhead{Exercice résolu \thepopexr}{poporange}{#1}},
  before lower={\poppill{Solution}{popink}{popcream}\par\vspace{6pt}}}
% Exercice (non résolu en place) — la correction est dans la partie Corrigés
\newtcolorbox{exercice}[1][]{popbase,colframe=popink,
  before upper={\stepcounter{popexo}\popboxhead{Exercice \thepopexo}{popink}{#1}}}
% Corrigé
\newtcolorbox{corrige}[1][]{popbase,colframe=popgreen,colback=popgreenL,
  before upper={\popboxhead{Corrigé}{popgreen}{#1}}}
% Objectifs / prérequis / bilan
\newtcolorbox{objectifs}{popbase,colframe=popink,colback=poporangeL,
  before upper={\popboxhead{Objectifs de la leçon}{popink}{}}}
\newtcolorbox{prerequis}{popbase,colframe=popink,colback=popblueL,
  before upper={\popboxhead{Prérequis}{popblue}{}}}
\newtcolorbox{bilan}{popbase,colframe=popink,colback=white,boxrule=2.8pt,
  before upper={\popboxhead{Fiche bilan}{poporange}{L'essentiel à connaître pour l'examen}}}
% Figure encadrée avec légende
\newtcolorbox{popfigure}[1][]{enhanced,breakable=false,colback=white,colframe=popline,boxrule=1.4pt,arc=8pt,
  left=10pt,right=10pt,top=10pt,bottom=8pt,before skip=10pt,after skip=12pt,halign=center,
  lower separated=false,halign lower=center,
  fontlower=\popsemi\fontsize{9}{11.5}\selectfont\color{popmuted},
  #1}
\newcommand{\legende}[1]{\par\vspace{6pt}{\popsemi\fontsize{9}{11.5}\selectfont\color{popmuted}#1}}

% ============================================================
% Signature OnBuch+
% ============================================================
\newcommand{\poplogoinline}[1]{{\poplogo\fontsize{#1}{#1}\selectfont\color{popink}OnBuch\color{poporange}+}}
\newcommand{\signatureonbuch}{%
  \begin{tcolorbox}[enhanced,colback=popink,colframe=popink,boxrule=2.2pt,arc=10pt,
      drop shadow=poporange,left=14pt,right=14pt,top=10pt,bottom=10pt,before skip=0pt,after skip=0pt]
    \tikz[baseline=-0.7ex]\node[circle,fill=popgreen,draw=popcream,line width=1.3pt,minimum size=22pt,inner sep=0]{\popcheck{white}};%
    \hspace{9pt}\begin{minipage}[c]{0.62\linewidth}
      {\popxbold\fontsize{12}{13}\selectfont\color{popcream}Signé\ \textbullet\ {\poplogo OnBuch\color{poporange}+}}\par\vspace{2pt}
      {\raggedright\popsemi\fontsize{8}{10.5}\selectfont\color{popline}\MakeUppercase{Équipe pédagogique OnBuch+}\\\MakeUppercase{Conforme au programme officiel MINESEC}\par}
    \end{minipage}\hfill
    {\poplogo\fontsize{22}{22}\selectfont\color{popcream}OnBuch\color{poporange}+}
  \end{tcolorbox}%
}

% ============================================================
% Page de garde
% #1 titre · #2 matière · #3 niveau · #4 module · #5 n° leçon · #6 durée ·
% #7 description / objectifs (texte ou itemize)
% ============================================================
\newcommand{\pagegarde}[7]{%
  \thispagestyle{empty}
  \newgeometry{a4paper,margin=1.7cm,top=1.6cm,bottom=1.4cm}%
  \begin{tikzpicture}[remember picture,overlay]
    \fill[poporange!18] ([xshift=-3.2cm,yshift=-3.6cm]current page.north east) circle (4.6cm);
    \fill[poppurple!14] ([xshift=2.4cm,yshift=7.6cm]current page.south west) circle (3.8cm);
  \end{tikzpicture}%
  {\poplogo\fontsize{30}{30}\selectfont\color{popink}OnBuch\color{poporange}+}\hfill
  \raisebox{0.6ex}{\poppill{Cours signé \textbullet\ Premium}{popink}{popcream}}\par
  \vspace{1pt}\popsquiggle{poppurple}\par
  \vspace{8pt}
  \begin{tcolorbox}[enhanced,colback=poporange,colframe=popink,boxrule=2.8pt,arc=13pt,
      drop shadow=popink,left=18pt,right=18pt,top=12pt,bottom=13pt,after skip=10pt]
    \poppill{#2 \textbullet\ #3}{popink}{popcream}\hspace{5pt}\poppill{#4}{white}{popink}\par\vspace{8pt}
    {\popdisplay\fontsize{14}{14}\selectfont\color{popink}LEÇON #5}\par\vspace{4pt}
    {\raggedright\hyphenpenalty=10000\exhyphenpenalty=10000\popdisplay\fontsize{25}{27}\selectfont\color{popcream}\MakeUppercase{#1}\par}
    \vspace{8pt}
    {\popxbold\fontsize{9.5}{9.5}\selectfont\color{popink}\MakeUppercase{Durée conseillée : #6}}
  \end{tcolorbox}
  \begin{tcolorbox}[enhanced,colback=white,colframe=popink,boxrule=2.2pt,arc=10pt,
      drop shadow=popink,left=16pt,right=16pt,top=10pt,bottom=10pt,after skip=8pt]
    \poppill{Dans cette leçon}{poppurple}{white}\par\vspace{5pt}
    \popsemi\fontsize{9.3}{12.6}\selectfont\color{popink}#7
  \end{tcolorbox}
  \noindent\begin{minipage}[t]{0.487\linewidth}
    \begin{tcolorbox}[enhanced,colback=popgreen,colframe=popink,boxrule=2.2pt,arc=10pt,
        drop shadow=popink,left=11pt,right=11pt,top=10pt,bottom=10pt,height=3.1cm,valign=center]
      \tcbox[colback=white,colframe=popink,boxrule=1.3pt,arc=5pt,boxsep=0pt,nobeforeafter,
        left=3pt,right=3pt,top=3pt,bottom=3pt,valign=center]{\qrcode[height=1.9cm]{https://play.google.com/store/apps/details?id=com.luvvix.onbuch&hl=fr}}%
      \hspace{8pt}\begin{minipage}[c]{0.5\linewidth}
        {\popdisplay\fontsize{11}{12.5}\selectfont\color{white}\MakeUppercase{Télécharge OnBuch}}\par\vspace{4pt}
        {\raggedright\popsemi\fontsize{8.4}{11}\selectfont\color{white}Gratuit \textbullet\ Google Play\\Tuteur IA, annales, concours, résultats\par}
      \end{minipage}
    \end{tcolorbox}
  \end{minipage}\hfill
  \begin{minipage}[t]{0.487\linewidth}
    \begin{tcolorbox}[enhanced,colback=poppurple,colframe=popink,boxrule=2.2pt,arc=10pt,
        drop shadow=popink,left=11pt,right=11pt,top=10pt,bottom=10pt,height=3.1cm,valign=center]
      \tikz[baseline=-0.7ex]\node[circle,fill=white,draw=popink,line width=1.3pt,minimum size=1.6cm,inner sep=0]{\popdisplay\fontsize{24}{24}\selectfont\color{poppurple}+};%
      \hspace{8pt}\begin{minipage}[c]{0.6\linewidth}
        {\popdisplay\fontsize{11}{12.5}\selectfont\color{white}\MakeUppercase{Passe à OnBuch+}}\par\vspace{4pt}
        {\raggedright\popsemi\fontsize{8.4}{11}\selectfont\color{white}Tous les cours signés, Tuteur IA illimité, fiches PDF — onglet OnBuch+ de l'app\par}
      \end{minipage}
    \end{tcolorbox}
  \end{minipage}
  \vspace{8pt}
  \popcontenu
  \vfill
  \signatureonbuch
  \restoregeometry
  \newpage
}

% Rangée « Ce cours contient » (page de garde)
\newcommand{\poptile}[3]{% #1 couleur #2 gros texte #3 légende
  \begin{tcolorbox}[enhanced,colback=#1,colframe=popink,boxrule=2pt,arc=9pt,shadow={2.5pt}{-2.5pt}{0pt}{popink},
      left=6pt,right=6pt,top=9pt,bottom=9pt,halign=center,valign=center,height=1.75cm,width=\linewidth,nobeforeafter]
    {\popdisplay\fontsize{12.5}{13}\selectfont\color{white}\MakeUppercase{#2}}\par\vspace{3pt}
    {\centering\popsemi\fontsize{7.8}{9.5}\selectfont\color{white}#3\par}
  \end{tcolorbox}}
\newcommand{\popcontenu}{%
  \par\noindent{\popxboldls\fontsize{7.5}{7.5}\selectfont\color{popmuted}CE COURS SIGNÉ CONTIENT}\par\vspace{7pt}
  \noindent\begin{minipage}[t]{0.235\linewidth}\poptile{popblue}{Cours}{complet, illustré, pas à pas}\end{minipage}\hfill
  \begin{minipage}[t]{0.235\linewidth}\poptile{poporange}{Méthodes}{et exercices résolus}\end{minipage}\hfill
  \begin{minipage}[t]{0.235\linewidth}\poptile{poppink}{Exercices}{type examen + corrigés}\end{minipage}\hfill
  \begin{minipage}[t]{0.235\linewidth}\poptile{popgreen}{Fiche bilan}{l'essentiel à retenir}\end{minipage}\par}

% Sommaire
\newtcolorbox{sommaire}{enhanced,colback=white,colframe=popink,boxrule=2.2pt,arc=10pt,
  drop shadow=popink,left=18pt,right=18pt,top=13pt,bottom=13pt,before skip=6pt,after skip=20pt,
  before upper={\poppill{Sommaire}{poppurple}{white}\par\vspace{9pt}\popsemi\fontsize{10.6}{14.5}\selectfont\color{popink}}}

% Signature de fin de cours — poussée tout en bas de la dernière page
\newcommand{\findecours}{%
  \par\vfill\nopagebreak
  \begin{center}\popsquiggle{poporange}\end{center}\vspace{4pt}
  \signatureonbuch
}
\AtBeginDocument{\def\llbracket{\mathopen{[\mkern-3.5mu[}}\def\rrbracket{\mathclose{]\mkern-3.5mu]}}} % intervalles d'entiers (glyphe ⟦ absent de la police)
% Glyphes absents de Fira Math : redéfinis à partir de glyphes disponibles
\AtBeginDocument{%
  \def\setminus{\mathbin{\backslash}}\let\smallsetminus\setminus
  \def\vdots{\mathord{\vbox{\baselineskip3.2pt\lineskiplimit0pt\kern4pt\hbox{.}\hbox{.}\hbox{.}}}}%
  \def\ddots{\mathinner{\mkern1mu\raise6.5pt\hbox{.}\mkern2mu\raise3.6pt\hbox{.}\mkern2mu\raise0.7pt\hbox{.}\mkern1mu}}%
  \def\triangle{\mathord{\scalebox{0.9}{$\symup{\Delta}$}}}%
  \def\nabla{\mathord{\raisebox{\depth}{\scalebox{1}[-1]{$\symup{\Delta}$}}}}%
  \def\checkmark{\popcheck{popdark}}%
  \def\star{\mathbin{\ast}}\def\bigstar{\mathord{\ast}}%
  \def\diamond{\mathbin{\scalebox{0.6}{$\Diamond$}}}%
  \def\bigcup{\mathop{\mathchoice{\raisebox{-0.3em}{\scalebox{1.7}{$\cup$}}}{\scalebox{1.25}{$\cup$}}{\cup}{\cup}}\displaylimits}%
  \def\bigcap{\mathop{\mathchoice{\raisebox{-0.3em}{\scalebox{1.7}{$\cap$}}}{\scalebox{1.25}{$\cap$}}{\cap}{\cap}}\displaylimits}%
  \def\lesseqqgtr{\mathrel{\lessgtr}}\def\bigoplus{\mathop{\oplus}\displaylimits}%
}
\providecommand{\ssi}{si et seulement si} % abréviation fréquente
\AtBeginDocument{\providecommand{\wideparen}{\overparen}} % arc de cercle

%% FIGFIX-BEGIN v4 — correctifs globaux des figures (docs/onbuch-generation/tools/figfix)
\usepackage{adjustbox}\usepackage{etoolbox}
% toute figure TikZ/pgfplots plus large que la ligne est réduite à la largeur disponible
\BeforeBeginEnvironment{tikzpicture}{\begin{adjustbox}{max width=\linewidth}}
\AfterEndEnvironment{tikzpicture}{\end{adjustbox}}
% légendes pgfplots lisibles par-dessus les courbes
\pgfplotsset{every axis/.append style={legend style={fill=white,fill opacity=0.92,text opacity=1,draw=black!15,inner sep=3pt}}}
% étiquettes d'arêtes (syntaxe edge["..."]) avec fond blanc
\tikzset{every edge quotes/.append style={fill=white,inner sep=1.2pt}}
%% FIGFIX-END
\begin{document}

\pagegarde{VI : rédiger un Curriculum Vitae}{Français}{Tle Tech}{Français – classe de Terminale technique}{N27}{9 séances (fiche de progression harmonisée)}{Cette leçon apprend à lire, analyser et rédiger un Curriculum Vitae adapté à une situation d'embauche. Elle mobilise l'étude des textes fonctionnels et iconiques, les valeurs du subjonctif et du conditionnel, et aboutit à la production et à la présentation d'un CV personnel.
\vspace{4pt}
\begin{multicols}{2}\small
\begin{itemize}
\item À la fin de cette leçon, je sais définir le CV et identifier sa fonction dans une situation de communication professionnelle
\item À la fin de cette leçon, je sais repérer la structure externe d'un CV (état civil, formation, expérience, compétences, centres d'intérêt)
\item À la fin de cette leçon, je sais analyser la structure interne d'un CV (mise en page, rubriques, ordre chronologique ou thématique)
\item À la fin de cette leçon, je sais employer correctement le subjonctif et le conditionnel dans les formulations de souhait, d'hypothèse et de politesse propres à la candidature
\item À la fin de cette leçon, je sais rédiger un CV complet, lisible et sans faute
\item À la fin de cette leçon, je sais adapter un CV à une offre d'embauche précise en sélectionnant les informations pertinentes
\end{itemize}\end{multicols}}

\begin{sommaire}
\begin{multicols}{2}
\begin{enumerate}
\item Découverte : qu'est-ce qu'un CV ?
\item La structure externe du CV
\item La structure interne du CV
\item Langue : les valeurs du subjonctif et du conditionnel
\item Textes iconiques et adaptation du CV à une offre d'embauche
\item Rédaction, présentation et évaluation du CV
\item Activité d'intégration
\item Méthodes et exercices résolus
\item Exercices
\item Corrigés des exercices
\item Fiche bilan
\end{enumerate}
\end{multicols}
\end{sommaire}

\begin{objectifs}
\begin{itemize}
\item À la fin de cette leçon, je sais définir le CV et identifier sa fonction dans une situation de communication professionnelle
\item À la fin de cette leçon, je sais repérer la structure externe d'un CV (état civil, formation, expérience, compétences, centres d'intérêt)
\item À la fin de cette leçon, je sais analyser la structure interne d'un CV (mise en page, rubriques, ordre chronologique ou thématique)
\item À la fin de cette leçon, je sais employer correctement le subjonctif et le conditionnel dans les formulations de souhait, d'hypothèse et de politesse propres à la candidature
\item À la fin de cette leçon, je sais rédiger un CV complet, lisible et sans faute
\item À la fin de cette leçon, je sais adapter un CV à une offre d'embauche précise en sélectionnant les informations pertinentes
\item À la fin de cette leçon, je sais exploiter des documents iconiques (photo d'identité, logos, pictogrammes) dans un CV
\item À la fin de cette leçon, je sais présenter oralement mon CV devant un jury simulant un entretien d'embauche
\item À la fin de cette leçon, je sais justifier mes choix de rédaction et évaluer un CV selon une grille de critères
\end{itemize}
\end{objectifs}

\begin{prerequis}
\begin{itemize}
\item Les types de textes et les textes fonctionnels (lettre administrative, lettre de motivation) vus en Première
\item Les modes verbaux : indicatif, impératif, infinitif (notions de base)
\item La mise en page d'un document administratif (en-tête, corps, signature)
\item Le vocabulaire du monde du travail (poste, employeur, candidat, recrutement)
\item La lecture de documents iconiques (photographies, schémas) vue en classe de Seconde
\end{itemize}
\end{prerequis}

\newpage

\coursec{Découverte : qu'est-ce qu'un CV ?}

\courssub{Situation d'accroche : 350 candidatures, 12 retenues}

À Douala, une entreprise de BTP publie une offre d'emploi de technicien en génie civil. En une semaine, elle reçoit 350 candidatures, mais le recruteur ne retient que 12 CV pour les entretiens. Parmi les écartés : des diplômés compétents dont les CV étaient mal structurés, truffés de fautes ou inadaptés au poste. Comment rédiger un \cle{Curriculum Vitae} clair, honnête et convaincant qui donne envie à un employeur camerounais de vous recevoir ? C'est tout l'enjeu de cette leçon.

Cette situation n'a rien d'exceptionnel. Chaque année, des milliers de jeunes techniciens camerounais — titulaires d'un Baccalauréat technique, d'un BTS ou d'une qualification professionnelle — postulent à des emplois, des stages ou des concours. Or, avant même de rencontrer le candidat, l'employeur lit un document : le \cle{Curriculum Vitae}. Ce document est donc votre premier ambassadeur. S'il est négligé, vous ne serez jamais convoqué, aussi compétent soyez-vous. Toute cette unité vise à vous apprendre à construire ce document pas à pas. Commençons par le comprendre.

\courssub{Définition et étymologie du Curriculum Vitae}

Imaginons d'abord l'intuition. Vous ne pouvez pas raconter toute votre vie à un employeur en quelques minutes. Il vous faut donc un document qui \emph{résume} votre parcours, qui le \emph{structure} pour qu'on le lise vite, et qui le \emph{prouve} par des faits vérifiables. Ce document existe : c'est le CV.

L'expression \emph{curriculum vitae} vient du latin : \emph{curriculum} signifie « course, carrière » et \emph{vitae} signifie « de la vie ». Le Curriculum Vitae est donc littéralement la \cle{course de la vie} : le résumé du chemin parcouru par une personne, de sa naissance à sa situation actuelle, du point de vue de sa formation et de son travail.

\begin{definition}[Curriculum Vitae]
Le \cle{Curriculum Vitae} (abrégé \cle{CV}) est un document écrit qui présente de manière synthétique et structurée l'\cle{état civil}, la \cle{formation}, l'\cle{expérience professionnelle} et les \cle{compétences} d'une personne candidate à un emploi, à un stage ou à un concours.
\end{definition}

Décomposons cette définition pour bien la comprendre :

\begin{itemize}
\item \emph{document écrit} : le CV n'est pas une parole, c'est un texte que l'on rédige, que l'on met en page et que l'on transmet (sur papier ou par voie numérique) ;
\item \emph{synthétique} : il tient en une à deux pages ; on ne raconte pas tout, on sélectionne l'essentiel ;
\item \emph{structuré} : les informations sont rangées en rubriques claires (état civil, formation, expérience, compétences, centres d'intérêt) ;
\item \emph{candidate à un emploi} : le CV n'existe que dans une situation précise — la recherche d'un emploi, d'un stage ou l'inscription à un concours.
\end{itemize}

\begin{aretenir}[L'essentiel]
Le CV est un \cle{texte fonctionnel} : il n'est pas écrit pour le plaisir ou pour exprimer des sentiments, mais pour obtenir un résultat concret dans la vie sociale et professionnelle. Comme le certificat, la facture ou la demande d'emploi, il appartient à la famille des écrits utilitaires.
\end{aretenir}

\courssub{La fonction du CV : à quoi sert-il ?}

Pourquoi rédige-t-on un CV ? La réponse semble évidente — « pour trouver du travail » — mais elle est à préciser. Le CV n'obtient pas directement l'emploi : il obtient \cle{l'entretien d'embauche}. C'est une distinction capitale. Le CV est une \emph{porte d'entrée} : sa fonction est de convaincre le recruteur, en quelques dizaines de secondes de lecture, que votre profil mérite d'être rencontré. L'embauche elle-même se joue ensuite, lors de l'entretien.

Le CV sert aussi dans d'autres situations de la vie courante :

\begin{itemize}
\item se présenter à un \cle{concours} administratif ou professionnel (concours de la fonction publique, écoles de formation) ;
\item demander un \cle{stage} académique ou professionnel en entreprise ;
\item postuler à une \emph{bourse} ou à une formation à l'étranger ;
\item se faire connaître d'un partenaire ou d'un client lorsqu'on est travailleur indépendant.
\end{itemize}

\begin{methode}[Identifier la fonction d'un texte fonctionnel]
Pour déterminer la fonction d'un texte fonctionnel (CV, lettre, attestation, facture...), posez-vous trois questions :
\begin{enumerate}
\item \textbf{À qui} le texte s'adresse-t-il ? (identifier le destinataire)
\item \textbf{Dans quelle situation} est-il produit ? (identifier le contexte : embauche, concours, démarche administrative...)
\item \textbf{Pour obtenir quel résultat} ? (identifier le but : être convoqué, obtenir un document, être admis...)
\end{enumerate}
La réponse à ces trois questions donne la fonction du texte. Pour le CV : il s'adresse à un recruteur, dans une situation de candidature, pour obtenir un entretien.
\end{methode}

\courssub{La situation d'énonciation du CV}

Tout texte est produit par quelqu'un, pour quelqu'un, dans un but. C'est ce qu'on appelle la \cle{situation d'énonciation}. Analysons celle du CV.

\begin{itemize}
\item \textbf{L'émetteur} est le \cle{candidat} : c'est vous, élève de Terminale technique devenu demandeur d'emploi ou de stage. Il parle de lui-même, mais à la troisième personne des faits : le CV n'est pas un récit à la première personne, c'est une fiche d'informations.
\item \textbf{Le récepteur} est le \cle{recruteur} : chef d'entreprise, directeur des ressources humaines, jury de concours. C'est un lecteur pressé, qui compare des dizaines de candidatures.
\item \textbf{Le but} est \cle{persuasif} : le candidat cherche à convaincre le recruteur qu'il possède le profil recherché. Le CV est donc un texte à visée argumentative, même s'il n'argumente pas ouvertement : il persuade par la sélection et la présentation des faits.
\item \textbf{Le canal} est écrit : papier remis en main propre, courrier, ou fichier envoyé par courriel.
\end{itemize}

\begin{popfigure}
\begin{center}
\begin{tikzpicture}[>=Stealth, node distance=1.2cm]
\node[draw=popblue, fill=popblueL, rounded corners, minimum width=3.2cm, minimum height=1.4cm, align=center, font=\small\bfseries] (cand) {CANDIDAT\\ \normalfont\footnotesize (émetteur)};
\node[draw=poporange, fill=poporangeL, rounded corners, minimum width=3.2cm, minimum height=1.4cm, align=center, font=\small\bfseries, right=3.4cm of cand] (cv) {CURRICULUM VITAE\\ \normalfont\footnotesize (message écrit)};
\node[draw=popgreen, fill=popgreenL, rounded corners, minimum width=3.2cm, minimum height=1.4cm, align=center, font=\small\bfseries, right=3.4cm of cv] (rec) {RECRUTEUR\\ \normalfont\footnotesize (récepteur)};
\draw[->, very thick, popblue] (cand) -- node[above, font=\footnotesize, align=center]{rédige et\\ transmet} (cv);
\draw[->, very thick, poporange] (cv) -- node[above, font=\footnotesize, align=center]{informe et\\ persuade} (rec);
\draw[->, very thick, popgreen, dashed] (rec.south) to[bend left=25] node[below, font=\footnotesize, align=center]{convocation à l'entretien\\ (résultat attendu)} (cand.south);
\end{tikzpicture}
\end{center}
\legende{Schéma de la situation de communication du CV : le candidat (émetteur) rédige un CV (message écrit) qu'il transmet au recruteur (récepteur) ; la réponse attendue est une convocation à l'entretien.}
\end{popfigure}

Ce schéma montre une particularité : le recruteur ne répond pas toujours. Si le CV ne convainc pas, il est simplement écarté, sans explication. D'où l'importance de soigner ce document : il parle \emph{à votre place}, en votre absence.

\courssub{Ne pas confondre : CV, lettre de motivation, carte professionnelle}

Dans la vie professionnelle, plusieurs documents se ressemblent mais n'ont ni la même fonction ni le même contenu. Distinguons-les clairement.

\begin{center}
\begin{tabularx}{\linewidth}{|l|X|X|}
\hline
\rowcolor{popblueL} \textbf{Critère} & \textbf{Curriculum Vitae} & \textbf{Lettre de motivation} \\
\hline
Destinataire & Le recruteur (entreprise, jury de concours) & Le recruteur (même personne, même situation) \\
\hline
Contenu & Des \textbf{faits objectifs} : état civil, diplômes, expériences, compétences, présentés en rubriques & Une \textbf{argumentation subjective} : motivations, projet, qualités mises en valeur, rédigée en phrases \\
\hline
Forme & Fiche structurée, style télégraphique, pas de phrases complètes obligatoires & Lettre rédigée (formules de politesse, paragraphes, signature) \\
\hline
Objectif & Prouver que le profil correspond au poste & Convaincre que le candidat est le bon choix et donner envie de le rencontrer \\
\hline
\end{tabularx}
\end{center}

\begin{attention}[Erreur fréquente]
Ne confondez jamais le CV avec la lettre de motivation. Le CV présente des \cle{faits objectifs} (dates, diplômes, postes occupés) ; la lettre de motivation développe une \cle{argumentation subjective} (pourquoi vous voulez ce poste, pourquoi vous êtes fait pour lui). À l'examen comme dans la vie, transformer son CV en lettre rédigée — ou l'inverse — est une faute de genre textuel lourdement sanctionnée. Autre piège : la \cle{carte professionnelle} est un document administratif délivré \emph{après} l'embauche (elle atteste qu'on est salarié d'une entreprise) ; elle ne sert pas à chercher un emploi.
\end{attention}

\courssub{Analyse d'un premier document : le CV d'un jeune technicien camerounais}

Observons maintenant un document réel-type. Voici, reproduit de manière simplifiée, le CV d'un jeune technicien camerounais du secteur agropastoral qui postule à un emploi dans une exploitation agricole modernisée de l'Ouest.

\begin{exemplebox}[CV de M. Etoundi Basile, technicien agropastoral]
\textbf{ETOUNDI BASILE}\\
Né le 12 mars 2001 à Bafoussam — Célibataire\\
Tél. : 6 90 00 00 00 — Courriel : b.etoundi@mail.cm\\
Quartier : Akwa-Nord, Douala

\textbf{FORMATION}
\begin{itemize}
\item 2024 : Baccalauréat technique série F4 (Génie rural), Lycée technique de Bafoussam
\item 2022 : Probatoire technique F4, Lycée technique de Bafoussam
\end{itemize}

\textbf{EXPÉRIENCE PROFESSIONNELLE}
\begin{itemize}
\item Juin--sept. 2024 : stage ouvrier à la SODEPALM (entretien des palmeraies, suivi de production)
\item 2023--2024 : responsable de la basse-cour du club d'agropastoral du lycée
\end{itemize}

\textbf{COMPÉTENCES}
\begin{itemize}
\item Conduite d'élevage avicole et porcin ; techniques de culture vivrière
\item Maîtrise de l'outil informatique (Word, Excel)
\end{itemize}

\textbf{LANGUES} : français (courant), anglais (scolaire), ghomala (langue maternelle)

\textbf{CENTRES D'INTÉRÊT} : football, lecture agricole
\end{exemplebox}

Observons ce document ensemble, comme des enquêteurs. Que remarquons-nous ?

\begin{enumerate}
\item \textbf{La longueur} : tout tient sur \cle{une page}. C'est la règle pour un jeune candidat : une page, deux au maximum pour un cadre expérimenté. Un recruteur n'a pas le temps de lire un roman.
\item \textbf{Le ton} : il est \cle{neutre} et factuel. Aucune phrase comme « je suis le meilleur » ou « j'adore les animaux » : seulement des faits (dates, diplômes, postes). L'éloge de soi-même appartient à la lettre de motivation.
\item \textbf{Les informations sont vérifiables} : chaque diplôme correspond à un établissement réel, chaque stage à une entreprise identifiable. Un recruteur peut téléphoner pour confirmer. Mentir sur un CV (faux diplôme, fausse expérience) est une faute grave qui peut coûter l'emploi, même après l'embauche.
\item \textbf{L'ordre} : les informations les plus récentes apparaissent en premier dans chaque rubrique (ordre \emph{antichronologique}) : le recruteur veut savoir d'abord où vous en êtes aujourd'hui.
\item \textbf{Les rubriques} : état civil et coordonnées en tête, puis formation, expérience, compétences, langues, centres d'intérêt. Nous étudierons cette structure en détail dans la section suivante.
\end{enumerate}

\begin{savaistu}[Moins d'une minute pour convaincre]
Des études sur le recrutement montrent qu'un employeur consacre en moyenne \cle{moins d'une minute} — parfois trente secondes — à la première lecture d'un CV. Il ne lit pas : il \emph{balaye} le document des yeux, à la recherche du diplôme, de l'expérience et des compétences qui correspondent au poste. La clarté de la présentation (rubriques visibles, aération, orthographe irréprochable) est donc aussi décisive que le contenu. Un CV excellent mais confus part à la corbeille ; un CV honnête et limpide décroche l'entretien.
\end{savaistu}

\courssub{Exercice résolu : identifier la fonction d'un document}

\begin{exoresolu}[Reconnaître un texte fonctionnel]
\textbf{Énoncé.} Voici trois documents produits par Aminatou, jeune diplômée d'un BTS en comptabilité : (a) une fiche d'une page présentant ses diplômes et son stage à la SABC ; (b) une lettre rédigée expliquant au directeur d'une PME de Yaoundé pourquoi elle souhaite travailler dans son entreprise ; (c) la carte que l'entreprise lui a remise après son embauche, portant son nom, son poste et son matricule. Identifiez chaque document et justifiez votre réponse en précisant sa fonction.
\tcblower
\textbf{Solution détaillée.} Appliquons la méthode des trois questions (à qui ? dans quelle situation ? pour quel résultat ?).

(a) Il s'agit du \textbf{Curriculum Vitae}. Il s'adresse à un recruteur, dans une situation de candidature, pour obtenir un entretien. Ses indices formels : une page, des rubriques, des faits objectifs (diplômes, stage), un ton neutre.

(b) Il s'agit de la \textbf{lettre de motivation}. Même destinataire et même situation que le CV, mais fonction différente : argumenter et convaincre. Ses indices : texte rédigé en phrases, formules de politesse, expression des motivations personnelles.

(c) Il s'agit de la \textbf{carte professionnelle}. Elle est délivrée par l'employeur \emph{après} l'embauche ; sa fonction est d'attester de la qualité de salarié (preuve d'appartenance à l'entreprise), non de chercher un emploi. Ses indices : nom, poste, matricule, émise par l'entreprise.

Conclusion : les trois documents appartiennent à la vie professionnelle, mais seul le CV résume objectivement le parcours du candidat \emph{avant} l'embauche.
\end{exoresolu}

\begin{exercice}[À vous de jouer]
Moussa, élève de Terminale technique spécialité F3 (Électrotechnique), veut postuler à un stage dans une entreprise d'installation électrique de Garoua. 1) Qui est l'émetteur et qui est le récepteur du CV qu'il va rédiger ? 2) Quel est le but précis de ce document ? 3) Citez trois informations objectives que son CV devra contenir. 4) Pourquoi ne doit-il pas écrire dans son CV : « Je suis le meilleur électricien de ma génération » ?
\end{exercice}

\begin{corrige}[Exercice : le stage de Moussa]
1) L'émetteur est Moussa (le candidat) ; le récepteur est le recruteur de l'entreprise d'installation électrique de Garoua.

2) Le but du CV est d'obtenir une convocation à un entretien (ou l'acceptation du stage) : persuader le recruteur que le profil de Moussa correspond au stage proposé.

3) Informations objectives attendues : son état civil et ses coordonnées ; sa formation (classe de Terminale F3, établissement, diplômes déjà obtenus comme le Probatoire) ; ses compétences techniques (câblage, lecture de schémas électriques, sécurité électrique) ; éventuellement ses langues et centres d'intérêt.

4) Cette phrase est un jugement subjectif et invérifiable : elle relève de l'argumentation personnelle, qui a sa place dans la lettre de motivation, pas dans le CV. Le CV n'énonce que des faits objectifs et vérifiables ; une telle affirmation donnerait en outre une impression de prétention.
\end{corrige}

\begin{aretenir}[Bilan de la découverte]
\begin{itemize}
\item Le Curriculum Vitae (« course de la vie » en latin) est un texte fonctionnel qui présente de façon synthétique et structurée l'état civil, la formation, l'expérience et les compétences d'un candidat.
\item Sa fonction : obtenir un entretien d'embauche, un stage ou l'admission à un concours.
\item Sa situation d'énonciation : un candidat (émetteur) s'adresse par écrit à un recruteur (récepteur) dans un but persuasif.
\item Il se distingue de la lettre de motivation (argumentation subjective) et de la carte professionnelle (document délivré après l'embauche).
\item Un bon CV est court (1 à 2 pages), de ton neutre, et ne contient que des informations vérifiables.
\end{itemize}
\end{aretenir}

Maintenant que nous savons ce qu'est un CV et à quoi il sert, la prochaine étape consiste à examiner de près sa \cle{structure externe} : comment ce document est-il physiquement organisé sur la page ? C'est l'objet de la section suivante.

\coursec{La structure externe du CV}

Dans la section précédente, nous avons découvert ce qu'est un \cle{curriculum vitae} : un document bref et fonctionnel qui présente un candidat à un employeur. Nous avons vu que le CV n'est pas une biographie complète, mais un outil de \emph{sélection} : le recruteur y cherche, en quelques dizaines de secondes, la preuve que le candidat correspond au poste. Encore faut-il que cette preuve soit \emph{facile à trouver}. C'est précisément l'objet de cette section : étudier la \cle{structure externe} du CV, c'est-à-dire son plan visible, ses parties, ses rubriques et leur ordonnancement. Un bon contenu mal organisé vaut moins qu'un contenu moyen bien présenté : apprenons donc à construire le squelette du document.

\courssub{La notion de rubrique}

Observez deux candidats qui postulent au même emploi de magasinier à Douala. Le premier écrit un long paragraphe continu où se mélangent son adresse, son BEPC, son stage à la SABC et sa passion pour le football. Le second présente les mêmes informations, mais rangées sous des titres clairs. Lequel sera lu jusqu'au bout ? Évidemment le second. Cette observation fondamentale nous conduit à la notion clé de cette leçon.

\begin{definition}[Rubrique]
Une \cle{rubrique} est une partie d'un document regroupant des informations de même nature sous un titre commun. Dans un CV, chaque rubrique porte un intitulé visible (en gras, souvent en majuscules ou en couleur) : « Formation », « Expérience professionnelle », « Compétences », etc.
\end{definition}

La rubrique joue un double rôle. Pour le \emph{rédacteur}, elle impose une discipline : on ne met dans une rubrique que ce qui lui appartient. Pour le \emph{lecteur}, elle sert de repère : le recruteur qui cherche les diplômes va directement à « Formation » sans lire le reste.

\courssub{L'en-tête : l'identité du candidat}

L'\cle{en-tête} est la première chose que voit le recruteur. Il occupe le haut de la page et contient :

\begin{itemize}
\item le \textbf{nom} (en majuscules) et le \textbf{prénom} du candidat ;
\item les \textbf{coordonnées} : adresse (quartier, ville, boîte postale éventuelle), numéro de téléphone, adresse électronique ;
\item éventuellement, le \textbf{titre du poste visé} ou un titre professionnel (« Technicien en froid et climatisation », « Candidate au poste de secrétaire comptable »), qui montre immédiatement que le CV répond à une offre précise.
\end{itemize}

\begin{exemplebox}[Un en-tête correct]
\textbf{MBALLA ESSOMBA Jean-Claude}\par
Technicien en maintenance industrielle\par
Quartier Nkolbisson, Yaoundé — B.P. 4521\par
Tél. : 6 90 12 34 56 — E-mail : jc.mballa@mail.com
\end{exemplebox}

\begin{attention}[Erreur fréquente : des coordonnées négligées]
Deux fautes ruinent de nombreuses candidatures dès l'en-tête. Premièrement, \textbf{oublier de mettre à jour ses coordonnées} : un numéro de téléphone qui ne sonne plus ou une ancienne adresse empêchent le recruteur de convoquer le candidat — le CV, même excellent, devient inutile. Deuxièmement, donner une \textbf{adresse e-mail fantaisiste} du type « petitprince237@... » ou « beaugossedouala@... » : elle nuit gravement au sérieux de la candidature. Préférez une adresse sobre formée de votre prénom et de votre nom.
\end{attention}

\courssub{La rubrique « Formation / Études et diplômes »}

Cette rubrique présente le parcours scolaire et les diplômes obtenus. La règle d'or est l'\cle{ordre antichronologique} : on commence par le diplôme \emph{le plus récent} et l'on remonte vers le plus ancien. Pourquoi ? Parce que ce qui intéresse le recruteur, c'est votre niveau \emph{actuel}, pas votre CEP.

\begin{exemplebox}[Rubrique « Formation » bien ordonnée]
\textbf{FORMATION}\par
2025 : Probatoire technique, série F3 (Génie civil), Lycée technique de Nkongsamba.\par
2024 : BEPC, Collège bilingue de Bangangté.\par
2023 : stage de découverte du métier de maçon (attestation de fin de stage).
\end{exemplebox}

Chaque ligne comporte idéalement trois éléments : la \textbf{date} (ou la période), le \textbf{diplôme ou la formation}, l'\textbf{établissement}. On peut ajouter la mention (« avec mention assez bien ») si elle valorise le candidat.

\begin{exemplebox}[Exemple chiffré : l'importance d'une rubrique « Formation » claire]
Sur 100 CV reçus pour un poste de secrétaire comptable à Bafoussam, \textbf{30 sont écartés dès la première lecture} pour absence de rubrique « Formation » claire : diplômes noyés dans un paragraphe, dates absentes, ordre illogique. Cela représente 30 \% des candidatures, soit près d'un candidat sur trois, éliminé non pas parce qu'il est incompétent, mais parce que sa structure externe est défaillante. La forme protège le fond.
\end{exemplebox}

\courssub{La rubrique « Expérience professionnelle »}

Cette rubrique recense les \textbf{emplois}, \textbf{stages} et travaux déjà effectués. Pour un élève de Terminale technique, les stages scolaires et les petits boulots de vacances comptent pleinement : ils prouvent le contact avec le monde du travail. Chaque expérience se présente avec :

\begin{itemize}
\item les \textbf{dates} (période de début et de fin) ;
\item le \textbf{poste occupé} ou la nature de la mission ;
\item la \textbf{structure d'accueil} (entreprise, administration, cabinet).
\end{itemize}

\begin{exemplebox}[Rubrique « Expérience professionnelle »]
\textbf{EXPÉRIENCE PROFESSIONNELLE}\par
Juillet--août 2025 : stage ouvrier à la brasserie SABC de Douala (chaîne d'embouteillage).\par
Juillet 2024 : stage d'observation à la SONARA de Limbé (service sécurité industrielle).\par
Décembre 2023 : aide-comptable dans un cabinet comptable de Yaoundé (classement des pièces, saisie).
\end{exemplebox}

Là encore, on respecte l'ordre antichronologique : l'expérience la plus récente d'abord. On peut ajouter une ligne de \emph{tâches réalisées} pour chaque poste, mais avec sobriété : deux ou trois verbes d'action suffisent (« saisie des factures, accueil des clients »).

\courssub{La rubrique « Compétences »}

La rubrique \cle{Compétences} répond à la question : « Que savez-vous \emph{faire} concrètement ? » On y distingue deux familles :

\begin{itemize}
\item les \textbf{compétences techniques} : maîtrise de l'outil informatique (traitement de texte, tableur), conduite d'engins, soudure, câblage électrique, utilisation d'un logiciel de comptabilité, permis de conduire (catégorie B, C...) ;
\item les \textbf{compétences linguistiques} : français (lu, écrit, parlé), anglais (niveau scolaire, courant...), langues nationales (ewondo, duala, fulfuldé...), précieuses dans un pays multilingue comme le Cameroun.
\end{itemize}

\begin{exemplebox}[Rubrique « Compétences »]
\textbf{COMPÉTENCES}\par
Informatique : Word, Excel, saisie rapide au clavier.\par
Technique : lecture de plans, conduite de chariot élévateur (attestation SABC).\par
Langues : français (courant), anglais (niveau scolaire), ewondo (langue maternelle).\par
Permis de conduire catégorie B.
\end{exemplebox}

\courssub{La rubrique « Centres d'intérêt / Divers »}

Cette dernière rubrique, facultative mais recommandée, humanise le candidat. On y mentionne des \textbf{activités valorisantes} : sport pratiqué en club ou en compétition (football, athlétisme), engagement associatif (membre d'une chorale, d'un club de jeunes, d'une association de quartier), passions utiles au poste (lecture, mécanique). Évitez les mentions vagues ou sans valeur professionnelle (« regarder la télévision »).

\begin{attention}[Ce qu'il faut éviter dans « Centres d'intérêt »]
Ne citez que des activités que vous pouvez assumer en entretien. Affirmer « capitaine de l'équipe de handball » alors qu'on n'y a jamais joué expose à une question piège du recruteur. Une passion réelle et modeste vaut mieux qu'un mensonge spectaculaire.
\end{attention}

\courssub{La photo, la date et la signature}

Trois éléments complètent la structure externe :

\begin{itemize}
\item la \textbf{photo d'identité} : si elle est utilisée, elle doit être récente, de bonne qualité, sur fond neutre, le candidat en tenue correcte et regardant l'objectif. Une photo de fête ou floue décrédibilise toute la candidature ;
\item la mention \textbf{« Fait à..., le... »} : selon les usages administratifs camerounais, on indique le lieu et la date de rédaction (« Fait à Yaoundé, le 12 mars 2026 ») ;
\item la \textbf{signature} : apposée à la main sur la version papier, elle engage le candidat sur l'exactitude des informations fournies.
\end{itemize}

\begin{savaistu}[La photo d'identité au Cameroun]
Au Cameroun, la photo d'identité reste souvent \emph{exigée} dans les candidatures administratives (concours, recrutements publics, dossiers d'embauche), contrairement à d'autres pays — comme la France ou le Canada — où elle est déconseillée pour lutter contre les discriminations à l'embauche. En contexte camerounais, mieux vaut donc prévoir une photo sérieuse, sauf si l'offre précise le contraire.
\end{savaistu}

\courssub{Maquette d'un CV vierge}

La figure suivante synthétise la structure externe complète : chaque cadre correspond à une rubrique, placée à sa position habituelle sur la page.

\begin{popfigure}
\begin{center}
\begin{tikzpicture}
\draw[thick,popdark] (0,0) rectangle (10.5,13.2);
\draw[fill=popblueL,draw=popblue,thick] (0.4,11.2) rectangle (7.6,12.8);
\node[align=left,anchor=west] at (0.6,12.4) {\textbf{EN-T\^ETE}};
\node[align=left,anchor=west] at (0.6,11.8) {\footnotesize Nom, pr\'enom, titre du poste vis\'e};
\node[align=left,anchor=west] at (0.6,11.45) {\footnotesize Adresse, t\'el\'ephone, e-mail};
\draw[fill=popcream,draw=popgold,thick] (8.0,11.2) rectangle (10.1,12.8);
\node[align=center] at (9.05,12.0) {\footnotesize Photo};
\draw[fill=popgreenL,draw=popgreen,thick] (0.4,8.2) rectangle (10.1,10.9);
\node[align=left,anchor=north west] at (0.6,10.7) {\textbf{FORMATION / \'ETUDES ET DIPL\^OMES}};
\node[align=left,anchor=north west] at (0.6,10.2) {\footnotesize Du dipl\^ome le plus r\'ecent au plus ancien};
\node[align=left,anchor=north west] at (0.6,9.8) {\footnotesize 2025 : Probatoire technique, Lyc\'ee technique de...};
\node[align=left,anchor=north west] at (0.6,9.4) {\footnotesize 2024 : BEPC, Coll\`ege de...};
\draw[fill=poporangeL,draw=poporange,thick] (0.4,5.2) rectangle (10.1,7.9);
\node[align=left,anchor=north west] at (0.6,7.7) {\textbf{EXP\'ERIENCE PROFESSIONNELLE}};
\node[align=left,anchor=north west] at (0.6,7.2) {\footnotesize Dates -- poste -- structure (stages, emplois)};
\node[align=left,anchor=north west] at (0.6,6.8) {\footnotesize 2025 : stage \`a la SABC de Douala...};
\node[align=left,anchor=north west] at (0.6,6.4) {\footnotesize 2024 : stage \`a la SONARA de Limb\'e...};
\draw[fill=poppurpleL,draw=poppurple,thick] (0.4,2.6) rectangle (10.1,4.9);
\node[align=left,anchor=north west] at (0.6,4.7) {\textbf{COMP\'ETENCES}};
\node[align=left,anchor=north west] at (0.6,4.2) {\footnotesize Techniques : informatique, conduite d'engins, permis...};
\node[align=left,anchor=north west] at (0.6,3.8) {\footnotesize Linguistiques : fran\c{c}ais, anglais, langues nationales};
\draw[fill=poppinkL,draw=poppink,thick] (0.4,0.9) rectangle (10.1,2.3);
\node[align=left,anchor=north west] at (0.6,2.1) {\textbf{CENTRES D'INT\'ER\^ET / DIVERS}};
\node[align=left,anchor=north west] at (0.6,1.65) {\footnotesize Sport, associatif, activit\'es valorisantes};
\node[align=right,anchor=east] at (10.0,0.45) {\footnotesize Fait \`a ..., le ... \quad Signature};
\end{tikzpicture}
\end{center}
\legende{Maquette d'un CV vierge : les cinq cadres étiquetés correspondent aux rubriques de la structure externe, de l'en-tête (en haut) à la signature (en bas).}
\end{popfigure}

\courssub{Ordonner les rubriques : une stratégie}

L'ordre présenté ci-dessus (formation puis expérience) est le plus courant pour un jeune diplômé, mais il n'est pas figé. Le principe directeur est stratégique.

\begin{methode}[Ordonner les rubriques de son CV]
\begin{enumerate}
\item Identifier ce que l'offre d'emploi valorise le plus : le diplôme, l'expérience ou une compétence précise.
\item Repérer dans son propre parcours l'atout le plus fort \emph{par rapport à ce poste}.
\item Placer cet atout en premier, juste après l'en-tête : un jeune sans expérience met « Formation » d'abord ; un technicien ayant déjà beaucoup travaillé met « Expérience professionnelle » d'abord.
\item Terminer par les rubriques secondaires (compétences, centres d'intérêt) et la signature.
\end{enumerate}
\end{methode}

\courssub{Lecture guidée de deux CV contrastés}

Comparons maintenant deux CV réels (anonymisés) reçus pour un même poste d'aide-comptable.

\textbf{CV A (bien structuré).} En-tête complet avec e-mail sobre ; rubrique « Formation » en ordre antichronologique avec dates et établissements ; rubrique « Expérience » mentionnant un stage dans un cabinet comptable de Yaoundé ; compétences informatiques précisées ; photo neutre ; signature et mention « Fait à Yaoundé, le... ». Le recruteur trouve chaque information en moins de dix secondes.

\textbf{CV B (désordonné).} Le nom apparaît en bas de page ; les diplômes sont mélangés aux stages dans un paragraphe unique ; le numéro de téléphone est absent ; l'e-mail est « beaugossedouala@... » ; aucune date n'accompagne les expériences ; la photo montre le candidat en tenue de plage. Le contenu réel du candidat est peut-être bon, mais il est illisible.

De cette comparaison, on dégage la classification suivante.

\begin{popfigure}
\begin{center}
\begin{tikzpicture}[grow=down,level distance=1.6cm,
every node/.style={draw,rounded corners,align=center,font=\small},
level 1/.style={sibling distance=6.5cm},
level 2/.style={sibling distance=3.2cm}]
\node[fill=popblueL,draw=popblue,thick]{\textbf{Rubriques du CV}}
  child { node[fill=popgreenL,draw=popgreen]{\textbf{Obligatoires}}
    child { node[fill=white, align=center]{En-t\^ete\\ (identit\'e, coordonn\'ees)} }
    child { node[fill=white, align=center]{Formation\\ (dipl\^omes)} }
    child { node[fill=white, align=center]{Exp\'erience\\ professionnelle} }
  }
  child { node[fill=poporangeL,draw=poporange]{\textbf{Facultatives}}
    child { node[fill=white, align=center]{Comp\'etences\\ d\'etaill\'ees} }
    child { node[fill=white, align=center]{Centres\\ d'int\'er\^et} }
    child { node[fill=white, align=center]{Photo, titre\\ du poste vis\'e} }
  };
\end{tikzpicture}
\end{center}
\legende{Arbre de classification des rubriques : les rubriques obligatoires (sans lesquelles le CV est incomplet) et les rubriques facultatives (qui valorisent mais ne sont pas exigées).}
\end{popfigure}

\begin{aretenir}[L'essentiel de la structure externe]
Un CV complet comprend : un \textbf{en-tête} (nom, coordonnées à jour, éventuellement poste visé), la rubrique \textbf{Formation} (ordre antichronologique), la rubrique \textbf{Expérience professionnelle} (dates, poste, structure), les \textbf{Compétences} (techniques et linguistiques), les \textbf{Centres d'intérêt}, et selon les usages la \textbf{photo}, la mention « Fait à..., le... » et la \textbf{signature}. L'ordre des rubriques se choisit en fonction du poste visé : ce qui valorise le plus le candidat vient en premier.
\end{aretenir}

\begin{exoresolu}[Reconstituer la structure d'un CV]
\textbf{Énoncé.} Voici, en désordre, les éléments du CV de Nadège, candidate à un poste de secrétaire à Bafoussam : (a) stage de deux mois dans une mairie, 2024 ; (b) nadege.fotso@mail.com ; (c) Probatoire technique G1, 2025, Lycée technique de Bafoussam ; (d) football en club ; (e) FOTSO Nadège ; (f) maîtrise de Word et Excel. Reconstituez les rubriques du CV dans le bon ordre.
\tcblower
\textbf{Solution détaillée.}
\begin{enumerate}
\item \textbf{En-tête} : (e) FOTSO Nadège, suivi des coordonnées (b) nadege.fotso@mail.com — adresse sobre et correcte.
\item \textbf{Formation} : (c) Probatoire technique G1, 2025, Lycée technique de Bafoussam — le diplôme le plus récent d'abord.
\item \textbf{Expérience professionnelle} : (a) 2024, stage de deux mois dans une mairie — date, nature du stage, structure.
\item \textbf{Compétences} : (f) maîtrise de Word et Excel — compétence technique informatique, précieuse pour un poste de secrétaire.
\item \textbf{Centres d'intérêt} : (d) football en club — activité sportive valorisante (esprit d'équipe).
\end{enumerate}
Le poste visé étant un emploi de secrétaire, on aurait aussi pu placer « Compétences » juste après l'en-tête, car la maîtrise de Word et Excel est l'atout décisif de la candidature.
\end{exoresolu}

\begin{exercice}[À vous de jouer]
Relevez dans un journal ou sur une annonce affichée une offre d'emploi de votre région. Dressez la liste des rubriques que devra contenir le CV d'un candidat à cette offre, dans l'ordre où vous les placeriez, et justifiez votre ordre en deux phrases. Le corrigé et la remédiation seront traités lors de la séance d'évaluation.
\end{exercice}

Ainsi structuré de l'extérieur, le CV possède son squelette. Mais un squelette ne suffit pas : il faut maintenant apprendre à \emph{rédiger} le contenu de chaque rubrique — choix des mots, formulations, verbes d'action. Ce sera l'objet de la section suivante, consacrée à la structure interne du CV.

\coursec{La structure interne du CV}

Dans la section précédente, nous avons étudié la \cle{structure externe} du CV : les grandes rubriques qui le composent (état civil, formation, expériences professionnelles, compétences, centres d'intérêt). Mais un CV ne se juge pas seulement à ce qu'il contient : il se juge aussi à la manière dont ces informations sont \emph{organisées et présentées à l'intérieur du document}. Un recruteur consacre en moyenne moins d'une minute à la première lecture d'un CV. Si le vôtre est confus, mal aligné ou trop dense, il sera écarté avant même d'être lu. C'est tout l'enjeu de la \cle{structure interne}.

\courssub{Qu'est-ce que la structure interne ?}

\begin{definition}[Structure interne du CV]
La \cle{structure interne} d'un CV désigne l'organisation \textbf{graphique} et \textbf{logique} des informations à l'intérieur du document : la mise en page (marges, alignements, aération), le choix typographique (polices, tailles, gras, puces), l'ordre de présentation des rubriques et des éléments (antichronologique ou thématique) et le style rédactionnel (concision, verbes d'action). Elle rend le document lisible, clair et convaincant.
\end{definition}

Pour bien comprendre, faisons une analogie : la structure externe est le \emph{plan d'une maison} (le salon, les chambres, la cuisine) ; la structure interne est \emph{l'aménagement de chaque pièce} (les meubles, la peinture, l'éclairage). Deux maisons peuvent avoir le même plan et pourtant l'une est agréable à visiter, l'autre non. De même, deux CV peuvent contenir les mêmes rubriques et produire des effets totalement différents sur le lecteur.

\courssub{La mise en page : aération, alignement, marges, hiérarchie visuelle}

\textbf{L'observation de départ.} Prenez deux feuilles contenant exactement le même texte : l'une remplie de bord à bord, sans espace, sans titre visible ; l'autre avec des marges, des blocs séparés par des blancs et des titres en gras. Laquelle avez-vous envie de lire ? La seconde, évidemment. Cette expérience simple montre que l'\textbf{œil du lecteur a besoin de repères visuels} pour se guider.

\begin{propriete}[Les quatre règles de la mise en page]
\begin{enumerate}
\item \textbf{L'aération} : le texte ne doit jamais occuper toute la page. Des espaces blancs séparent les rubriques et les paragraphes. Un CV « respire ».
\item \textbf{L'alignement} : les éléments de même nature sont alignés verticalement (les dates dans une même colonne à gauche, les intitulés au même niveau). On évite les alignements fantaisistes.
\item \textbf{Les marges} : on laisse des marges régulières (environ 2 cm) sur les quatre côtés de la feuille. Un texte qui touche les bords donne une impression d'étouffement.
\item \textbf{La hiérarchie visuelle} : les titres de rubriques (\textsc{Formation}, \textsc{Expérience professionnelle}) sont mis en évidence (gras, taille supérieure, éventuellement une couleur sobre) pour que le regard les repère immédiatement.
\end{enumerate}
\end{propriete}

\begin{popfigure}
\begin{tikzpicture}
% Maquette confuse (gauche)
\begin{scope}
\draw[fill=white,draw=popdark] (0,0) rectangle (5.6,7.6);
\node[font=\scriptsize\bfseries] at (2.8,7.2) {CV de Mballa Jean};
\node[font=\tiny,align=left,text width=5.2cm] at (2.8,5.9) {né le 12 mars 2004 à Douala célibataire tel 690000000 formation : 2023 BEPC collège de Deido 2024 probatoire 2025 baccalauréat G2 lycée technique expérience : 2022 stage menuiserie atelier Ndogpassi 2023 vendeur boutique 2024 stage comptabilité compétences : word excel soudure conduite langues : français anglais};
\node[font=\tiny,align=left,text width=5.2cm] at (2.8,2.6) {loisirs : football lecture musique voyage je suis dynamique motivé sérieux ponctuel disponible immédiatement et je cherche un emploi dans une entreprise sérieuse où je pourrai mettre mes compétences au service de la production et apprendre davantage auprès de professionnels expérimentés};
\node[font=\scriptsize\itshape,text=popink] at (2.8,0.5) {Version A : confuse, dense, sans repères};
\end{scope}
% Maquette aérée (droite)
\begin{scope}[xshift=7cm]
\draw[fill=white,draw=popdark] (0,0) rectangle (5.6,7.6);
\draw[fill=popblueL] (0.3,6.6) rectangle (5.3,7.4);
\node[font=\scriptsize\bfseries] at (2.8,7.0) {MBALLA Jean — Technicien};
\draw[popblue,line width=0.8pt] (0.3,6.35) -- (5.3,6.35);
\node[font=\tiny\bfseries,text=popblue,anchor=west] at (0.3,6.1) {FORMATION};
\node[font=\tiny,anchor=west] at (0.3,5.75) {2025 \quad Baccalauréat G2 — Lycée technique};
\node[font=\tiny,anchor=west] at (0.3,5.45) {2024 \quad Probatoire G2 — Lycée technique};
\draw[popblue,line width=0.8pt] (0.3,5.1) -- (5.3,5.1);
\node[font=\tiny\bfseries,text=popblue,anchor=west] at (0.3,4.85) {EXPÉRIENCE};
\node[font=\tiny,anchor=west] at (0.3,4.5) {2024 \quad Stage en comptabilité — Entreprise X};
\node[font=\tiny,anchor=west] at (0.3,4.2) {2023 \quad Vendeur — Boutique Y, Douala};
\draw[popblue,line width=0.8pt] (0.3,3.85) -- (5.3,3.85);
\node[font=\tiny\bfseries,text=popblue,anchor=west] at (0.3,3.6) {COMPÉTENCES};
\node[font=\tiny,anchor=west] at (0.3,3.25) {$\bullet$ Word, Excel \quad $\bullet$ Soudure};
\draw[popblue,line width=0.8pt] (0.3,2.9) -- (5.3,2.9);
\node[font=\tiny\bfseries,text=popblue,anchor=west] at (0.3,2.65) {LANGUES ET LOISIRS};
\node[font=\tiny,anchor=west] at (0.3,2.3) {Français, anglais — Football, lecture};
\node[font=\scriptsize\itshape,text=popgreen] at (2.8,0.5) {Version B : aérée, hiérarchisée, lisible};
\end{scope}
\end{tikzpicture}
\legende{Deux mises en page du même contenu. Version A : texte continu, aucune marge, aucun titre visible — le regard se perd. Version B : marges, rubriques signalées par des bandeaux et des filets, dates alignées à gauche — l'information se repère en quelques secondes.}
\end{popfigure}

\courssub{Le choix typographique}

La \cle{typographie} est l'art de composer le texte avec des caractères. Pour un CV, la règle d'or est la \textbf{sobriété}.

\begin{aretenir}[Règles typographiques du CV]
\begin{itemize}
\item Utiliser \textbf{une à deux polices} sobres au maximum (par exemple Arial, Calibri ou Times New Roman). Multiplier les polices donne une impression de désordre.
\item Choisir une \textbf{taille lisible} : 11 ou 12 points pour le texte courant, 14 à 16 points pour les titres de rubriques.
\item Employer le \textbf{gras} et les \textbf{puces} avec modération : le gras pour les titres et les éléments essentiels, les puces pour énumérer les tâches ou les compétences. Si tout est en gras, plus rien ne ressort.
\item Éviter les polices fantaisistes, les couleurs vives multiples et les majuscules abusives.
\end{itemize}
\end{aretenir}

\courssub{L'ordre de présentation : antichronologique ou thématique}

\textbf{L'ordre antichronologique} consiste à présenter les éléments \textbf{du plus récent au plus ancien}. C'est l'ordre le plus courant et le plus attendu par les recruteurs, car ce qui intéresse d'abord un employeur, c'est votre situation \emph{actuelle} : votre dernier diplôme, votre dernier emploi.

\textbf{L'ordre thématique} (ou CV par compétences) regroupe les informations \textbf{par domaines de compétences} (« Gestion et comptabilité », « Vente et accueil ») plutôt que par dates. Il est utile pour les profils à expériences variées ou discontinues, car il masque les périodes creuses et met en avant le savoir-faire.

\begin{popfigure}
\begin{tikzpicture}[font=\small]
% Frise
\draw[-{Stealth},thick,popdark] (0,0) -- (13,0);
\foreach \x/\a in {0.5/2018,3.5/2020,6.5/2022,9.5/2024,12.5/2026}{
  \draw[popdark] (\x,-0.12) -- (\x,0.12);
  \node[below] at (\x,-0.15) {\a};
}
% Ordre chronologique
\node[anchor=west,font=\small\bfseries,text=popblue] at (0,1.5) {Ordre chronologique (ancien $\rightarrow$ récent)};
\draw[-{Stealth},very thick,popblue] (0.5,0.9) -- (12.5,0.9);
\node[font=\scriptsize,text=popblue] at (6.5,0.55) {lecture dans le sens du temps};
% Ordre antichronologique
\node[anchor=west,font=\small\bfseries,text=poporange] at (0,-1.6) {Ordre antichronologique (récent $\rightarrow$ ancien)};
\draw[-{Stealth},very thick,poporange] (12.5,-1.0) -- (0.5,-1.0);
\node[font=\scriptsize,text=poporange] at (6.5,-0.65) {ordre recommandé pour le CV : on commence par 2026};
\end{tikzpicture}
\legende{Frise temporelle 2018--2026. En haut, l'ordre chronologique suit le sens du temps. En bas, l'ordre antichronologique, utilisé dans le CV, part de l'élément le plus récent et remonte vers le plus ancien.}
\end{popfigure}

\begin{methode}[Rédiger les rubriques « Formation » et « Expérience »]
\begin{enumerate}
\item Classer les éléments dans l'\textbf{ordre antichronologique} : le plus récent en premier.
\item Pour chaque élément, respecter le schéma : \textbf{date — intitulé — lieu/structure}.
\item Aligner toutes les dates dans la même colonne, à gauche.
\item Rédiger les tâches avec des verbes d'action à l'infinitif ou au participe passé.
\item Vérifier la cohérence des dates (pas de chevauchement impossible, pas de « trou » injustifié).
\end{enumerate}
\end{methode}

\begin{exemplebox}[La rubrique « Expérience » bien construite]
\begin{itemize}
\item \textbf{2025} — Stage de comptabilité — Entreprise SOTRAC, Douala : \emph{saisir les factures ; classer les pièces comptables}.
\item \textbf{2024} — Vendeur — Boutique Électronique Express, Deïdo : \emph{assurer l'accueil des clients ; gérer la caisse}.
\item \textbf{2023} — Stage d'ouvrier — Atelier de menuiserie Ndogpassi : \emph{participer à la fabrication de meubles ; réparer les chaises}.
\end{itemize}
On remarque : dates alignées, ordre antichronologique, verbes d'action à l'infinitif.
\end{exemplebox}

\courssub{La concision : phrases nominales et verbes d'action}

Un CV n'est pas une rédaction : on n'y écrit pas de phrases complètes avec sujet et subordonnées. On privilégie la \cle{phrase nominale} (sans verbe conjugué) et le \cle{verbe d'action} à l'infinitif ou au participe passé.

\begin{center}
\begin{tabularx}{\linewidth}{|X|X|}
\hline
\rowcolor{popblueL} \textbf{À éviter (phrase lourde)} & \textbf{À préférer (concis)} \\
\hline
« J'étais chargé de la gestion de la caisse. » & « Gérer la caisse » ou « Caisse gérée » \\
\hline
« Je m'occupais de l'accueil des clients. » & « Assurer l'accueil des clients » \\
\hline
« J'ai participé à la réparation des machines. » & « Réparer les machines » ou « Réparations effectuées » \\
\hline
\end{tabularx}
\end{center}

\begin{savaistu}[Le pouvoir des verbes d'action]
Les verbes d'action comme \emph{organiser}, \emph{gérer}, \emph{concevoir}, \emph{réparer}, \emph{installer}, \emph{former} rendent les expériences plus dynamiques que les tournures passives (« j'étais chargé de... », « il m'a été demandé de... »). Ils présentent le candidat comme quelqu'un qui \textbf{agit}, et non qui subit. Les recruteurs y sont très sensibles : un même stage décrit avec des verbes d'action paraît plus riche et plus professionnel.
\end{savaistu}

\courssub{L'honnêteté, la vérifiabilité et l'unité graphique}

\begin{attention}[Ne jamais mentir dans un CV]
Erreur fréquente et grave : \textbf{mentir sur un diplôme ou une expérience} (inventer un stage, gonfler une durée, se prétendre titulaire d'un diplôme non obtenu). Toute fausse déclaration peut entraîner le \textbf{rejet du dossier} lors de la vérification des références, ou même un \textbf{licenciement} si la fraude est découverte après l'embauche. Les dates doivent être cohérentes et les diplômes réels : un recruteur peut demander les attestations.
\end{attention}

L'\textbf{unité graphique} est la dernière exigence : toutes les puces doivent être identiques, toutes les dates alignées et écrites de la même façon (« 2024 » partout, pas « 2024 » puis « deux mille vingt-quatre »), et l'\textbf{orthographe irréprochable}. Une faute d'orthographe dans un CV est souvent éliminatoire, car elle suggère le manque de rigueur.

\begin{exemplebox}[Un chiffre à retenir]
Un CV de plus de \textbf{2 pages} pour un jeune diplômé est jugé \textbf{trop long} par la majorité des recruteurs. Pour un élève de Terminale technique, \textbf{une seule page} suffit largement : mieux vaut une page claire et dense en informations utiles que deux pages diluées.
\end{exemplebox}

\courssub{Activité : comparer deux mises en page}

\begin{exoresolu}[Quelle mise en page est la plus efficace ?]
Énoncé. On donne aux élèves les deux versions du CV de Mballa Jean reproduites dans la figure ci-dessus : la version A (texte continu, sans marges, sans titres) et la version B (rubriques signalées, dates alignées, puces). Questions : 1) Relever trois différences de présentation entre les deux versions. 2) Dire quelle version est la plus efficace et justifier en trois arguments.
\tcblower
Solution détaillée.
\textbf{1) Différences relevées :} (a) la version B présente des titres de rubriques en évidence (bandeaux et filets), la version A n'en a pas ; (b) la version B aligne les dates dans une colonne à gauche, la version A noie les dates dans le texte ; (c) la version B utilise des puces et des espaces blancs, la version A est un bloc compact sans marges.
\textbf{2) Justification :} la version B est la plus efficace. Premier argument : la \emph{hiérarchie visuelle} permet au recruteur de repérer chaque rubrique en quelques secondes. Deuxième argument : l'\emph{aération} et les marges rendent la lecture agréable et invitent à lire jusqu'au bout. Troisième argument : les \emph{dates alignées} et l'ordre antichronologique permettent de suivre immédiatement le parcours du candidat. Conclusion : à contenu identique, c'est la structure interne (mise en page, ordre, concision) qui fait la différence entre un CV lu et un CV écarté.
\end{exoresolu}

\begin{exercice}[À vous de jouer]
Voici trois lignes extraites d'un CV : « 2022 — stage en électricité — ENEO, Yaoundé » ; « 2025 — Baccalauréat G2 — Lycée technique de Nkongsamba » ; « 2024 — Probatoire G2 — Lycée technique de Nkongsamba ». 1) Remettez ces éléments dans l'ordre correct pour un CV. 2) Réécrivez la phrase « J'ai été chargé de l'installation des câbles électriques » sous forme concise. 3) Citez deux règles d'unité graphique à respecter.
\end{exercice}

\begin{corrige}[Exercice : mise en ordre et concision]
1) Ordre antichronologique : d'abord « 2025 — Baccalauréat G2 — Lycée technique de Nkongsamba », ensuite « 2024 — Probatoire G2 — Lycée technique de Nkongsamba », enfin « 2022 — stage en électricité — ENEO, Yaoundé » (du plus récent au plus ancien).
2) Forme concise : « Installer les câbles électriques » (infinitif) ou « Câbles électriques installés » (participe passé).
3) Deux règles d'unité graphique : utiliser des puces identiques dans tout le document ; aligner toutes les dates dans la même colonne et les écrire de la même manière. On peut ajouter : soigner une orthographe irréprochable.
\end{corrige}

\begin{aretenir}[L'essentiel de la structure interne]
La structure interne du CV repose sur cinq piliers : une \textbf{mise en page aérée} (marges, alignements, titres en évidence), une \textbf{typographie sobre} (une à deux polices, gras et puces avec modération), un \textbf{ordre antichronologique} (le plus récent d'abord) ou thématique selon le profil, une \textbf{rédaction concise} (phrases nominales, verbes d'action) et une \textbf{honnêteté totale} (dates cohérentes, diplômes réels, unité graphique et orthographe irréprochables). Un jeune diplômé tient son CV sur une page.
\end{aretenir}

\coursec{Langue : les valeurs du subjonctif et du conditionnel}

Après avoir maîtrisé la \textbf{structure interne du CV} (rubriques, mise en page, vocabulaire technique), il est indispensable de soigner la \textbf{langue} pour donner à votre candidature la crédibilité et la nuance attendues par un recruteur. Un CV et une lettre de motivation ne sont pas de simples listes de faits ; ils expriment une \textbf{intention}, un \textbf{projet}, une \textbf{disponibilité}. C'est là qu'interviennent deux modes majeurs : le \cle{subjonctif} (mode de la subjectivité, du souhait, de la nécessité) et le \cle{conditionnel} (mode de l'hypothèse, de la politesse, de l'atténuation). Maîtriser leurs valeurs, c'est éviter les maladresses qui peuvent coûter un entretien.

\courssub{Le subjonctif : mode de l'éventuel, du souhait, de l'obligation et du sentiment}

\begin{definition}[Le subjonctif]
Mode verbal qui exprime un fait \textbf{envisagé, souhaité, ordonné, ressenti ou mis en doute}, et non un fait accompli ou certain. Il est le plus souvent employé dans une \textbf{proposition subordonnée} introduite par la conjonction \og que \fg{}, dont le verbe de la principale exprime une subjectivité (volonté, émotion, nécessité, doute).
\end{definition}

\textbf{Intuition.}\par
L'indicatif dit : \og C'est vrai / C'est fait \fg{} (\emph{Je viens}). Le subjonctif dit : \og Je veux que ce soit vrai / Il faut que ce soit fait \fg{} (\emph{Je veux que tu viennes}). Dans un dossier de candidature, vous ne décrivez pas seulement ce qui \emph{est} (indicatif), vous formulez ce qui \emph{doit être} ou ce que vous \emph{souhaitez} (subjonctif).

\begin{methode}[Repérer les contextes déclencheurs du subjonctif]
Au Probatoire / Baccalauréat technique et dans la rédaction professionnelle, le subjonctif est obligatoire après les constructions suivantes (liste non exhaustive mais couvrante pour le programme) :
\begin{enumerate}
    \item \textbf{Volonté / Souhait / Ordre / Conseil :} \emph{je veux que, je souhaite que, je demande que, je préfère que, j'aimerais que.}
    \item \textbf{Nécessité / Obligation (tournure impersonnelle) :} \emph{il faut que, il est nécessaire que, il est indispensable que, il est urgent que.}
    \item \textbf{Sentiment / Émotion / Appréciation :} \emph{je suis content que, je regrette que, c'est dommage que, c'est bien que, je trouve étonnant que.}
    \item \textbf{Doute / Négation / Incertitude :} \emph{je doute que, il n'est pas sûr que, il n'est pas certain que, je ne pense pas que, je ne crois pas que.}
    \item \textbf{Finalité (but) :} \emph{pour que, afin que, de peur que, de sorte que.}
    \item \textbf{Concession (opposition) :} \emph{bien que, quoique.}
\end{enumerate}
\end{methode}

\begin{exemplebox}[Exemples en contexte de candidature]
\begin{itemize}
    \item \textbf{Volonté :} \og \emph{Je souhaite vivement que mon dossier \textbf{soit} examiné avec attention.} \fg{} (subjonctif présent, voix passive)
    \item \textbf{Obligation :} \og \emph{Il est indispensable que le candidat \textbf{maîtrise} l'anglais technique.} \fg{} (subjonctif présent)
    \item \textbf{Finalité :} \og \emph{Je me forme \textbf{afin que} mes compétences \textbf{soient} reconnues.} \fg{} (subjonctif présent)
    \item \textbf{Concession :} \og \emph{\textbf{Bien que} je \textbf{sois} débutant, je suis très motivé.} \fg{} (subjonctif présent)
    \item \textbf{Sentiment :} \og \emph{Je suis ravi que cette opportunité \textbf{se présente}.} \fg{}
\end{itemize}
\end{exemplebox}

\begin{attention}[Pièges fréquents au Bac et en rédaction professionnelle]
\begin{itemize}
    \item \textbf{Erreur :} \og Après que j'aie obtenu mon diplôme... \fg{} $\rightarrow$ \textbf{Correction :} \og \textbf{Après que j'ai obtenu} mon diplôme... \fg{} (\textbf{Indicatif} obligatoire après \og après que \fg{}, car le fait est présenté comme accompli).
    \item \textbf{Erreur :} \og Il faut que je \textbf{viens} demain. \fg{} $\rightarrow$ \textbf{Correction :} \og Il faut que je \textbf{vienne} demain. \fg{} (Subjonctif présent : \emph{que je vienne}, \emph{que tu viennes}, \emph{qu'il vienne}, \emph{que nous venions}, \emph{que vous veniez}, \emph{qu'ils viennent}).
    \item \textbf{Erreur :} \og Je pense qu'il \textbf{viendrait}. \fg{} $\rightarrow$ \textbf{Correction :} \og Je pense qu'il \textbf{vient} \fg{} (indicatif, certitude) \textbf{OU} \og Je ne pense pas qu'il \textbf{vienne} \fg{} (subjonctif, doute).
    \item \textbf{Accord du participe passé (subjonctif passé) :} \og Il faut que tu \textbf{sois parti} avant midi. \fg{} (auxiliaire \emph{être} au subjonctif + participe accordé avec le sujet \emph{tu}).
\end{itemize}
\end{attention}

\courssub{Le conditionnel : hypothèse, politesse et information non confirmée}

\begin{definition}[Le conditionnel]
Mode verbal qui exprime un fait \textbf{soumis à une condition} (hypothèse), un \textbf{souhait atténué} (politesse) ou une \textbf{information non confirmée} (valeur journalistique). Il possède deux temps : le \textbf{conditionnel présent} (valeur principale) et le \textbf{conditionnel passé} (antériorité dans l'hypothèse).
\end{definition}

\textbf{Intuition.}\par
L'indicatif affirme : \og Je \textbf{viens} \fg{}. Le conditionnel module : \og Je \textbf{viendrais} \emph{si j'étais libre} \fg{} (hypothèse) ou \og Je \textbf{viendrais} \emph{volontiers} \fg{} (politesse, souhait atténué). Dans un CV, une lettre ou un entretien, le conditionnel est la marque de l'\textbf{humilité professionnelle} et de la \textbf{projection dans l'avenir}.

\begin{propriete}[Les trois valeurs majeures du conditionnel présent]
\begin{enumerate}
    \item \textbf{Valeur hypothétique (la plus grammaticale) :} le conditionnel exprime une action qui se réaliserait \textbf{si une condition était remplie}. La condition est souvent introduite par \og si \fg{} suivi de l'\textbf{imparfait de l'indicatif}. Exemple : \og \emph{Si j'\textbf{obtenais} ce poste, je \textbf{m'investirais} pleinement.} \fg{}
    \item \textbf{Valeur de politesse / atténuation (la plus fréquente en candidature) :} le conditionnel adoucit une demande, un souhait ou une opinion pour ne pas paraître exigeant. Exemples : \og \emph{Je \textbf{souhaiterais} vous exposer mon parcours} \fg{} ; \og \emph{J'\textbf{aimerais} rejoindre votre équipe} \fg{} ; \og \emph{\textbf{Pourriez}-vous me préciser les modalités ?} \fg{}
    \item \textbf{Valeur journalistique (information non confirmée) :} le conditionnel rapporte une rumeur ou une information non vérifiée. Exemple : \og \emph{Le directeur \textbf{aurait} annoncé une restructuration.} \fg{} (conditionnel passé : \emph{aurait annoncé} = on le dit, mais ce n'est pas confirmé).
\end{enumerate}
\end{propriete}

\begin{exemplebox}[La politesse en contexte de candidature]
Comparez les registres :
\begin{itemize}
    \item Trop direct : \og \emph{Je \textbf{veux} ce poste.} \fg{} $\rightarrow$ Atténué : \og \emph{Je \textbf{souhaiterais} occuper ce poste.} \fg{}
    \item Trop direct : \og \emph{\textbf{Pouvez}-vous me répondre vite ?} \fg{} $\rightarrow$ Atténué : \og \emph{\textbf{Pourriez}-vous me répondre dès que possible ?} \fg{}
    \item Trop affirmatif : \og \emph{C'\textbf{est} un atout pour votre structure.} \fg{} $\rightarrow$ Atténué : \og \emph{Cela \textbf{serait} un atout pour votre structure.} \fg{}
\end{itemize}
\end{exemplebox}

\begin{methode}[Concordance des temps : Si + Imparfait $\rightarrow$ Conditionnel présent]
C'est la structure type de l'hypothèse \textbf{irréelle du présent} (possible mais non réalisée au moment où l'on parle).
\begin{center}
\begin{tabular}{|c|c|c|}\hline
\textbf{Proposition introduite par \og si \fg{}} & \textbf{Proposition principale} & \textbf{Valeur} \\ \hline
\emph{Si} + \textbf{Imparfait de l'indicatif} & \textbf{Conditionnel présent} & Hypothèse présente ou future, possible \\ \hline
\emph{Si} + \textbf{Plus-que-parfait de l'indicatif} & \textbf{Conditionnel passé} & Hypothèse passée, non réalisée (regret, reproche) \\ \hline
\end{tabular}
\end{center}
\textbf{Attention :} après \og si \fg{} exprimant la condition, on n'emploie \textbf{jamais} le conditionnel : on dit \og Si j'\textbf{avais} le temps \fg{} et non \og Si j'\textbf{aurais} le temps \fg{}.
\end{methode}

\begin{exemplebox}[Application à la candidature]
\begin{itemize}
    \item \og \emph{Si j'\textbf{avais} le permis poids lourd (imparfait), je \textbf{postulerais} (conditionnel présent) pour le poste de livreur.} \fg{}
    \item \og \emph{Si j'\textbf{avais su} (plus-que-parfait) l'existence de ce concours, j'\textbf{aurais candidaté} (conditionnel passé).} \fg{} (regret)
\end{itemize}
\end{exemplebox}

\courssub{Choisir entre subjonctif et conditionnel : l'arbre de décision}

La confusion est fréquente car les deux modes expriment le \og non-réel \fg{}. Le critère décisif est la \textbf{structure syntaxique} et la \textbf{nature du verbe introducteur}.

\begin{popfigure}
\begin{tikzpicture}[
    node distance=1.4cm and 1.6cm,
    decision/.style={diamond, draw=popblue, thick, fill=popblue!10, text width=3.6cm, align=center, inner sep=3pt},
    process/.style={rectangle, draw=popgreen, thick, fill=popgreen!10, text width=4.2cm, align=center, inner sep=4pt, rounded corners},
    start/.style={ellipse, draw=poporange, thick, fill=poporange!10, text width=3cm, align=center, inner sep=4pt},
    arrow/.style={-Stealth, thick, popdark},
    yesno/.style={midway, fill=white, inner sep=2pt, font=\small, text=popdark}
]
\node[start] (start) {Intention : exprimer...};
\node[decision, below of=start] (q1) {Y a-t-il un verbe principal + \og que \fg{} ?};
\node[process, below left=1.6cm and 0.8cm of q1, align=center] (subj){\textbf{SUBJONCTIF}\\ Volonté, obligation, sentiment, doute, finalité, concession\\ \emph{Ex. : Il faut que je vienne}};
\node[decision, below right=1.6cm and 0.8cm of q1] (q2) {Est-ce une hypothèse en \og si \fg{} ?};
\node[process, below left=1.6cm and 0.6cm of q2, align=center] (condhyp){\textbf{CONDITIONNEL}\\ Valeur hypothétique\\ \emph{Ex. : Si j'avais le temps, je le ferais}};
\node[process, below right=1.6cm and 0.6cm of q2, align=center] (condpol){\textbf{CONDITIONNEL}\\ Valeur de politesse / atténuation\\ \emph{Ex. : Je souhaiterais... Pourriez-vous...}};

\draw[arrow] (start) -- (q1);
\draw[arrow] (q1) -- node[yesno, left] {OUI} (subj);
\draw[arrow] (q1) -- node[yesno, right] {NON} (q2);
\draw[arrow] (q2) -- node[yesno, left] {OUI} (condhyp);
\draw[arrow] (q2) -- node[yesno, right] {NON} (condpol);
\end{tikzpicture}
\legende{Arbre de choix : quel mode employer selon l'intention de communication ?}
\end{popfigure}

\begin{exoresolu}[Analyse de phrases : subjonctif ou conditionnel ?]
\textbf{Énoncé :} Complétez avec le mode et le temps qui conviennent, puis justifiez.
\begin{enumerate}
    \item Il est urgent que le technicien \_\_\_\_\_\_\_\_\_\_ (intervenir) au plus vite.
    \item \_\_\_\_\_\_\_\_\_\_ (avoir)-vous l'amabilité de me recevoir ?
    \item Je doute qu'il \_\_\_\_\_\_\_\_\_\_ (connaître) ce logiciel spécifique.
    \item Si le salaire \_\_\_\_\_\_\_\_\_\_ (être) négociable, j'accepterais.
    \item Bien que le travail \_\_\_\_\_\_\_\_\_\_ (être) difficile, je le ferai.
\end{enumerate}

\tcblower
\textbf{Solution détaillée :}
\begin{enumerate}
    \item \textbf{intervienne} (subjonctif présent). \textbf{Justification :} \og Il est urgent que \fg{} exprime une \textbf{nécessité} $\rightarrow$ subjonctif. Verbe \emph{intervenir} : \emph{qu'il intervienne}.
    \item \textbf{Auriez}-vous (conditionnel présent). \textbf{Justification :} demande de service directe $\rightarrow$ \textbf{politesse / atténuation} $\rightarrow$ conditionnel. Avec inversion du sujet : \emph{Auriez-vous}.
    \item \textbf{connaisse} (subjonctif présent). \textbf{Justification :} \og Je doute que \fg{} exprime le \textbf{doute} $\rightarrow$ subjonctif. Verbe \emph{connaître} : \emph{qu'il connaisse}.
    \item \textbf{était} (imparfait de l'indicatif). \textbf{Justification :} dans la proposition introduite par \og si \fg{} (condition), on emploie l'\textbf{imparfait}, jamais le conditionnel ; la principale \og j'accepterais \fg{} porte le conditionnel.
    \item \textbf{soit} (subjonctif présent). \textbf{Justification :} \og Bien que \fg{} exprime la \textbf{concession} $\rightarrow$ subjonctif. Verbe \emph{être} (irrégulier) : \emph{qu'il soit}.
\end{enumerate}
\end{exoresolu}

\courssub{Tableau récapitulatif : modes, valeurs et exemples professionnels}

\begin{center}
\begin{tabularx}{\linewidth}{|>{\bfseries}l|X|X|}\hline
\rowcolor{popblueL}
\textbf{Mode / Valeur} & \textbf{Contexte déclencheur} & \textbf{Exemple en candidature} \\ \hline
Subjonctif présent : souhait, volonté & \emph{vouloir, souhaiter, préférer, demander} + \og que \fg{} & \og \emph{Je souhaite que mon profil \textbf{corresponde} à vos attentes.} \fg{} \\ \hline
Subjonctif présent : obligation, nécessité & \emph{il faut que, il est indispensable que} & \og \emph{Il faut que le candidat \textbf{maîtrise} les normes de sécurité.} \fg{} \\ \hline
Subjonctif présent : finalité & \emph{pour que, afin que, de sorte que} & \og \emph{Je me perfectionne \textbf{afin que} mes compétences \textbf{soient} reconnues.} \fg{} \\ \hline
Subjonctif présent : concession & \emph{bien que, quoique} & \og \emph{\textbf{Bien que} je \textbf{sois} junior, ma motivation est totale.} \fg{} \\ \hline
Subjonctif présent : sentiment & \emph{être content que, regretter que, c'est dommage que} & \og \emph{Je suis ravi que cette offre \textbf{corresponde} à mon projet.} \fg{} \\ \hline
Conditionnel présent : hypothèse & \og si \fg{} + imparfait $\rightarrow$ conditionnel & \og \emph{Si j'\textbf{intégrais} votre service, j'\textbf{apporterais} ma rigueur.} \fg{} \\ \hline
Conditionnel présent : politesse & \emph{souhaiter, aimer, pouvoir, vouloir} (sans \og que \fg{}) & \og \emph{Je \textbf{souhaiterais} vous rencontrer. \textbf{Pourriez}-vous m'indiquer la suite ?} \fg{} \\ \hline
Conditionnel passé : regret, reproche & \og si \fg{} + plus-que-parfait $\rightarrow$ conditionnel passé & \og \emph{Si j'\textbf{avais relu} ma lettre, je n'\textbf{aurais pas fait} cette faute.} \fg{} \\ \hline
Conditionnel passé : information non confirmée & verbes déclaratifs : \emph{dire, annoncer, affirmer} & \og \emph{L'entreprise \textbf{aurait recruté} cinquante techniciens cette année.} \fg{} \\ \hline
\end{tabularx}
\end{center}

\courssub{Application au dossier de candidature : mise en situation}

\begin{aretenir}[Règles d'or pour le CV et la lettre de motivation]
\begin{enumerate}
    \item \textbf{CV (faits accomplis) :} utilisez l'\textbf{indicatif} (passé composé, imparfait) pour décrire vos expériences réelles. \emph{Ex. : \og J'\textbf{ai géré} une équipe de cinq personnes. \fg{}}
    \item \textbf{Lettre de motivation (projection, souhait) :} utilisez le \textbf{conditionnel de politesse} pour votre demande d'entretien. \emph{Ex. : \og Je \textbf{serais} honoré de vous présenter mes motivations. \fg{}}
    \item \textbf{Lettre / entretien (exigences du poste, objectifs) :} utilisez le \textbf{subjonctif} après \og il faut que, je souhaite que, pour que \fg{}. \emph{Ex. : \og Je souhaite que mes compétences \textbf{soient} mises à profit. \fg{}}
    \item \textbf{Entretien (réponse hypothétique) :} utilisez \og si + imparfait $\rightarrow$ conditionnel \fg{}. \emph{Ex. : \og \textbf{Si} on me \textbf{confiait} ce dossier, je \textbf{prioriserais} la sécurité. \fg{}}
\end{enumerate}
\end{aretenir}

\begin{exemplebox}[Exemple chiffré : l'impact du conditionnel de politesse]
Lors d'une simulation d'entretien organisée en classe de Terminale technique (grille d'évaluation sur 20 points, item \og Communication professionnelle \fg{}), les candidats utilisant le \textbf{conditionnel de politesse} (\og \emph{Je \textbf{souhaiterais} évoluer vers un poste de chef d'équipe} \fg{}) obtiennent en moyenne \textbf{16/20}, contre \textbf{10/20} pour ceux qui emploient l'injonction directe (\og \emph{Je \textbf{veux} ce poste} \fg{}) ou l'affirmation sans nuance (\og \emph{Je \textbf{serai} chef d'équipe} \fg{}). Soit un écart de \textbf{6 points}, c'est-à-dire une note supérieure de $60\,\%$ : le conditionnel signale l'adaptabilité et le respect de la hiérarchie.
\end{exemplebox}

\courssub{Exercices d'entraînement (type Probatoire / Baccalauréat technique)}

\begin{exercice}[Transformation : indicatif $\rightarrow$ conditionnel de politesse]
Transformez les phrases suivantes en mettant le verbe en gras au \textbf{conditionnel présent} (valeur de politesse / atténuation).
\begin{enumerate}
    \item Je \textbf{veux} vous parler de mon stage.
    \item Je \textbf{peux} commencer dès lundi.
    \item Il \textbf{faut} que je parte tôt vendredi. (Attention : structure \og il faut que \fg{} $\rightarrow$ subjonctif ! Piège.)
    \item J'\textbf{aimerai} avoir vos nouvelles. (Corrigez aussi le temps.)
    \item Nous \textbf{espérons} que vous nous recevrez.
\end{enumerate}
\end{exercice}

\begin{exercice}[Complétion : subjonctif ou indicatif ?]
Conjuguez le verbe entre parenthèses au mode et au temps qui conviennent (subjonctif présent ou indicatif). Justifiez en une phrase.
\begin{enumerate}
    \item Il est possible que le directeur \_\_\_\_\_\_\_\_\_\_ (recevoir) les candidatures lui-même.
    \item Après que le jury \_\_\_\_\_\_\_\_\_\_ (délibérer), les résultats seront affichés.
    \item Je cherche un poste qui me \_\_\_\_\_\_\_\_\_\_ (permettre) de voyager.
    \item Bien que ce métier \_\_\_\_\_\_\_\_\_\_ (être) fatigant, je l'aime.
    \item Je ne pense pas qu'il \_\_\_\_\_\_\_\_\_\_ (avoir) le diplôme requis.
\end{enumerate}
\end{exercice}

\begin{corrige}[Exercice : Transformation]
\begin{enumerate}
    \item Je \textbf{voudrais} vous parler de mon stage. (\emph{Vouloir} $\rightarrow$ \emph{je voudrais} ; on peut aussi écrire \emph{je souhaiterais}).
    \item Je \textbf{pourrais} commencer dès lundi. (\emph{Pouvoir} $\rightarrow$ \emph{je pourrais}).
    \item Il \textbf{faut que je parte} tôt vendredi. (\textbf{Piège} : \og il faut que \fg{} appelle le \textbf{subjonctif}, pas le conditionnel. \emph{Partir} $\rightarrow$ \emph{que je parte}.)
    \item J'\textbf{aimerais} avoir vos nouvelles. (\og J'aimerai \fg{} est un futur, incorrect ici ; la politesse exige le conditionnel : \emph{j'aimerais}.)
    \item Nous \textbf{souhaiterions} que vous nous \textbf{receviez}. (Pour la politesse, on préfère \emph{souhaiter} au conditionnel, suivi du subjonctif : \emph{que vous receviez}.)
\end{enumerate}
\end{corrige}

\begin{corrige}[Exercice : Complétion]
\begin{enumerate}
    \item \textbf{reçoive} (subjonctif présent). \textbf{Justification :} \og il est possible que \fg{} exprime la \textbf{possibilité, l'incertitude} $\rightarrow$ subjonctif.
    \item \textbf{a délibéré} (passé composé de l'indicatif) ou \textbf{aura délibéré} (futur antérieur). \textbf{Justification :} \og après que \fg{} appelle \textbf{l'indicatif} (fait présenté comme accompli), jamais le subjonctif.
    \item \textbf{permette} (subjonctif présent). \textbf{Justification :} la subordonnée relative porte sur un antécédent \textbf{indéfini, recherché} (\og un poste qui... \fg{}, encore hypothétique) $\rightarrow$ subjonctif.
    \item \textbf{soit} (subjonctif présent). \textbf{Justification :} \og bien que \fg{} exprime la \textbf{concession} $\rightarrow$ subjonctif. Verbe \emph{être}, irrégulier : \emph{qu'il soit}.
    \item \textbf{ait} (subjonctif présent). \textbf{Justification :} \og je ne pense pas que \fg{} (opinion niée) exprime le \textbf{doute} $\rightarrow$ subjonctif. Verbe \emph{avoir} : \emph{qu'il ait}.
\end{enumerate}
\end{corrige}

\begin{attention}[Dernier rappel avant les activités d'intégration]
Ne confondez pas :
\begin{itemize}
    \item \og \textbf{Si j'aurais} su... \fg{} $\rightarrow$ \textbf{INTERDIT}. On écrit : \og \textbf{Si j'avais} su (imparfait), \textbf{j'aurais} agi autrement (conditionnel passé). \fg{}
    \item \og \textbf{Après que j'aie} obtenu... \fg{} $\rightarrow$ \textbf{INTERDIT}. On écrit : \og \textbf{Après que j'ai} obtenu (indicatif, passé composé). \fg{}
    \item \og \textbf{Il faut que je viens}... \fg{} $\rightarrow$ \textbf{INTERDIT}. On écrit : \og \textbf{Il faut que je vienne} (subjonctif présent). \fg{}
\end{itemize}
La maîtrise de ces distinctions est un marqueur fort de \cle{compétence communicationnelle} attendue au Probatoire / Baccalauréat technique.
\end{attention}

\coursec{Textes iconiques et adaptation du CV à une offre d'embauche}

Dans la section précédente, nous avons étudié les valeurs du subjonctif et du conditionnel, ces modes qui permettent d'exprimer le souhait, l'hypothèse ou la politesse — outils précieux pour rédiger une lettre de motivation ou formuler ses prétentions. Mais un CV ne se réduit pas à des phrases : c'est aussi un document \emph{visuel}, fait d'images, de signes et de mise en page. Et surtout, un CV n'est jamais envoyé dans le vide : il répond à une \cle{offre d'embauche} précise. C'est à ces deux compétences — lire les textes iconiques et adapter son CV à une situation d'embauche — qu'est consacrée cette section.

\courssub{Lecture des textes iconiques}

\textbf{Qu'est-ce qu'un texte iconique ?}\par

Observez le CV d'un camarade : à côté des rubriques rédigées, vous voyez une photographie d'identité, un petit dessin de téléphone devant un numéro, une enveloppe devant une adresse électronique, parfois le logo d'une entreprise ou d'un lycée. Ces images ne décorent pas : elles \emph{informent}. Le dessin du téléphone dispense d'écrire le mot « Téléphone » ; la photographie permet au recruteur d'identifier le candidat.

\begin{definition}[Texte iconique]
Un \cle{texte iconique} est un document où l'\cle{image} — photographie, pictogramme, logo, schéma — porte une partie de l'information, seule ou en complément du texte écrit. Dans un CV, les principaux éléments iconiques sont la \cle{photographie d'identité}, les \cle{pictogrammes} (téléphone, enveloppe, maison) et les \cle{logos} des structures fréquentées (établissements, entreprises).
\end{definition}

\begin{popfigure}
\begin{center}
\begin{tikzpicture}[scale=1]
% Pictogramme téléphone
\draw[popblue, line width=1.2pt] (0,0.1) -- (0,0.9) arc[start angle=180, end angle=0, radius=0.3] -- (0.6,0.1);
\draw[popblue, fill=popblueL, line width=1pt] (-0.12,0.05) rectangle (0.12,0.35);
\draw[popblue, fill=popblueL, line width=1pt] (0.48,0.05) rectangle (0.72,0.35);
\node[below, popdark] at (0.3,-0.15) {\small Téléphone};
\node[below, popmuted] at (0.3,-0.6) {\small 6 90 00 00 00};
% Pictogramme enveloppe
\begin{scope}[xshift=3.5cm]
\draw[poporange, line width=1.2pt] (0,0) rectangle (1.2,0.8);
\draw[poporange, line width=1.2pt] (0,0.8) -- (0.6,0.3) -- (1.2,0.8);
\node[below, popdark] at (0.6,-0.15) {\small E-mail};
\node[below, popmuted] at (0.6,-0.6) {\small nom@mail.com};
\end{scope}
% Pictogramme maison
\begin{scope}[xshift=7cm]
\draw[popgreen, line width=1.2pt] (0,0) rectangle (1,0.6);
\draw[popgreen, line width=1.2pt] (-0.15,0.6) -- (0.5,1.1) -- (1.15,0.6);
\draw[popgreen, line width=1pt] (0.4,0) rectangle (0.6,0.35);
\node[below, popdark] at (0.5,-0.15) {\small Adresse};
\node[below, popmuted] at (0.5,-0.6) {\small Bonabéri, Douala};
\end{scope}
% Pictogramme photo
\begin{scope}[xshift=10.5cm]
\draw[poppurple, line width=1.2pt] (0,0) rectangle (1,1.2);
\draw[poppurple, fill=poppurpleL] (0.5,0.75) circle (0.22);
\draw[poppurple, fill=poppurpleL] (0.5,0.2) circle (0.3);
\node[below, popdark] at (0.5,-0.15) {\small Photo d'identité};
\end{scope}
\end{tikzpicture}
\end{center}
\legende{Les pictogrammes usuels d'un CV : chaque petit dessin remplace une étiquette écrite et guide l'œil du lecteur vers l'information cherchée (contact, adresse, identité).}
\end{popfigure}

\textbf{Fonctions de l'image dans un CV}\par

L'image dans un CV remplit trois fonctions principales :

\begin{itemize}
\item \textbf{Identifier} : la photographie d'identité et les logos permettent au recruteur d'associer un visage au candidat et de reconnaître rapidement les structures mentionnées (université, entreprise connue).
\item \textbf{Illustrer} : un logo ou une photo rend le document concret et crédible ; il atteste visuellement d'une appartenance ou d'une expérience.
\item \textbf{Rendre lisible} : les pictogrammes structurent l'état civil et les coordonnées ; l'œil repère instantanément le numéro de téléphone ou l'adresse sans lire tout le bloc.
\end{itemize}

\begin{attention}[Limites et risques de l'image]
L'image peut aussi desservir le candidat. Une \cle{photo inappropriée} (photo de fête, tenue négligée, arrière-plan confus, selfie) donne une mauvaise première impression et peut éliminer une candidature. De même, la \cle{surcharge} d'images — trop de logos, de couleurs, de pictogrammes décoratifs — nuit à la lisibilité et détourne l'attention du contenu. Règle d'or : une photo d'identité nette, récente, sur fond neutre, et seulement les pictogrammes utiles aux coordonnées.
\end{attention}

\courssub{Analyser une offre d'embauche}

Avant de toucher à son CV, il faut lire l'offre d'embauche comme on analyse un texte : méthodiquement. Toute annonce contient quatre types d'informations à repérer :

\begin{enumerate}
\item le \cle{diplôme} exigé (BT, BTS, DUT, licence professionnelle...) ;
\item l'\cle{expérience} demandée (durée, domaine, type de poste) ;
\item les \cle{compétences} recherchées (techniques : conduire une maintenance ; comportementales : rigueur, travail en équipe) ;
\item les \cle{conditions} : lieu de travail, type de contrat, parfois le salaire.
\end{enumerate}

\begin{exemplebox}[Annonce à analyser]
« \emph{Entreprise industrielle basée à Douala recrute un technicien en froid. Profil : BT ou BTS en froid et climatisation ; expérience d'au moins 2 ans en maintenance industrielle ; maîtrise des installations frigorifiques et des normes de sécurité ; sens de l'organisation. Poste à pourvoir immédiatement.} »

Critères dégagés : diplôme = BT/BTS froid ; expérience = 2 ans minimum en maintenance industrielle ; compétences = installations frigorifiques, normes de sécurité, organisation ; lieu = Douala.
\end{exemplebox}

\begin{exemplebox}[Pourquoi l'expérience datée est décisive]
Une offre exigeant « 2 ans d'expérience en maintenance » élimine d'office les CV qui ne mentionnent \emph{aucune expérience datée}. Un recruteur qui lit « stage de 3 mois » sans dates, ou rien du tout dans la rubrique « Expérience professionnelle », ne peut pas vérifier le critère : le CV part à la corbeille, même si le candidat est compétent. D'où l'importance de dater précisément chaque expérience (ex. : « juin 2023 -- août 2025 : technicien stagiaire puis titulaire, SODECOTON »).
\end{exemplebox}

\courssub{Le principe d'adaptation du CV à l'offre}

\textbf{Faire correspondre ses rubriques aux critères de l'annonce}\par

Le principe est simple : le CV n'est pas une biographie, c'est une \emph{réponse}. À chaque critère de l'annonce doit correspondre, dans le CV, un élément visible qui y répond. Cela suppose de sélectionner et de reformuler ses informations : mettre en avant l'expérience pertinente (même si ce n'est pas la plus longue), adapter le \cle{titre du CV} au poste visé (« Technicien en maintenance frigorifique industrielle » plutôt que le vague « Technicien »), et reprendre le vocabulaire de l'annonce.

\begin{methode}[Adapter son CV à une offre en trois étapes]
\begin{enumerate}
\item \textbf{Souligner les critères de l'annonce} : diplôme, expérience, compétences, lieu. Les recopier dans une colonne.
\item \textbf{Repérer dans son parcours ce qui y répond} : pour chaque critère, chercher dans ses études, stages et emplois l'élément correspondant, même modeste (un stage compte comme une expérience).
\item \textbf{Réorganiser et reformuler les rubriques} : placer en premier ce qui répond à l'offre, ajuster le titre du CV, employer les mots exacts de l'annonce, dater les expériences.
\end{enumerate}
\end{methode}

\begin{popfigure}
\begin{center}
\begin{tabularx}{\linewidth}{|X|X|}
\hline
\rowcolor{popblueL} \textbf{Critère de l'annonce} & \textbf{Élément du CV à valoriser} \\
\hline
BT ou BTS en froid et climatisation & Rubrique « Formation » : diplôme, établissement, année d'obtention \\
\hline
2 ans d'expérience en maintenance industrielle & Rubrique « Expérience » placée en premier, avec dates précises et nom de l'entreprise \\
\hline
Maîtrise des installations frigorifiques & Rubrique « Compétences » : reprendre le terme exact « installations frigorifiques » \\
\hline
Normes de sécurité & Mention d'une formation ou d'une pratique de la sécurité en atelier \\
\hline
Poste à Douala & Adresse ou mention « disponible à Douala » dans l'état civil \\
\hline
\end{tabularx}
\end{center}
\legende{Tableau de correspondance offre $\leftrightarrow$ CV : chaque critère de l'annonce doit trouver une réponse visible dans une rubrique du CV.}
\end{popfigure}

\begin{attention}[Ce qu'on peut ajuster, ce qu'on ne change jamais]
On peut ajuster : l'\cle{ordre} des rubriques (l'expérience avant la formation si elle est le point fort), la \cle{mise en évidence} (développer l'expérience pertinente, réduire les autres), le \cle{vocabulaire} (reprendre les termes de l'annonce), le titre du CV. On ne change \emph{jamais} : les \cle{faits}. Inventer un diplôme, allonger une durée d'emploi, exagérer une fonction, c'est un faux qui peut être vérifié (appel à l'ancien employeur, demande d'attestation) et qui disqualifie le candidat, parfois définitivement. L'honnêteté des faits est la règle absolue.
\end{attention}

\begin{attention}[Erreur fréquente : le CV « passe-partout »]
Envoyer le même CV à toutes les offres, sans adaptation, est l'erreur la plus répandue. Le recruteur, qui lit des dizaines de candidatures, repère immédiatement un CV « passe-partout » : titre générique, rubriques dans le désordre par rapport à ses critères, aucun mot de son annonce. Un CV adapté montre au contraire que le candidat a compris le poste.
\end{attention}

\begin{savaistu}[Les logiciels de tri des CV]
Certains cabinets de recrutement et grandes entreprises utilisent des logiciels qui filtrent les CV par \cle{mots-clés} avant même qu'un humain ne les lise. Le logiciel compare les mots du CV aux mots de l'annonce. Reprendre le vocabulaire exact de l'offre (« maintenance industrielle », « froid commercial ») augmente donc mécaniquement les chances de passer ce premier filtre. Adapter son CV, ce n'est pas de la flatterie : c'est parler la langue du recruteur... et de ses machines.
\end{savaistu}

\courssub{Étude de cas guidée : un technicien en froid face à deux offres}

Monsieur Essono, titulaire d'un BT en froid et climatisation, a travaillé 2 ans comme technicien de maintenance dans une usine de Douala, puis 1 an comme dépanneur de chambres froides pour des commerces. Il répond à deux offres.

\textbf{Offre A} (maintenance industrielle, Douala) : « 2 ans d'expérience en maintenance industrielle, maîtrise des installations frigorifiques. »

\textbf{Offre B} (froid commercial, Bafoussam) : « technicien itinérant pour la maintenance de chambres froides et meubles frigorifiques commerciaux, sens du contact client. »

\begin{exoresolu}[Adapter le CV d'Essono aux deux offres]
Énoncé : indiquer, pour chaque offre, le titre du CV, la rubrique à placer en premier et l'expérience à développer.
\tcblower
\textbf{Pour l'offre A} : titre = « Technicien en maintenance frigorifique industrielle » ; rubrique « Expérience professionnelle » en premier ; développer les 2 ans en usine (dates, tâches de maintenance, normes de sécurité) ; l'année de dépannage commercial est réduite à une ligne. Le lieu (Douala) est un atout à signaler.

\textbf{Pour l'offre B} : titre = « Technicien en froid commercial » ; l'expérience de dépanneur de chambres froides passe en premier et est détaillée (types d'appareils, relation avec les commerçants = sens du contact client) ; l'expérience industrielle reste mentionnée, car elle prouve la solidité technique, mais en second. La disponibilité à se déplacer (poste itinérant) et la mobilité vers Bafoussam sont précisées.

Dans les deux cas, les faits sont identiques : seuls l'ordre, le titre et le niveau de détail changent. C'est exactement cela, adapter un CV.
\end{exoresolu}

\begin{exercice}[À vous de jouer]
Voici une offre : « Supermarché de Yaoundé recrute un technicien comptable. Profil : BT en comptabilité, 1 an d'expérience en caisse ou en gestion des stocks, maîtrise du tableur, rigueur. »
\begin{enumerate}
\item Soulignez les quatre critères de l'annonce.
\item Rédigez le titre du CV adapté à cette offre.
\item Un candidat a fait 6 mois de stage en caisse et 1 an de vente au marché : indiquez ce qu'il doit mettre en avant et comment le formuler honnêtement.
\end{enumerate}
\end{exercice}

\begin{corrige}[Exercice : adapter un CV à l'offre de technicien comptable]
\begin{enumerate}
\item Critères : diplôme = BT en comptabilité ; expérience = 1 an en caisse ou gestion des stocks ; compétences = maîtrise du tableur, rigueur ; lieu = Yaoundé.
\item Titre possible : « Technicien comptable — caisse et gestion des stocks » (les mots de l'annonce sont repris).
\item Le stage de 6 mois en caisse répond directement au critère « caisse » : il doit être daté et détaillé en premier dans la rubrique « Expérience ». L'année de vente au marché peut être formulée honnêtement comme « gestion d'un point de vente : encaissements, suivi des stocks » — ce qui correspond à la « gestion des stocks » sans mentir. La maîtrise du tableur doit apparaître dans « Compétences », avec un exemple concret si possible.
\end{enumerate}
\end{corrige}

\begin{aretenir}[L'essentiel de la section]
\begin{itemize}
\item Un texte iconique porte une information par l'image : photo d'identité, pictogrammes, logos. L'image identifie, illustre et rend lisible, mais une photo inappropriée ou une surcharge nuit à la candidature.
\item Une offre d'embauche s'analyse en quatre critères : diplôme, expérience, compétences, lieu.
\item Adapter son CV, c'est faire correspondre chaque rubrique aux critères de l'annonce : titre ajusté, expérience pertinente mise en avant et datée, vocabulaire de l'offre repris.
\item On ajuste l'ordre, la mise en évidence et le vocabulaire ; on ne falsifie jamais les faits.
\end{itemize}
\end{aretenir}

\coursec{Rédaction, présentation et évaluation du CV}

Après avoir analysé les textes iconiques et appris à adapter son CV à une offre d'embauche précise, nous abordons maintenant l'étape finale et décisive : la concrétisation. Un CV parfait sur le papier mais mal relu, mal présenté à l'oral ou non conforme aux attentes du recruteur perd toute sa valeur. Cette section synthétise l'ensemble du processus de production, de la collecte des informations jusqu'à l'évaluation finale, en intégrant les outils linguistiques (subjonctif, conditionnel) étudiés précédemment.

\courssub{Les étapes de la rédaction : de la collecte à la version finale}

La rédaction d'un CV ne s'improvise pas ; elle suit un processus itératif rigoureux. Chaque étape conditionne la qualité de la suivante.

\begin{methode}[Processus de production d'un CV]
\begin{enumerate}
    \item \textbf{Collecte et tri des informations :} rassembler tous les diplômes, formations, expériences (stages, jobs d'été, bénévolat), compétences techniques (logiciels, langues, permis), centres d'intérêt pertinents. \emph{Trier} : ne garder que ce qui valorise la candidature visée.
    \item \textbf{Choix du plan (structure externe) :} décider entre CV \emph{chronologique} (par dates), \emph{fonctionnel} (par compétences) ou \emph{mixte}. Pour un jeune diplômé technique, le CV chronologique inversé (du plus récent au plus ancien) est souvent le plus lisible.
    \item \textbf{Rédaction (structure interne) :} remplir les rubriques avec des phrases nominales, des verbes d'action à l'infinitif ou au participe passé, du vocabulaire précis. Appliquer les règles de mise en page (police unique, alignements, espacements).
    \item \textbf{Relecture et correction :} vérifier le fond (cohérence, exhaustivité), la forme (orthographe, typographie, mise en page) et l'adéquation à l'offre.
    \item \textbf{Finalisation et présentation :} conversion en PDF (format figé), nommage du fichier clair (\texttt{Nom\_Prenom\_CV\_Poste.pdf}), préparation de la présentation orale.
\end{enumerate}
\end{methode}

\begin{popfigure}
\begin{tikzpicture}[
    node distance=1.2cm and 1.2cm,
    every node/.style={font=\small, align=center},
    box/.style={draw=popblue, thick, rounded corners, fill=popblueL, minimum width=2.6cm, minimum height=1.2cm},
    arrow/.style={-Stealth, thick, popdark}
]
\node[box, align=center] (collecte){1. Collecte \\ et tri des infos};
\node[box, right=of collecte, align=center] (plan){2. Choix du plan \\ (structure externe)};
\node[box, right=of plan, align=center] (redaction){3. Rédaction \\ (structure interne)};
\node[box, right=of redaction, align=center] (relecture){4. Relecture \\ et correction};
\node[box, right=of relecture, align=center] (final){5. Finalisation \\ et présentation};
\draw[arrow] (collecte) -- (plan);
\draw[arrow] (plan) -- (redaction);
\draw[arrow] (redaction) -- (relecture);
\draw[arrow] (relecture) -- (final);
% Boucle de rétroaction
\draw[arrow, poporange, dashed] (relecture.south) -- ++(0,-0.8) -| node[below, pos=0.25, poporange, font=\tiny] {Si erreurs ou incohérences} (redaction.south);
\end{tikzpicture}
\legende{Organigramme des étapes de production du CV : processus itératif avec boucle de rétroaction sur la relecture.}
\end{popfigure}

\courssub{La relecture : une étape cruciale, trop souvent négligée}

La relecture n'est pas une simple vérification orthographique ; c'est un contrôle qualité en trois temps. Un CV truffé de fautes signale un manque de rigueur, qualité pourtant centrale dans les métiers techniques.

\begin{methode}[Relecture efficace en trois passes]
\begin{enumerate}
    \item \textbf{Passe 1 : le fond (cohérence et exhaustivité).} Vérifier les dates (début/fin), la chronologie inverse, l'absence de trous inexpliqués, la présence des mots-clés de l'offre, la pertinence des compétences listées.
    \item \textbf{Passe 2 : la forme (orthographe, grammaire, typographie).} Accords (participes passés, adjectifs), homophones (\emph{à/a}, \emph{et/est}, \emph{son/sont}), majuscules (diplômes, noms propres), ponctuation. Utiliser un correcteur automatique \emph{puis} relire soi-même.
    \item \textbf{Passe 3 : la fluidité (lecture à voix haute).} L'oreille détecte les lourdeurs, les répétitions, les phrases trop longues. Faire relire par un tiers (enseignant, camarade, professionnel) : un regard neuf repère ce que l'auteur ne voit plus.
\end{enumerate}
\end{methode}

\begin{attention}[Cohérence des dates : le piège fatal]
Négliger la cohérence chronologique éveille immédiatement la méfiance du recruteur.
\begin{itemize}
    \item \textbf{Chevauchements impossibles :} ex. \emph{Stage à Douala (janvier--juin 2023)} et \emph{Formation à Yaoundé (mars--août 2023)}.
    \item \textbf{Trous inexpliqués :} une année blanche sans mention (voyage, maladie, recherche active, service civique) laisse place à l'interprétation négative.
    \item \textbf{Correction :} tout trou supérieur à 3 mois doit être justifié en une ligne (ex. \emph{2021--2022 : année de césure -- projet humanitaire dans le Nord, permis B obtenu}).
\end{itemize}
\end{attention}

\courssub{La présentation orale du CV : convaincre en deux minutes}

Lors d'un entretien ou d'un oral de Probatoire, le candidat doit \emph{présenter} son CV, non le \emph{réciter}. C'est un exercice de synthèse orale structurée, mobilisant le conditionnel de politesse et le subjonctif du souhait.

\begin{methode}[Présenter son CV à l'oral (2 minutes maximum)]
\begin{enumerate}
    \item \textbf{Accroche / Identité (15 s) :} « Bonjour, je m'appelle \textbf{Prénom Nom}, j'ai \textbf{âge} ans, je suis titulaire d'un \textbf{Baccalauréat technique, série [série], option [option]} obtenu en \textbf{année} avec mention \textbf{[mention]}. »
    \item \textbf{Parcours et compétences clés (60 s) -- structure en 3 points :}
    \begin{itemize}
        \item \emph{Formation technique :} projet de fin d'études, modules forts, certifications éventuelles.
        \item \emph{Expérience de terrain :} stages, jobs, missions précises, résultats chiffrés (ex. \emph{« maintenance de 15 postes informatiques, taux de satisfaction de 95\,\% »}).
        \item \emph{Compétences transverses :} langues, outils bureautiques, qualités personnelles (rigueur, travail en équipe, autonomie).
    \end{itemize}
    \item \textbf{Motivation et projet (30 s) :} lier son profil au poste. \emph{« Ce poste de technicien de maintenance correspond à mon projet professionnel car il me permettrait de valoriser mes compétences en [compétence] tout en développant [autre compétence]. Je souhaite vivement intégrer votre structure. »}
    \item \textbf{Conclusion / Ouverture (15 s) :} « Je suis à votre disposition pour tout complément d'information. »
\end{enumerate}
\end{methode}

\begin{aretenir}[Outils linguistiques pour l'oral]
\begin{itemize}
    \item \textbf{Conditionnel de politesse / hypothèse :} \emph{« Je \textbf{souhaiterais} intégrer... », « Cela me \textbf{permettrait} de... », « Je \textbf{serais} ravi de... »}. Il atténue l'affirmation et marque le respect.
    \item \textbf{Subjonctif du souhait / de la nécessité :} \emph{« Il est important que je \textbf{puisse} mettre en œuvre... », « Je veux que mon profil \textbf{corresponde} à vos attentes. »}. Il exprime la volonté et l'adéquation recherchée.
    \item \textbf{Indicatif (faits réels, acquis) :} \emph{« J'\textbf{ai réalisé}..., j'\textbf{ai acquis}..., je \textbf{maîtrise}... »}.
\end{itemize}
\end{aretenir}

\begin{exemplebox}[Simulation de présentation orale (extrait)]
\textbf{Candidat :} « Bonjour, je m'appelle \textbf{Armand Ngono}, j'ai 19 ans et je suis titulaire d'un Baccalauréat technique F4 (Génie électrique), mention Assez Bien, session 2024. \\
\textbf{1. Formation :} mon projet de fin d'études portait sur l'automatisme d'une chaîne de conditionnement (automate programmable, capteurs). J'ai validé le module \emph{Électrotechnique industrielle} avec 16/20. \\
\textbf{2. Expérience :} stage de 3 mois à ENEO Douala-Bassa : maintenance préventive sur postes HTA/BT, rédaction de 12 fiches d'intervention. J'ai \textbf{acquis} la rigueur des consignes de sécurité. \\
\textbf{3. Motivation :} \emph{je \textbf{souhaiterais} (conditionnel) rejoindre votre service maintenance pour \textbf{valoriser} mes bases en automatisme. \textbf{Il est essentiel que je \textbf{progresse}} (subjonctif) sur le terrain auprès de vos techniciens confirmés. »} \\
\textbf{Posture :} debout, regard franc vers le jury, mains libres, débit posé, pas de lecture de notes.
\end{exemplebox}

\courssub{Grille d'évaluation du CV : critères et pondération}

L'évaluation (Probatoire, concours, recrutement) repose sur une grille critériée. La connaître permet de s'auto-évaluer avant de rendre sa copie.

\begin{popfigure}
\begin{center}
\begin{tabularx}{\linewidth}{|l|X|c|}
\hline
\rowcolor{popblueL} \textbf{Critère} & \textbf{Indicateurs de réussite} & \textbf{Points} \\
\hline
\textbf{1. Rubriques complètes et pertinentes} & Identité, formation, expériences, compétences techniques et transverses, centres d'intérêt utiles. Aucune rubrique vide. & 6 \\
\hline
\textbf{2. Mise en page et lisibilité} & Police unique, alignements parfaits, espacements réguliers, une page, PDF propre, photo professionnelle (si requise). & 4 \\
\hline
\textbf{3. Qualité de la langue} & Aucune faute d'orthographe ni de grammaire, syntaxe soignée (infinitifs, participes), vocabulaire technique précis, pas de jargon abusif. & 4 \\
\hline
\textbf{4. Adaptation à l'offre / au poste} & Mots-clés de l'offre présents, compétences mises en avant cohérentes, objectif professionnel ciblé. & 3 \\
\hline
\textbf{5. Présentation orale (si oral)} & Structure en 3 points, moins de 2 minutes, posture, regard, voix, emploi du conditionnel et du subjonctif, réponses argumentées. & 3 \\
\hline
\rowcolor{popblueL} \textbf{TOTAL} & & \textbf{20} \\
\hline
\end{tabularx}
\end{center}
\legende{Grille d'évaluation pondérée sur 20 points. Les critères 1 à 4 concernent le document écrit ; le critère 5 concerne l'oral.}
\end{popfigure}

\begin{exoresolu}[Application de la grille : cas d'un élève de Terminale F3]
\textbf{Profil :} élève de Terminale F3 (Construction mécanique), stage dans une entreprise métallurgique de Douala, projet de fin d'études : « Fabrication d'une presse hydraulique ». \textbf{CV soumis :} 1 page, photo, police Calibri 11. \\
\textbf{Évaluez ce CV à l'aide de la grille et proposez une remédiation.}

\tcblower

\textbf{Évaluation détaillée :}
\begin{itemize}
    \item \textbf{Critère 1 (6 pts) :} rubriques toutes remplies. Expérience détaillée (tâches, machines : plieuse, cisaille). Compétences : lecture de plans, soudure TIG/MIG, CAO. $\rightarrow$ \textbf{5,5/6} (manque : permis B non mentionné, pourtant utile sur les chantiers).
    \item \textbf{Critère 2 (4 pts) :} mise en page impeccable, alignements parfaits, en-têtes en gras 14 pt, corps 11 pt. $\rightarrow$ \textbf{4/4}.
    \item \textbf{Critère 3 (4 pts) :} 2 fautes (\emph{« maintenance préventives »} ; \emph{« les stages que j'ai fait »} -- accord du participe passé manquant). Vocabulaire technique correct. $\rightarrow$ \textbf{2,5/4}.
    \item \textbf{Critère 4 (3 pts) :} CV générique, non adapté à l'offre « Technicien de maintenance industrielle » (mots-clés \emph{GMAO}, \emph{diagnostic de panne} absents). $\rightarrow$ \textbf{1,5/3}.
    \item \textbf{Critère 5 (3 pts) -- oral simulé :} présentation de 2 min 30 s (dépassement), regard fuyant, hésitations sur la motivation. Emploi de l'indicatif au lieu du conditionnel (\emph{« Je veux travailler ici »}). $\rightarrow$ \textbf{1/3}.
\end{itemize}
\textbf{Total : 14,5/20 (Assez Bien).} \\
\textbf{Remédiation prioritaire :} relecture orthographique (accords du participe passé), intégration des mots-clés de l'offre, entraînement oral chronométré.
\end{exoresolu}

\courssub{Compte rendu collectif et remédiation}

L'analyse des productions de la classe permet d'identifier des tendances lourdes et d'y remédier collectivement avant l'examen.

\begin{savaistu}[Le dossier de candidature au Cameroun]
Lors des concours administratifs camerounais (ENAM, ENSET, concours professionnels, Fonction publique), le \textbf{dossier de candidature incomplet ou mal présenté est la première cause de rejet avant même l'épreuve écrite}. Une pièce manquante (acte de naissance non légalisé, diplôme non certifié conforme, photo non conforme, CV sur deux pages au lieu d'une), une rature ou une enveloppe mal libellée entraînent l'irrecevabilité immédiate. La rigueur du CV et du dossier est donc une compétence citoyenne et professionnelle vitale.
\end{savaistu}

\begin{aretenir}[Erreurs récurrentes observées en classe (diagnostic type)]
\begin{itemize}
    \item \textbf{Structure :} rubrique « Compétences » vide ou générique (\emph{« Informatique : Word, Excel »} sans niveau), absence de rubrique « Objectif professionnel ».
    \item \textbf{Langue :} accords du participe passé (ex. \emph{« les stages que j'ai fait »} au lieu de \emph{« les stages que j'ai faits »}), confusion \emph{a/à}, \emph{et/est}, majuscules aléatoires sur les noms de diplômes.
    \item \textbf{Modes verbaux :} utilisation exclusive de l'indicatif à l'oral (\emph{« Je veux... »} au lieu de \emph{« Je souhaiterais... »}), absence de subjonctif pour exprimer la nécessité (\emph{« Il faut que je suis... »} au lieu de \emph{« Il faut que je sois... »}).
    \item \textbf{Adaptation :} CV « passe-partout » envoyé à toutes les offres sans modification.
    \item \textbf{Oral :} lecture du CV, dépassement de temps, regard baissé, absence de lien entre la formation et le poste.
\end{itemize}
\end{aretenir}

\begin{methode}[Remédiation ciblée : 3 ateliers de réécriture]
\begin{enumerate}
    \item \textbf{Atelier « Accords et modes » (30 min) :} à partir de CV anonymes de la classe, repérer et corriger 5 fautes d'accords et transformer 3 phrases à l'indicatif en conditionnel ou subjonctif de politesse ou de souhait. \emph{Exercice :} « Je veux intégrer votre équipe » $\rightarrow$ « \textbf{Je souhaiterais} intégrer votre équipe » ; « Il faut que je suis opérationnel » $\rightarrow$ « Il faut que je \textbf{sois} opérationnel ».
    \item \textbf{Atelier « Mots-clés et adaptation » (30 min) :} donner une offre d'emploi réelle. En binômes, surligner 10 mots-clés techniques et comportementaux dans l'offre. Réécrire les rubriques « Compétences » et « Objectif » d'un CV type pour y intégrer ces mots-clés.
    \item \textbf{Atelier « Oral express » (1 h) :} simulation de \emph{speed-recrutement}. Chaque élève présente son CV en 2 minutes chronométrées face à un camarade (jury). Grille d'évaluation simplifiée (structure / langue / posture / motivation). Rotation. Retour collectif sur les 3 points forts et les 3 points faibles globaux.
\end{enumerate}
\end{methode}

\begin{exercice}[Entraînement final : mon CV prêt pour le Probatoire ou un concours]
\begin{enumerate}
    \item Sur la base de votre parcours réel (ou fictif mais réaliste pour un élève de Terminale technique), rédigez votre CV complet sur \textbf{une page A4} (format PDF).
    \item Cochez la check-list avant remise :
    \begin{itemize}
        \item[$\square$] Rubriques obligatoires présentes + 1 rubrique « Atouts / Certifications » ?
        \item[$\square$] Chronologie inverse respectée, dates cohérentes (pas de trou supérieur à 3 mois non justifié) ?
        \item[$\square$] Zéro faute (relecture à voix haute + correcteur + camarade) ?
        \item[$\square$] Mise en page : police unique, alignements, 1 page, fichier \texttt{Nom\_Prenom\_CV.pdf} ?
        \item[$\square$] Adaptation : 5 mots-clés d'une offre cible intégrés ?
    \end{itemize}
    \item Préparez votre présentation orale (2 minutes maximum) en appliquant la méthode « Présenter son CV à l'oral ». Entraînez-vous devant un miroir ou en enregistrant votre voix.
    \item Auto-évaluez-vous sur la \textbf{grille sur 20 points} (critères 1 à 4). Notez votre score cible : \_\_\_/17 (partie écrite).
\end{enumerate}
\end{exercice}

\begin{corrige}[Exercice -- Auto-évaluation type]
La correction est individuelle. L'enseignant collecte les CV (version papier ou numérique), les évalue sur la grille de 20 points, et organise une séance de remise individuelle (5 minutes par élève) : \emph{« Ton point fort : la rubrique Expériences, chiffrée et précise. Ton point faible : 3 fautes d'accords et une adaptation insuffisante à l'offre. Objectif : 16/20 au Probatoire. Travail : fiche de révision sur l'accord du participe passé + réécriture de la rubrique Objectif pour l'offre visée. »}
\end{corrige}

\courssub{Synthèse de l'unité : le CV, miroir de votre professionnalisme}

Rédiger un CV en Terminale technique, c'est mobiliser l'ensemble des compétences du programme :
\begin{itemize}
    \item \textbf{Réception :} analyser une offre d'embauche (texte fonctionnel), décoder des icônes et pictogrammes (textes iconiques).
    \item \textbf{Langue :} maîtriser la morphosyntaxe (accords, valeurs du subjonctif et du conditionnel) pour une communication précise et polie.
    \item \textbf{Production écrite :} structurer l'information (structure externe et interne), synthétiser son parcours, adapter le lexique technique.
    \item \textbf{Production orale :} présenter, argumenter, interagir avec un jury.
\end{itemize}

Le CV n'est pas un document statique ; c'est un \cle{outil vivant} qui évolue avec chaque stage, chaque compétence acquise, chaque candidature. La rigueur technique que vous appliquez en atelier (précision, norme, sécurité), appliquez-la à votre CV : c'est votre première \cle{réalisation professionnelle} évaluée par le monde de l'entreprise.

\newpage

\coursec{Activité d'intégration}

\coursec{Leçon 27 : Rédiger un Curriculum Vitae}

\courssub{Objectifs de l'activité}
\begin{itemize}
    \item Mobiliser l'ensemble des savoirs de la leçon (structure externe et interne du CV, choix lexicaux, modes verbaux) pour produire un document professionnel complet.
    \item Analyser une offre d'emploi réelle afin d'identifier les compétences-clés attendues par le recruteur et d'y adapter son profil.
    \item Rédiger un CV cohérent, sans fautes, respectant l'ordre antichronologique et utilisant des verbes d'action précis.
    \item Rédiger un paragraphe de présentation orale intégrant correctement le subjonctif (souhait) et le conditionnel (politesse/hypothèse).
    \item Développer des compétences transversales : auto-évaluation, évaluation par les pairs, prise de parole en public, gestion du stress face à un jury.
\end{itemize}

\courssub{Situation}
\textbf{Contexte :} Vous êtes élève en Terminale Technique (spécialités : Maintenance Industrielle, Comptabilité/Gestion, ou Productions Animales/Végétales) au Lycée Technique de Bafoussam. La fin de l'année scolaire approche et l'insertion professionnelle devient une priorité. Votre établissement organise une \emph{« Journée de l'Insertion Professionnelle »} en partenariat avec trois structures économiques de la région. Un jury simulé (professeur + deux élèves recruteurs) sélectionnera les trois meilleurs dossiers pour un stage rémunéré de deux mois pendant les grandes vacances.

\textbf{Les trois offres d'embauche (adaptées pour l'exercice) :}

\begin{center}
\begin{tabularx}{\linewidth}{|>{\bfseries}l|X|X|X|}\hline
\textbf{Critères} & \textbf{Offre A : SABC Douala} & \textbf{Offre B : PME « Horizon Services » Yaoundé} & \textbf{Offre C : Coopérative « Terroir Vert » Bafang} \\ \hline
\textbf{Poste} & Technicien de Maintenance Industrielle (H/F) & Secrétaire Comptable (H/F) & Technicien Agricole / Encadreur de proximité (H/F) \\ \hline
\textbf{Secteur} & Agro-industrie (Brasserie) & Services / BTP & Agriculture / Coopérative \\ \hline
\textbf{Diplôme exigé} & Baccalauréat Technique F2/F3 (Maintenance) ou BTS Maintenance & Baccalauréat Technique G2/G3 (Comptabilité) ou BTS Comptabilité & Baccalauréat Technique F4/F5 (Agro) ou BTS Productions Animales/Végétales \\ \hline
\textbf{Expérience} & Stage validé en milieu industriel (6 mois min.) apprécié. Maintenance préventive/curative. & Maîtrise logiciels comptables (Sage, QuickBooks), classeurs Excel avancés. Expérience caisse/banque. & Connaissance cultures de rente (cacao, café, plantain) et élevage porcin/avicole. Permis B (moto) indispensable. \\ \hline
\textbf{Compétences techniques} & Lecture plans PID, soudure TIG/MIG, automatismes (API Siemens/Schneider), GMAO (Logiciel DIMO Maint). & Établissement états financiers, déclaration TVA/IRPP (télédéclaration), suivi trésorerie, archivage numérique. & Techniques compostage, itinéraires techniques, vulgarisation, gestion stocks intrants, rapport d'activité hebdo. \\ \hline
\textbf{Savoir-être} & Rigueur sécurité (HSE), réactivité panne, travail en 3x8, esprit d'équipe. & Discrétion, organisation, sens du contact client/fournisseurs, résistance pression clôture mensuelle. & Pédagogie auprès producteurs, sens terrain, intégrité gestion intrants, adaptation zones rurales. \\ \hline
\textbf{Contact} & DRH SABC, BP 4088 Douala - \texttt{recrutement@sabc.cm} & DG Horizon Services, Bastos Yaoundé - \texttt{rh@horizon-services.cm} & Président Coop. Terroir Vert, Bafang - \texttt{contact@terroirvert.coop} \\ \hline
\end{tabularx}
\end{center}

\textbf{Votre mission :} Choisir \textbf{une seule} de ces offres correspondant à votre filière. Analyser les exigences. Rédiger votre CV « sur mesure » et préparer une présentation orale de 3 minutes maximum pour convaincre le jury.

\courssub{Consignes}
\begin{enumerate}
    \item \textbf{Analyse de l'offre (15 min) :} Surlignez dans l'offre choisie : (a) 3 mots-clés techniques, (b) 2 savoir-être prioritaires, (c) le diplôme exact exigé. Notez-les dans un tableau de correspondance « Exigence $\leftrightarrow$ Mon profil ».
    \item \textbf{Rédaction du CV (45 min) :} Produisez un CV complet sur une page A4 maximum.
        \begin{itemize}
            \item \textit{Structure externe :} En-tête (identité, photo pro, contacts, lien LinkedIn/portfolio si pertinent), titres de rubriques clairs, mise en page aérée, police lisible (Arial/Calibri 11), fichier PDF final.
            \item \textit{Structure interne (ordre antichronologique) :} Formation (diplômes, options, mentions, projets de fin d'études), Expériences (stages, jobs d'été, bénévolat technique) avec \textbf{verbes d'action} au passé composé ou infinitif (ex: \emph{Réalisé, Programmé, Géré, Encadré}), Compétences Techniques (logiciels, machines, normes), Compétences Transverselles (langues, permis, soft skills), Centres d'intérêt (en lien avec le poste si possible).
        \end{itemize}
    \item \textbf{Rédaction du paragraphe de présentation orale (15 min) :} Rédigez un script de 150 mots maximum pour l'oral. \textbf{Contrainte linguistique :} Intégrez obligatoirement :
        \begin{itemize}
            \item Une phrase au \textbf{subjonctif présent} exprimant un \textbf{souhait} ou une \textbf{volonté} (ex: \emph{« Je souhaite que mon intégration se fasse... »}, \emph{« Il faut que je mette mes compétences au service... »}).
            \item Une phrase au \textbf{conditionnel présent} exprimant la \textbf{politesse} ou l'\textbf{hypothèse} (ex: \emph{« Je souhaiterais vous présenter... »}, \emph{« Cette expérience me permettrait de... »}).
        \end{itemize}
        Surlignez ces deux formes verbales dans votre brouillon.
    \item \textbf{Évaluation par les pairs (20 min) :} En groupe de 4, échangez vos CV et grilles d'évaluation (fourni par le professeur). Notez sur 20 (5 pts Structure/Présentation, 5 pts Adaptation à l'offre, 5 pts Langue/Verbes d'action, 5 pts Cohérence/Argumentation). Annotez les forces/faiblesses.
    \item \textbf{Simulation orale (30 min) :} Les 3 meilleurs CV par groupe (selon moyenne pairs) passent devant le jury (Professeur + 2 élèves recruteurs). Présentation 3 min + Questions 2 min. Jury note sur 20 (Oral : aisance, vocabulaire technique, justifications, respect modes verbaux).
    \item \textbf{Remédiation (15 min) :} Retour en classe. Le professeur commente les 3 erreurs majeures récurrentes observées. Chaque élève réécrit \textbf{une rubrique faible} de son CV (souvent « Expériences » ou « Compétences ») en appliquant les corrections.
\end{enumerate}

\courssub{Ressources à mobiliser}
\begin{definition}[Structure externe type du CV Technique]
\textbf{En-tête :} Nom, Prénom, Adresse (Ville, Quartier), Téléphone (WhatsApp pro), Email (prénom.nom@domaine.pro), Photo (format identité, fond neutre, tenue soignée), \textbf{Titre ciblé} (ex: \emph{Technicien Maintenance Industrielle - Spécialité Automatismes}).
\textbf{Rubriques ordonnées :} 1. Formation (Diplômes en préparation/obtenus), 2. Expériences Professionnelles (Stages, Jobs), 3. Compétences Techniques (Dur), 4. Compétences Transverselles (Souple), 5. Langues / Informatique / Permis, 6. Centres d'intérêt (ciblés).
\end{definition}

\begin{propriete}[Structure interne : Règle de l'antichronologie]
Dans chaque rubrique, l'expérience ou le diplôme \textbf{le plus récent} se place \textbf{en premier}. Format date : \emph{09/2023 -- 06/2024} ou \emph{2023--2024}. Description : \textbf{Verbe d'action} + \textbf{Objet} + \textbf{Contexte/Résultat chiffré si possible}.
\end{propriete}

\begin{methode}[Verbes d'action par filière (échantillon)]
\begin{itemize}
    \item \textbf{Maintenance (F2/F3) :} Diagnostiqué, Réparé, Paramétré (API), Soudé, Prévu (GMAO), Lu (Plans PID), Respecté (HSE).
    \item \textbf{Comptabilité (G2/G3) :} Saisi, Rapproché (Banque), Établi (Bilan, TVA), Déclaré (Télédéclaration), Archivé, Analysé (Ratios).
    \item \textbf{Agricole (F4/F5) :} Conduit (Itinéraire technique), Vacciné, Épandu, Conseillé (Vulgarisation), Géré (Stocks intrants), Rédigé (Rapports).
\end{itemize}
\end{methode}

\begin{aretenir}[Valeurs du Subjonctif et du Conditionnel (Rappel)]
\begin{itemize}
    \item \textbf{Subjonctif présent (Souhait / Volonté / Nécessité / Doute) :} Après \emph{il faut que, je veux que, je souhaite que, il est important que, bien que, pour que}. \textbf{Formation :} Radical du 3ème pers. plur. indic. présent + \emph{-e, -es, -e, -ions, -iez, -ent}. Ex: \emph{Je souhaite que tu \textbf{sois} (être) opérationnel rapidement.}
    \item \textbf{Conditionnel présent (Politesse / Hypothèse / Atténuation) :} Radical du futur simple + terminaisons imparfait (\emph{-ais, -ais, -ait, -ions, -iez, -aient}). Ex: \emph{Je \textbf{souhaiterais} (souhaiter) vous rencontrer. / Cette formation me \textbf{permettrait} (permettre) d'évoluer.}
\end{itemize}
\end{aretenir}

\begin{attention}[Pièges fréquents à éviter]
\begin{itemize}
    \item \textbf{CV « passe-partout » :} Ne pas adapter les intitulés de stages/compétences aux mots-clés de l'offre (ex: écrire « Stage en usine » au lieu de « Stage maintenance préventive ligne d'embouteillage » pour la SABC).
    \item \textbf{Verbes faibles :} « J'ai fait », « J'ai aidé », « J'ai regardé ». Remplacer par : « J'ai \textbf{réalisé} », « J'ai \textbf{assisté} le chef d'équipe », « J'ai \textbf{surveillé} ».
    \item \textbf{Fautes d'accord subjonctif :} *Je souhaite que je \textbf{vais} (indicatif) $\rightarrow$ Je souhaite que j'\textbf{aille}. *Il faut que tu \textbf{est} $\rightarrow$ Il faut que tu \textbf{sois}.
    \item \textbf{Conditionnel vs Futur :} *Je \textbf{voudrai} (futur) venir $\rightarrow$ Je \textbf{voudrais} (conditionnel politesse) venir.
\end{itemize}
\end{attention}

\courssub{Démarche et corrigé commenté}
\begin{exoresolu}[Cas pratique : Candidature Offre A -- SABC Douala (Profil F3 Maintenance)]
\textbf{Énoncé :} Rédaction du CV et préparation orale pour l'Offre A (SABC) à partir du profil d'un élève de Terminale F3.
\tcblower
\textbf{Étape 1 : Tableau de correspondance (Analyse offre $\leftrightarrow$ Profil élève)}

\begin{center}
\begin{tabularx}{\linewidth}{|>{\bfseries}X|X|X|}\hline
\textbf{Exigence Offre A (SABC)} & \textbf{Mots-clés surlignés} & \textbf{Mon profil réel (Élève F3 Lycée Tech Bafoussam)} \\ \hline
Diplôme BAC F2/F3 Maintenance & \cle{BAC F3 Maintenance} & \cle{BAC F3 Maintenance en préparation (Juin 2025)} -- Mention Assez Bien prévue. \\ \hline
Stage 6 mois milieu industriel & \cle{Stage industriel}, \cle{Maintenance préventive/curative} & \cle{Stage 8 semaines (Juil-Août 2024) : BRACAM Douala (Brasserie). Maintenance préventive compresseurs, curatif convoyeurs.} \\ \hline
Lecture plans PID, Soudure, API, GMAO & \cle{Plans PID}, \cle{Soudure TIG/MIG}, \cle{API Siemens}, \cle{GMAO DIMO Maint} & \cle{TP Lycée : Lecture PID, Initiation API Siemens S7-1200 (LOGO!), Soudure à l'arc (notion TIG). Pas encore GMAO DIMO.} \\ \hline
Sécurité HSE, Travail 3x8, Équipe & \cle{HSE}, \cle{3x8}, \cle{Équipe} & \cle{Module HSE validé (Attestation). Habitué travaux pratiques en binôme. Disponible pour horaires décalés.} \\ \hline
\end{tabularx}
\end{center}

\textbf{Étape 2 : Rédaction du CV (Extrait commenté)}

\begin{center}
\begin{tabularx}{\linewidth}{|X|}\hline
\textbf{NGUEFACK TEDONGMO Junior} \\
\textit{Technicien Maintenance Industrielle -- Spécialité Automatismes \& Instrumentation} \\
Douala, Bonabéri -- +237 6XX XX XX XX -- junior.nguefack@email.pro -- \cle{LinkedIn: linkedin.com/in/junior-nguefack} \\
\hline
\textbf{\textcolor{popblue}{FORMATION}} \\
\textbf{09/2022 -- 06/2025} \textbf{Baccalauréat Technique Série F3 (Maintenance des Systèmes Automatisés)} -- Lycée Technique de Bafoussam. \textit{Mention Assez Bien attendue. Projet de fin d'études : « Automatisation d'une chaîne de convoyage par API Siemens S7-1200 ».} \\
\textbf{09/2019 -- 06/2022} \textbf{Brevet de Technicien (BT) Maintenance Mécanique} -- CETIC de Bafoussam. \textit{Mention Bien.} \\
\hline
\textbf{\textcolor{popblue}{EXPÉRIENCES PROFESSIONNELLES (Stages)}} \\
\textbf{07/2024 -- 08/2024 (8 sem.)} \textbf{Stagiaire Technicien Maintenance -- BRACAM Douala (Groupe Castel).} \\
\emph{Missions :} \textbf{Réalisé} la maintenance préventive niveau 1 sur compresseurs d'air (vidange, filtres, contrôle pressions) selon gammes GMAO. \textbf{Diagnostiqué} et \textbf{réparé} 3 pannes courantes sur convoyeurs à bandes (changement roulements, tension courroies). \textbf{Lus} plans PID pour repérage vannes/Actionneurs. \textbf{Respecté} strictement consignes HSE (Port EPI, Consignation/LOTO). \textbf{Participé} à l'équipe de quart (simulation 3x8). \\
\textbf{07/2023 -- 08/2023 (6 sem.)} \textbf{Stagiaire Atelier -- Garage Moderne Bafoussam.} \\
\emph{Missions :} \textbf{Assisté} le chef d'atelier pour révision moteurs thermiques. \textbf{Effectué} soudures à l'arc sur châssis. \textbf{Géré} l'outillage (inventaire, rangement). \\
\hline
\textbf{\textcolor{popblue}{COMPÉTENCES TECHNIQUES (DUR)}} \\
\textbf{Automatismes :} Programmation de base API Siemens S7-1200 (TIA Portal V16), LOGO! Soft Comfort. Lecture plans électriques (NF EN 60617) \& PID (ISA S5.1). \\
\textbf{Mécanique / Soudure :} Maintenance préventive/curative pneumatique/hydraulique/mécanique. Soudure à l'arc (électrode enrobée) -- Notions TIG/MIG (TP Lycée). Lecture de plans de définition/ensemble. \\
\textbf{Outils/GMAO :} Notions GMAO (suivi OT, stocks pièces rechange). Pack Office (Excel avancé : TCD, RechercheV pour suivi interventions). \\
\hline
\textbf{\textcolor{popblue}{COMPÉTENCES TRANSVERSELLES (SOFT)}} \\
Rigueur \& Respect normes HSE (Attestation Lycée). Adaptabilité horaires (3x8). Esprit d'équipe (Projets collectifs, Sport : Handball). Communication technique orale/écrite (Rapports de stage, Soutenance PFE). \\
\hline
\textbf{\textcolor{popblue}{LANGUES / INFORMATIQUE / PERMIS}} \\
Français (Langue maternelle), Anglais (Niveau B1 -- Lu technique, Oral courant), Dschang (Maternel). \\
Informatique : Windows 10/11, Office 365, TIA Portal (Base), AutoCAD 2D (Notions). \\
Permis : B (Voiture) -- En cours (Obtention Juillet 2025). \\
\hline
\textbf{\textcolor{popblue}{CENTRES D'INTÉRÊT}} \\
Handball (Club Lycée -- Poste Arrière : Réactivité, Esprit équipe). Bricolage électronique (Réparation petits appareils, Arduino). Veille technologique (Revue \emph{Maintenance \& Entreprise}, Chaîne YouTube \emph{Automatisme Pro}). \\
\hline
\end{tabularx}
\end{center}

\textbf{Commentaires du correcteur (Professeur) :}
\begin{itemize}
    \item \textbf{Titre ciblé :} Reprend l'intitulé exact de l'offre + spécialité (Automatismes) $\rightarrow$ \textbf{Adaptation réussie}.
    \item \textbf{Expérience BRACAM :} Mise en avant (1er). Verbes d'action variés (\textbf{Réalisé, Diagnostiqué, Lus, Respecté}). Chiffres implicites (3 pannes, 8 sem.). Lien direct avec SABC (même secteur brasserie, même process compresseurs/convoyeurs). Mention HSE et 3x8 $\rightarrow$ \textbf{Savoir-être validés}.
    \item \textbf{Compétences Techniques :} Honnêteté sur GMAO (Notions) et TIG/MIG (Notions TP). Valorisation API Siemens (exigence offre). Excel avancé (transférable GMAO).
    \item \textbf{Permis B :} Noté « en cours » avec date $\rightarrow$ Transparence.
\end{itemize}

\textbf{Étape 3 : Paragraphe de présentation orale (Script 3 min -- $\approx$ 160 mots)}

\begin{quote}
\textit{« Monsieur le Président, Mesdames, Messieurs les membres du jury, bonjour. Je m'appelle Junior Nguefack, futur titulaire d'un Baccalauréat F3 Maintenance des Systèmes Automatisés.}

\textit{Mon choix pour la SABC n'est pas un hasard. \textbf{Je souhaite que mon intégration se fasse au sein d'une entreprise de référence} où la rigueur technique et la sécurité sont des valeurs quotidiennes. \textbf{Je souhaiterais} vous présenter brièvement mon parcours.}

\textit{Lors de mon stage de huit semaines à la BRACAM de Douala, j'ai été confronté à la réalité de la maintenance en milieu brassicole : gestion de l'air comprimé, suivi des convoyeurs, respect strict des procédures HSE et consignation. J'ai particulièrement apprécié le travail en équipe sur les quarts. Ce stage a confirmé mon goût pour la maintenance préventive, pilier de la performance industrielle.}

\textit{Ma formation m'a doté de solides bases en automatismes (API Siemens) et en lecture de plans PID, compétences que je suis prêt à approfondir sur vos lignes de production. \textbf{Il est important que je mette} rapidement mon énergie et ma capacité d'apprentissage au service de votre outil de production.}

\textit{Je vous remercie de votre attention et reste à votre disposition pour tout échange technique. »}
\end{quote}

\textbf{Analyse linguistique (Contraintes respectées) :}
\begin{itemize}
    \item \textbf{Subjonctif (Souhait) :} \emph{« Je \textbf{souhaite que mon intégration \textbf{se fasse}} »} (verbe \emph{se faire} -- 3ème pers. sg. subj. présent). \emph{« Il est important que je \textbf{mette} »} (verbe \emph{mettre} -- 1ère pers. sg. subj. présent). \checkmark
    \item \textbf{Conditionnel (Politesse) :} \emph{« Je \textbf{souhaiterais} vous présenter »} (verbe \emph{souhaiter} -- 1ère pers. sg. cond. présent). \checkmark
    \item \textbf{Vocabulaire technique :} Air comprimé, convoyeurs, HSE, consignation, API, PID, maintenance préventive. \checkmark
    \item \textbf{Structure :} Accroche $\rightarrow$ Motivation/Adéquation $\rightarrow$ Preuve (Stage) $\rightarrow$ Compétences transférables $\rightarrow$ Engagement (Subjonctif) $\rightarrow$ Conclusion polie. \checkmark
\end{itemize}

\textbf{Étape 4 : Simulation grille d'évaluation par les pairs (Extrait)}

\begin{center}
\begin{tabularx}{\linewidth}{|l|>{\centering\arraybackslash}X|>{\centering\arraybackslash}X|X|}\hline
\textbf{Critère (5 pts)} & \textbf{Note Pair 1} & \textbf{Note Pair 2} & \textbf{Commentaire synthétique} \\ \hline
Structure / Présentation (Mise en page, photo, clarté) & 5 / 5 & 4,5 / 5 & « Très pro, photo adaptée, rubriques claires. Juste interligne un peu juste rubrique Compétences. » \\ \hline
Adaptation à l'offre (Mots-clés SABC, Brasserie, HSE, 3x8) & 5 / 5 & 5 / 5 & « Parfait. Stage BRACAM mis en 1er. HSE et 3x8 cités explicitement. » \\ \hline
Langue / Verbes d'action (Variété, justesse, modes verbaux oral) & 4,5 / 5 & 4 / 5 & « Verbes variés. Subjonctifs \emph{se fasse/mette} corrects. Conditionnel \emph{souhaiterais} correct. Attention à "j'ai été confronté" (passif) $\rightarrow$ préférer "j'ai géré/résolu". » \\ \hline
Cohérence / Argumentation (Lien formation-emploi, projet pro) & 5 / 5 & 5 / 5 & « Projet pro clair : maintenance préventive/automatismes. Honnête sur lacunes (GMAO, TIG). » \\ \hline
\textbf{TOTAL MOYEN} & \textbf{19,5 / 20} & \textbf{18,5 / 20} & \textbf{Moyenne : 19 / 20 -- SÉLECTIONNÉ POUR JURY} \\ \hline
\end{tabularx}
\end{center}

\textbf{Étape 5 : Remédiation post-jury (Réécriture ciblée)}
\textit{Observation jury :} « Rubrique Compétences Techniques : "Notions GMAO" trop vague. Le recruteur SABC utilise DIMO Maint. L'élève a fait du Excel TCD/RechercheV pour suivi interventions. Il faut valoriser ce transfert. »
\textit{Réécriture corrigée (Élève) :}
\begin{quote}
\textbf{Outils / GMAO :} Utilisation avancée d'Excel (Tableaux Croisés Dynamiques, RechercheX, Mise en forme conditionnelle) pour \textbf{créer un tableau de bord de suivi des interventions} (MTTR, MTBF, Taux de disponibilité) durant le stage BRACAM. \textbf{Transférable immédiatement} sur logiciel GMAO type DIMO Maint (architecture base de données équipement/interventions comprise).
\end{quote}
\end{exoresolu}

\courssub{Bilan de l'activité}
Cette activité d'intégration a permis de valider l'acquisition des compétences terminales de l'unité « Rédiger un CV » dans une situation professionnalisante authentique.

\begin{itemize}
    \item \textbf{Sur le plan méthodologique :} Les élèves ont compris que le CV n'est pas une biographie figée mais un \textbf{document de vente} qui se \textbf{reconfigure} pour chaque offre. L'exercice du « Tableau de correspondance » (Consigne 1) s'est révélé le levier principal pour passer du CV « catalogue » au CV « ciblé ». L'ordre antichronologique et l'usage des verbes d'action sont désormais perçus comme des normes professionnelles et non de simples consignes scolaires.
    \item \textbf{Sur le plan linguistique :} L'obligation d'insérer un subjonctif de souhait et un conditionnel de politesse dans l'oral a forcé la manipulation consciente de ces modes. Les erreurs classiques (\emph{*je souhaite que je vais}, \emph{*je voudrai}) ont été corrigées \textit{in situ} lors de la préparation du script, ancrant la règle morphologique dans un usage communicationnel réel (la politesse du candidat, la projection dans l'entreprise).
    \item \textbf{Sur le plan transverse (Soft Skills) :} L'évaluation par les pairs (grille critériée) a développé l'esprit critique et la bienveillance exigée. La simulation de jury a géré le stress oral, la posture, le regard et la gestion du temps (3 min chrono). Les retours du jury (professeur + pairs recruteurs) ont mis en lumière l'écart entre le \emph{ressenti} de l'élève (« j'ai mis mes stages ») et la \emph{lecture} du recruteur (« je ne vois pas ce que tu as fait concrètement »).
    \item \textbf{Remédiation effective :} La réécriture immédiate d'une rubrique faible (souvent « Compétences » ou « Expériences » trop vagues) ancre la correction. L'élève repart avec un CV \textbf{finalisé, corrigé, validé par les pairs et le jury}, prêt à être déposé pour de vraies candidatures (stages, emplois, BTS).
\end{itemize}

\begin{aretenir}[Compétences certifiées en fin d'activité]
\begin{enumerate}
    \item Analyser une offre d'emploi pour en extraire les critères de sélection (Hard skills / Soft skills).
    \item Produire un CV technique normé (Structure externe/interne, antichronologie, verbes d'action, adaptation lexicale).
    \item Mobiliser le subjonctif (souhait/nécessité) et le conditionnel (politesse/hypothèse) à l'oral comme à l'écrit dans un contexte professionnel.
    \item Présenter son parcours de manière structurée, argumentée et chronométrée face à un jury.
    \item S'auto-évaluer et évaluer un pair à l'aide d'une grille critériée ; réviser sa production suite à un retour critique.
\end{enumerate}
\end{aretenir}

\newpage

\coursec{Méthodes et exercices résolus}

\courssub{Découverte : qu'est-ce qu'un CV ?}

\begin{methode}[Analyser et définir un Curriculum Vitae]
    \textbf{Objectif :} Comprendre la nature, la fonction et les caractéristiques d'un CV pour le distinguer des autres documents professionnels.
    \begin{enumerate}
        \item \textbf{Identifier le genre :} Repérer qu'il s'agit d'un \cle{texte fonctionnel} à visée professionnelle (recherche d'emploi, stage, formation).
        \item \textbf{Définir la finalité :} Le CV est une \cle{publicité de soi} : il doit convaincre un recruteur en quelques secondes (lecture diagonale) de vous convoquer en entretien.
        \item \textbf{Lister les critères de qualité (les 3 C) :} \textbf{C}oncis (1 page idéalement, 2 max), \textbf{C}laire (structure aérée, police lisible), \textbf{C}ohérent (parcours logique, pas de trous inexpliqués).
        \item \textbf{Distinguere CV et Lettre de motivation :} Le CV = \emph{faits, dates, données objectives} ; La Lettre = \emph{projet, motivation, subjectivité, « je »}.
    \end{enumerate}
    \textbf{Réflexe Bac :} Face à un sujet « Qu'est-ce qu'un CV ? », répondre par une définition structurée : Nature + Finalité + Structure type + Différence avec la lettre.
\end{methode}

\begin{exoresolu}[Analyse de la nature d'un CV]
\textbf{Énoncé :} Vous êtes chargé d'animer une séance d'information pour des élèves de Seconde sur le thème « Préparer son entrée dans la vie active ». Rédigez une fiche pédagogique définissant le CV en précisant sa nature, sa finalité, sa structure type et sa différence avec la lettre de motivation.
\tcblower
\textbf{Solution détaillée :}

\textbf{FICHE PÉDAGOGIQUE : LE CURRICULUM VITAE (CV)}

\textbf{1. Définition et nature}
Le Curriculum Vitae (locution latine signifiant « déroulement de la vie ») est un \textbf{texte fonctionnel} de type administratif et professionnel. C'est un document \textbf{normalisé} qui présente de façon synthétique et chronologique le parcours d'un candidat.

\textbf{2. Finalité (But visé)}
\begin{itemize}
    \item Obtenir un \textbf{entretien d'embauche} (et non l'emploi directement).
    \item Servir de \textbf{base de discussion} lors de l'entretien.
    \item Prouver l'\textbf{adéquation} Profil du poste / Profil du candidat.
\end{itemize}

\textbf{3. Structure type (Rubriques obligatoires)}
\begin{enumerate}
    \item \textbf{En-tête :} Identité (Nom, Prénom), Coordonnées (Tél, Email, Localisation), Titre visé / Profil professionnel.
    \item \textbf{Formation / Diplômes :} Du plus récent au plus ancien (antichronologique).
    \item \textbf{Expériences professionnelles :} Stages, jobs d'été, alternance, bénévolat significatif (Entreprise, Dates, Poste, 2-3 missions clés/résultats).
    \item \textbf{Compétences :} Techniques (logiciels, langues, permis), Transverses (management, gestion de projet).
    \item \textbf{Centres d'intérêt / Vie associative :} Sélectifs, valorisant des soft skills (sport collectif, engagement associatif).
\end{enumerate}

\textbf{4. Différence fondamentale CV vs Lettre de motivation}

\begin{center}
\begin{tabularx}{\linewidth}{|l|X|X|}\hline
\textbf{Critère} & \textbf{CV} & \textbf{Lettre de motivation} \\ \hline
\textbf{Nature} & Factuel, objectif, schématique & Argumentatif, subjectif, narratif \\ \hline
\textbf{Contenu} & « Quoi ? Quand ? Où ? » (Données brutes) & « Pourquoi ? Comment ? » (Projet, valeur ajoutée) \\ \hline
\textbf{Style} & Télégraphique, noms, verbes d'action, pas de « je » & Phrases complètes, « je », connecteurs logiques \\ \hline
\textbf{Longueur} & 1 page (standard jeune diplômé) & 1 page maximum \\ \hline
\end{tabularx}
\end{center}

\textbf{Conclusion :} Le CV est la \cle{carte d'identité professionnelle} ; il doit être adapté à chaque candidature (mots-clés de l'offre) et mis à jour régulièrement.
\end{exoresolu}

\courssub{La structure externe du CV}

\begin{methode}[Maîtriser la structure externe (Mise en page et présentation)]
    \textbf{Objectif :} Produire un document visuellement professionnel, lisible en 30 secondes (lecture en « F » ou « Z »).
    \begin{enumerate}
        \item \textbf{Format \& Marges :} Format A4, marges symétriques (2-2,5 cm), interlignes 1,15.
        \item \textbf{Police :} Sans empattement (Sans-serif) pour l'écran/impression rapide : \emph{Arial, Calibri, Helvetica, Roboto}. Taille 10-11 pour le corps, 14-16 pour le Nom, 12 pour les rubriques. \textbf{Une seule police} (gras/italique pour hiérarchiser).
        \item \textbf{Mise en page :} Alignement à gauche (standard), colonnes possibles (colonne étroite à gauche pour dates/icônes, large à droite pour contenu). Espace blanc suffisant (aération).
        \item \textbf{En-tête visuel :} Nom \textsc{PRÉNOM} en grand. Titre visé (ex: \emph{Technicien en Maintenance Industrielle}). Photo : \textbf{Professionnelle} (fond neutre, tenue adaptée, sourire) ou absente selon pays/secteur (au Cameroun : souvent attendue).
        \item \textbf{Enregistrement :} \textbf{PDF uniquement} (nom du fichier : \texttt{Nom\_Prenom\_PosteVisé.pdf}). Jamais Word (.docx) pour envoi final.
    \end{enumerate}
    \textbf{Formule à retenir :} \cle{Sobriété = Professionnalisme}. Pas de couleurs criardes, pas de cadres lourds, pas de cliparts. Une couleur d'accent maximum (bleu marine, anthracite).
\end{methode}

\begin{exoresolu}[Correction de la mise en page d'un CV]
\textbf{Énoncé :} Voici des extraits de commentaires d'un recruteur sur des CV reçus : \\
1. « Police fantaisie, illisible sur mobile. » \\
2. « 3 pages pour un jeune diplômé, trop long. » \\
3. « Photo de vacances / selfie. » \\
4. « Fichier reçu en .pages (Mac) ou .docx, mise en page décalée. » \\
5. « Coordonnées incomplètes (pas de téléphone, email fantaisiste \texttt{killerboy99@...}). » \\
Proposez pour chaque erreur la règle de correction correspondante issue de la méthode.
\tcblower
\textbf{Solution détaillée :}

\begin{center}
\begin{tabularx}{\linewidth}{|l|X|X|}\hline
\textbf{Erreur relevée} & \textbf{Règle violée (Méthode)} & \textbf{Correction à appliquer} \\ \hline
1. Police fantaisie & Étape 2 : Police Sans-serif unique, lisible & Remplacer par \emph{Arial, Calibri, Roboto} taille 11. \\ \hline
2. 3 pages (jeune diplômé) & Étape 1 : Concision (1 page standard) & Supprimer détails superflus (collège, jobs sans lien), réduire interlignes, synthétiser. \\ \hline
3. Photo non pro (selfie) & Étape 4 : Photo professionnelle ou aucune & Photo d'identité récente, fond uni, tenue soignée, cadrage visage/épaules. Ou supprimer. \\ \hline
4. Format .pages / .docx & Étape 5 : Export \textbf{PDF} impératif & Fichier $\to$ Exporter $\to$ PDF. Vérifier l'affichage sur un autre poste. \\ \hline
5. Coordonnées incomplètes / email fantaisiste & Étape 4 : En-tête complet \& pro & Tél + Email pro (\texttt{prenom.nom@domaine.com}), LinkedIn à jour si pertinent. \\ \hline
\end{tabularx}
\end{center}

\textbf{Conclusion :} La forme conditionne la lecture du fond. Un CV mal présenté est souvent écarté \emph{avant même d'être lu}.
\end{exoresolu}

\courssub{La structure interne du CV}

\begin{methode}[Construire la structure interne (Rubriques et rédaction du contenu)]
    \textbf{Objectif :} Remplir chaque rubrique avec des informations pertinentes, hiérarchisées et orientées « résultats ».
    \begin{enumerate}
        \item \textbf{Ordre des rubriques :} \textbf{Jeune diplômé / Étudiant :} Formation $\to$ Expériences (Stages/Jobs) $\to$ Compétences $\to$ Intérêts. \textbf{Expérimenté :} Expériences $\to$ Formation $\to$ Compétences.
        \item \textbf{Règle antichronologique :} Dans chaque rubrique, du plus récent au plus ancien.
        \item \textbf{Rédiger une expérience (Méthode ACTION / STAR) :}
            \begin{itemize}
                \item \textbf{Contexte :} Entreprise, Dates, Poste exact.
                \item \textbf{Actions / Missions :} Verbes d'action à l'infinitif ou participe passé (ex: \emph{Géré, Conçu, Réalisé, Animé}). \textbf{Quantifier} (chiffres, \%, budgets, effectifs).
                \item \textbf{Résultats :} Ce que cela a apporté (gain de temps, CA, satisfaction client).
            \end{itemize}
        \item \textbf{Compétences :} Grouper par catégories (Techniques, Linguistiques, Informatiques, Transverses). Indiquer le \textbf{niveau} (ex: Anglais \emph{B2/C1 - TOEIC 850}, Excel \emph{Avancé : TCD, Macros}).
        \item \textbf{Profil / Accroche (Optionnel mais recommandé) :} 3 lignes sous le nom. \emph{Ex: « Technicien supérieur en Électrotechnique, 2 ans d'expérience en maintenance préventive/curative sur lignes HT/BT. Rigoureux, habilité BR/BC, permis B. Recherche poste de Technicien Maintenance. »}
    \end{enumerate}
    \textbf{Piège :} Ne pas mentir (vérification facile), ne pas surcharger (sélectionner selon l'offre).
\end{methode}

\begin{exoresolu}[Rédaction d'une rubrique Expériences avec la méthode STAR]
\textbf{Énoncé :} Un élève de Terminale STI (Génie Mécanique) a effectué un stage de 8 semaines chez \texttt{CAMRAIL} (atelier de maintenance wagons). Il a : surveillé la maintenance de 15 wagons, participé au démontage/remontage de bogies, rédigé 3 fiches de non-conformité, utilisé la GMAO (logiciel Maximo). Il note bêtement : « Stage CAMRAIL : Maintenance wagons ». Réécrivez cette expérience correctement.
\tcblower
\textbf{Solution détaillée :}

\textbf{AVANT (Erreur classique) :}
\begin{quote}
\textbf{2023 :} Stage CAMRAIL -- Maintenance wagons.
\end{quote}

\textbf{APRÈS (Application méthode STAR / Verbes d'action + Chiffres) :}

\textbf{EXPÉRIENCES PROFESSIONNELLES}
\begin{itemize}
    \item \textbf{Avril -- Juin 2023 \hfill \emph{Stagiaire Technicien Maintenance Matériel Roulant}} \\
    \textbf{CAMRAIL -- Atelier Central de Douala} \hfill \emph{(8 semaines)}
    \begin{itemize}
        \item \textbf{Réalisé} la maintenance préventive niveau 1 sur \textbf{15 wagons} (contrôle freins, essieux, attelages) selon gammes constructeur.
        \item \textbf{Participé} au démontage, contrôle dimensionnel et remontage de \textbf{4 bogies} (respect des tolérances, graissage).
        \item \textbf{Rédigé} \textbf{3 fiches de non-conformité} (détection fissures, usure anormale) transmises au Chef d'équipe $\to$ \textbf{immobilisation préventive} évitant incidents.
        \item \textbf{Utilisé} la \textbf{GMAO Maximo} : consultation ordres de travail, clôture interventions, gestion stocks pièces de rechange.
        \item \textbf{Respecté} strictement les consignes S\&S (Port EPI, procédures verrouillage/étiquetage \emph{Lockout/Tagout}).
    \end{itemize}
\end{itemize}

\textbf{Analyse des gains :}
\begin{itemize}
    \item \textbf{Verbes d'action} variés (Réalisé, Participé, Rédigé, Utilisé, Respecté).
    \item \textbf{Chiffres clés} : 15 wagons, 4 bogies, 3 fiches, 8 semaines.
    \item \textbf{Vocabulaire technique} : GMAO, bogies, non-conformité, Lockout/Tagout, EPI.
    \item \textbf{Résultat implicite} : Sécurité, fiabilité, traçabilité.
\end{itemize}
\textbf{Conclusion :} Cette rédaction transforme une « ligne » en une \cle{preuve de compétence} exploitable par le recruteur.
\end{exoresolu}

\courssub{Langue : les valeurs du subjonctif et du conditionnel}

\begin{methode}[Utiliser le Subjonctif et le Conditionnel dans le cadre professionnel (CV / Lettre / Entretien)]
    \textbf{Objectif :} Maîtriser les valeurs modales de ces deux modes pour exprimer la volonté, l'hypothèse, la politesse, le doute ou la nécessité dans l'écrit professionnel.
    \begin{enumerate}
        \item \textbf{Le Subjonctif (Subjectivité / Nécessité / Volonté) :} Utilisé après verbes d'expression de la volonté, nécessité, doute, sentiment, impersonnels (\emph{il faut que, il est indispensable que, je souhaite que, je veux que, bien que}).
            \begin{itemize}
                \item \textbf{Valeur de nécessité/obligation :} « Il est \textbf{indispensable que le candidat \emph{maîtrise}} l'anglais technique. »
                \item \textbf{Valeur de volonté/souhait :} « Je \textbf{souhaite que mon profil \emph{retienne}} votre attention. »
                \item \textbf{Valeur de concession (Bien que + subj) :} « \textbf{Bien que je \emph{sois}} débutant, je suis opérationnel rapidement. »
            \end{itemize}
        \item \textbf{Le Conditionnel (Hypothèse / Politesse / Atténuation / Projet futur) :}
            \begin{itemize}
                \item \textbf{Politesse (Formules types) :} « \textbf{Je \emph{souhaiterais}} / \textbf{Je \emph{voudrais}} vous présenter ma candidature. » « \textbf{Serait-il possible d'}envisager un entretien ? »
                \item \textbf{Hypothèse / Condition (Si + Imparfait $\to$ Conditionnel Présent) :} « \textbf{Si j'\emph{intégrais}} votre équipe, \textbf{je \emph{m'investirais}} pleinement. »
                \item \textbf{Atténuation / Prudence :} « Cette expérience \textbf{me \emph{permettrait de}} développer mes compétences en gestion de projet. »
                \item \textbf{Futur dans le passé (Récit) :} Rare en CV, possible en lettre : « Il m'a confirmé qu'il \textbf{me \emph{recontacterait}} la semaine prochaine. »
            \end{itemize}
        \item \textbf{Distinction cruciale :} \emph{Je veux que tu viennes} (Subj : volonté sur autrui) vs \emph{Je voudrais venir} (Cond : mon souhait poli). \emph{Il faut que tu sois} (Subj : nécessité) vs \emph{Tu devrais être} (Cond : conseil).
    \end{enumerate}
    \textbf{Réflexe Bac :} Transformer une phrase indicative en subjonctif/conditionnel selon le contexte (politesse, hypothèse, nécessité).
\end{methode}

\begin{exoresolu}[Transformation de modes : Indicatif $\to$ Subjonctif / Conditionnel]
\textbf{Énoncé :} Réécrivez les phrases suivantes en utilisant le mode indiqué entre parenthèses pour les adapter à un contexte de lettre de motivation ou d'entretien.
\begin{enumerate}
    \item « Je \textbf{cherche} un stage. » (Conditionnel de politesse)
    \item « Tu \textbf{as} des compétences en soudure. » (Subjonctif après « Il faut que »)
    \item « Je \textbf{suis} disponible immédiatement. » (Conditionnel d'atténuation/hypothèse)
    \item « Tu \textbf{viendras} à l'entretien. » (Subjonctif après « Je souhaite que »)
    \item « Si j'\textbf{ai} le poste, je \textbf{travaillerai} dur. » (Conditionnel présent / Hypothèse)
\end{enumerate}
\tcblower
\textbf{Solution détaillée :}

\begin{enumerate}
    \item \textbf{Je \emph{souhaiterais} / \emph{voudrais} trouver un stage.} \\
    \textit{Justification :} Politesse, demande atténuée. « Je cherche » est trop direct/brutal.
    \item \textbf{Il faut que tu \emph{aies} des compétences en soudure.} \\
    \textit{Justification :} « Il faut que » + Subjonctif présent (nécessité). Verbe \emph{avoir} : que tu \emph{aies}.
    \item \textbf{Je \emph{serais} disponible immédiatement.} (ou \emph{Je \emph{serais} ravi d'être} disponible...). \\
    \textit{Justification :} Conditionnel présent (atténuation, prudence). « Je suis » = certitude absolue ; « Je serais » = sous réserve de confirmation/entente.
    \item \textbf{Je souhaite que tu \emph{viennes} à l'entretien.} \\
    \textit{Justification :} « Je souhaite que » + Subjonctif présent (volonté sur autrui). Verbe \emph{venir} : que tu \emph{viennes}.
    \item \textbf{Si j'\emph{avais} le poste, je \emph{travaillerais} dur.} \\
    \textit{Justification :} Structure hypothétique \textbf{Si + Imparfait $\to$ Conditionnel Présent}. « Si j'ai » (indicatif) $\to$ « Si j'avais » (imparfait). « Je travaillerai » (futur) $\to$ « Je travaillerais » (conditionnel).
\end{enumerate}

\textbf{Point de vigilance (Bac) :} Ne pas confondre \textbf{Conditionnel} (\emph{-ais, -ais, -ait, -ions, -iez, -aient}) et \textbf{Futur simple} (\emph{-ai, -as, -a, -ons, -ez, -ont}) à la 1ère pers. sg : \emph{Je travaillerai} (certitude) vs \emph{Je travaillerais} (hypothèse/politesse). À l'oral, même prononciation souvent, à l'écrit : distinction obligatoire.
\end{exoresolu}

\courssub{Textes iconiques et adaptation du CV à une offre d'embauche}

\begin{aretenir}[Le texte iconique au service du CV]
    \textbf{Définition :} Tout élément visuel non textuel porteur de sens : photo d'identité, icônes (téléphone, email, localisation, LinkedIn), logos d'entreprises/écoles, puces graphiques, barres de niveau de compétence, code QR vers portfolio/LinkedIn.
    \textbf{Règles d'usage :}
    \begin{itemize}
        \item \textbf{Pertinence :} Chaque icône doit avoir une fonction (repérage rapide), pas de décor gratuit.
        \item \textbf{Cohérence graphique :} Même style (plein/creux), même épaisseur, même couleur d'accent.
        \item \textbf{Lisibilité :} Taille suffisante, contraste élevé. Éviter les logos d'entreprises en fond de page (illisible à l'impression).
        \item \textbf{Photo :} Format carré ou portrait, fond neutre, tenue professionnelle, cadrage visage/épaules, sourire naturel. Au Cameroun, elle est fortement recommandée.
    \end{itemize}
    \textbf{Attention :} Les logiciels de tri automatique (ATS) ne « lisent » pas les images. Les informations critiques (coordonnées, compétences) doivent \textbf{aussi} être en texte clair.
\end{aretenir}

\begin{methode}[Adapter son CV à une offre d'embauche (Analyse d'offre \& Mots-clés)]
    \textbf{Objectif :} Passer les filtres ATS (logiciels de tri automatique) et montrer l'adéquation immédiate.
    \begin{enumerate}
        \item \textbf{Analyser l'offre (Déconstruction) :} Surligner : \textbf{Mots-clés techniques} (logiciels, normes, langues, diplômes), \textbf{Compétences comportementales} (Soft skills : autonomie, esprit d'équipe, rigueur), \textbf{Verbes d'action} attendus.
        \item \textbf{Créer une « Banque de mots-clés » :} Lister les termes exacts de l'offre.
        \item \textbf{Adapter le Titre / Profil :} Reprendre l'intitulé exact du poste visé.
        \item \textbf{Adapter la rubrique Compétences :} Placer en premier les compétences \textbf{exigées} dans l'offre. Utiliser \textbf{le vocabulaire de l'offre} (ex: l'offre dit « GMAO », vous écrivez « GMAO », pas « Logiciel maintenance »).
        \item \textbf{Adapter les Expériences :} Réorganiser les puces : mettre en premier les missions qui répondent aux besoins de l'offre. Supprimer ou réduire les expériences sans lien (sauf trou dans le parcours).
        \item \textbf{Vérification finale :} Le recruteur doit voir la correspondance en \textbf{6 secondes} (règle des 6 secondes).
    \end{enumerate}
    \textbf{Formule :} \cle{CV Master (complet) $\to$ CV Cible (adapté)}. On ne poste JAMAIS le CV Master.
\end{methode}

\begin{exoresolu}[Adaptation d'un CV à une offre réelle]
\textbf{Énoncé :} Extrait d'offre : \emph{« Entreprise agroalimentaire recherche \textbf{Technicien Qualité H/F}. Missions : Contrôle matières premières/finis (normes ISO 22000/HACCP), Gestion non-conformités, Audits internes, Rédaction procédures. Profil : BTS/DUT QHSE ou Agro, \textbf{Maîtrise HACCP impérative}, Anglais technique lu/écrit, \textbf{Rigueur, Réactivité}. »} \\
Un candidat a un BTS Industries Alimentaires, stage en brasserie (contrôle hygrométrie, pH, nettoyage), stage en conserverie (HACCP : suivi CCP, traçabilité), niveau anglais B1. Son CV Master mentionne « Stage Brasserie : Contrôles qualité » et « Stage Conserverie : Suivi production ». Proposez la version adaptée des deux expériences.
\tcblower
\textbf{Solution détaillée :}

\textbf{Analyse mots-clés offre :} \cle{ISO 22000 / HACCP}, \cle{Non-conformités}, \cle{Audits}, \cle{Procédures}, \cle{Anglais technique}, \cle{Rigueur}.

\textbf{CV CIBLÉ -- Rubrique EXPÉRIENCES (Réordonnancée et réécrite) :}

\textbf{Fév -- Juin 2024 \hfill \emph{Stagiaire Assistant Qualité (HACCP)}} \\
\textbf{Conserverie « Saveurs du Cameroun » -- Douala} \hfill \emph{(4 mois)}
\begin{itemize}
    \item \textbf{Surveillé} les \textbf{Points Critiques de Maîtrise (CCP)} du plan \textbf{HACCP} (stérilisation, scellage) : relevés T°/Pression toutes les 30 min, \textbf{zéro dérive} sur la période.
    \item \textbf{Géré} \textbf{2 non-conformités} (étiquettes non conformes, scellage défectueux) : isolation lots, analyse causes (5M), actions correctives, traçabilité complète sous \textbf{GMAO}.
    \item \textbf{Participé} à la préparation de l'\textbf{audit de certification ISO 22000} : mise à jour des fiches de poste, vérification affichages, contrôle documents fournisseurs.
    \item \textbf{Rédigé} \textbf{3 modes opératoires} (Nettoyage/Désinfection ligne 2, Réception matières premières) validés par le Responsable Qualité.
\end{itemize}

\textbf{Sept -- Déc 2023 \hfill \emph{Stagiaire Contrôle Qualité / Laboratoire}} \\
\textbf{Brasseries du Cameroun -- Douala} \hfill \emph{(3 mois)}
\begin{itemize}
    \item \textbf{Effectué} les contrôles physico-chimiques matières premières (malt, houblon) et produits finis (bière) : \textbf{pH, densité, alcool, CO2, turbidité} selon normes internes (référentiel ISO 17025).
    \item \textbf{Assuré} la traçabilité des échantillons (étiquetage, conservation, registre) avec \textbf{rigueur} et respect des délais de rendu résultats (J+1).
    \item \textbf{Lu et exploité} des \textbf{fiches techniques en anglais} (fournisseurs malt/levures) pour valider la conformité des lots reçus. \textit{(Niveau B1 technique validé)}.
\end{itemize}

\textbf{Compétences TECHNIQUES (Remontées en tête) :}
\textbf{HACCP / ISO 22000} (Maîtrise opérationnelle CCP, Non-conformités, Audits) -- \textbf{Analyses labo agro} (pH, T°, Densité, Microbiologie bases) -- \textbf{GMAO / Traçabilité} -- \textbf{Anglais technique} (Lecture fiches fournisseurs, normes) -- \textbf{Rédaction procédures / Modes opératoires}.

\textbf{Conclusion :} Le candidat prouve qu'il \emph{parle le langage de l'offre}. Il a sélectionné, renommé, chiffré pour matcher exactement les exigences.
\end{exoresolu}

\courssub{Rédaction, présentation et évaluation du CV}

\begin{methode}[Rédiger un CV complet à partir d'une situation donnée (Intégration)]
    \textbf{Objectif :} Produire un document finalisé, cohérent, sans faute, prêt à être envoyé.
    \begin{enumerate}
        \item \textbf{Collecte des données :} Lister TOUT (diplômes, stages, jobs, bénévolat, compétences informatiques, langues, permis, centres d'intérêt, références).
        \item \textbf{Choix du modèle :} Sobre, une colonne (plus sûr pour ATS) ou deux colonnes (si design maîtrisé). Titre = Poste visé.
        \item \textbf{Remplissage rubrique par rubrique (Ordre antichronologique) :}
            \begin{itemize}
                \item \textbf{Formation :} Diplôme exact, Établissement, Ville, Année, \textbf{Mention/Note si $\ge$ 12/14 ou mention Bien/TB}. Options/spécialités. Mémoire/Projet de fin d'études (titre + 1 phrase).
                \item \textbf{Expériences :} Appliquer méthode STAR (Verbe + Contexte + Chiffre + Résultat). Inclure stages, jobs étudiants, bénévolat significatif (trésorier association, coach sportif).
                \item \textbf{Compétences :} Techniques (métier), Informatiques (logiciels métier + bureautique), Langues (Niveau CECRL + Score si dispo), Permis (B, C, Poids lourd...).
                \item \textbf{Intérêts :} 3 max. Précis (pas « Lecture » mais « Lecture : Romans africains, Histoire » ; pas « Sport » mais « Football : Milieu de terrain, Club X, 3 ans »).
            \end{itemize}
        \item \textbf{Relecture « Zéro Faute » (Checklist) :}
            \begin{itemize}
                \item Orthographe / Grammaire (Correcteur + Relecture à voix haute + Tierce personne).
                \item Cohérence dates (pas de chevauchements bizarres, trous expliqués).
                \item Coordonnées valides (Tester tel/email).
                \item Mots-clés de l'offre présents ?
                \item Fichier PDF nommé correctement.
            \end{itemize}
    \end{enumerate}
    \textbf{Astuce :} Faire relire par un prof / pro du secteur avant envoi.
\end{methode}

\begin{exoresolu}[Rédaction intégrée d'un CV (Cas pratique)]
\textbf{Énoncé :} \textbf{M. Ngono Jean}, 22 ans, Terminale Génie Civil (BTP). Diplômes : Probatoire Série F4 (2022, Mention AB), BEP BTP (2021). Stages : 1) Entreprise \texttt{RAZEL} (3 mois, 2023) : Suivi chantier route, lecture plans, béton, géomètre aide. 2) Mairie de sa commune (1 mois, 2022) : Suivi travaux voirie. Compétences : AutoCAD (notions), Covadis (notions), Word/Excel (bon), Lecture plans (bon), Béton (maîtrise dosage). Langues : Français (Langue maternelle), Anglais (A2 scolaire), Bulu (courant). Permis B (en cours). Centres d'intérêt : Foot (club local, capitaine), Dessin technique (maquettes ponts). Objectif : Stage de fin d'études / Premier emploi \textbf{Technicien Chantier / Géomètre Topographe}. Rédigez son CV complet (structure texte).
\tcblower
\textbf{Solution détaillée :}

\begin{center}
\textbf{\LARGE NGONO \textsc{Jean}} \\[0.5em]
\textbf{Technicien Chantier / Géomètre Topographe (Débutant)} \\[0.5em]
\begin{tabular}{ll}
📍 Douala, Cameroun & 📞 +237 6XX XX XX XX \\
✉️ jean.ngono@email.pro & 🔗 linkedin.com/in/jean-ngono-btp
\end{tabular}
\end{center}

\textbf{PROFIL}
\emph{Futur titulaire du Baccalauréat Technique Génie Civil (Juin 2024), rigoureux et terrain. Expérience significative en suivi de travaux routiers (RAZEL) et voirie communale. Maîtrise du dosage béton, lecture plans EXE/DAO (AutoCAD/Covadis notions). Permis B en cours. Mobilité nationale. Recherche stage de fin d'études (3-6 mois) ou premier poste Technicien Chantier.}

\textbf{FORMATION}
\begin{itemize}
    \item \textbf{2023 -- 2024 : Baccalauréat Technique Génie Civil (BTP)} -- Lycée Technique de [Ville]. \textit{En cours. Projet de fin d'études : « Étude d'un pont en béton armé sur la rivière X ».}
    \item \textbf{2021 -- 2022 : Probatoire Série F4 (Génie Civil)} -- Lycée Technique de [Ville]. \textbf{Mention Assez Bien (13,2/20).}
    \item \textbf{2020 -- 2021 : BEP Bâtiment Travaux Publics} -- CFA / Lycée Technique de [Ville]. \textbf{Admis.}
\end{itemize}

\textbf{EXPÉRIENCES PROFESSIONNELLES \& STAGES}
\begin{itemize}
    \item \textbf{Janv -- Mars 2023 (3 mois) \hfill \emph{Stagiaire Technicien Chantier}} \\
    \textbf{RAZEL -- Chantier Route Nationale 3 (Section Y)} \\
    \begin{itemize}
        \item \textbf{Assisté} le Chef de chantier dans le suivi quotidien : pointage personnel/engins, tenue journal de chantier.
        \item \textbf{Réalisé} les \textbf{contrôles réception béton} (affaissement Abrams, moulage éprouvettes 16cm, 3 séries/jour) -- \textbf{100\% conformes}.
        \item \textbf{Participé} aux \textbf{relevés topographiques} (station totale) pour implantation bordures/caniveaux -- Précision $\pm$ 2cm.
        \item \textbf{Lecture} et exploitation de \textbf{plans d'exécution (AutoCAD papier/numérique)} : profil en long, sections types, détails ferraillage.
        \item \textbf{Appliqué} les consignes S\&S (Port EPI, balisage zone travail, sensibilisation équipes).
    \end{itemize}
    \item \textbf{Juil 2022 (1 mois) \hfill \emph{Stagiaire Agent Technique Voirie}} \\
    \textbf{Mairie de [Commune] -- Service Voirie} \\
    \begin{itemize}
        \item \textbf{Effectué} des relevés de dégradés chaussées (nids-de-poule, ornières) pour programme d'entretien.
        \item \textbf{Surveillé} la reprise de tranchées (remblais, compactage) par entreprises sous-traitantes.
    \end{itemize}
\end{itemize}

\textbf{COMPÉTENCES TECHNIQUES}
\begin{tabular}{ll}
\textbf{Chantier / Béton :} & Dosage/Contrôle béton (normes), Coffrage/Ferraillage (lecture plans), Matériaux (granulats, ciments). \\
\textbf{Topographie / DAO :} & Station totale (notions implantation), \textbf{AutoCAD 2D} (lecture, calques, cotation), \textbf{Covadis} (notions profil en long). \\
\textbf{Bureautique :} & Word/Excel (Suivi tableaux de bord, planning simple), PowerPoint.
\end{tabular}

\textbf{LANGUES}
\begin{itemize}
    \item Français : Langue maternelle (Rédaction rapports techniques).
    \item Anglais : Niveau \textbf{A2} (Scolaire -- Lecture plans/notes techniques basiques).
    \item Bulu : Courant (Oral).
\end{itemize}

\textbf{PERMIS \& MOBILITÉ}
Permis B (En cours -- Examen prévu Juillet 2024) -- Véhicule personnel / Mobilité nationale totale.

\textbf{CENTRES D'INTÉRÊT}
\begin{itemize}
    \item \textbf{Football :} Capitaine équipe senior Club [Nom], 3 ans (Esprit d'équipe, Leadership, Gestion stress).
    \item \textbf{Dessin technique / Maquettisme :} Conception maquettes ponts/poutres bois (Précision, Spatialisation, Passion métier).
\end{itemize}

\textbf{RÉFÉRENCES}
Disponibles sur demande (M. [Nom], Chef Chantier RAZEL / M. [Nom], Chef Service Voirie Mairie).

\vspace{1em}
\textbf{Conclusion :} Ce CV respecte la structure (Profil $\to$ Formation $\to$ Expériences $\to$ Compétences $\to$ Langues/Permis $\to$ Intérêts), utilise l'antichronologie, la méthode STAR (verbes + chiffres), le vocabulaire métier (Abrams, éprouvettes, station totale, profil en long) et cible le poste visé.
\end{exoresolu}

\begin{methode}[Présenter et défendre son CV à l'oral (Entretien / Soutenance)]
    \textbf{Objectif :} Savoir présenter son parcours en 2-3 minutes (Pitch) et répondre aux questions sur le CV.
    \begin{enumerate}
        \item \textbf{Structure du Pitch (2-3 min) :}
            \begin{enumerate}
                \item \textbf{Accroche (15s) :} « Bonjour, je m'appelle [Prénom Nom], j'ai [âge] ans, je suis en Terminale [Spécialité] au Lycée [Nom], je vise un poste de [Titre]. »
                \item \textbf{Le fil conducteur (1min) :} « Mon parcours s'articule autour de [Compétence clé 1] et [Compétence clé 2]. » Parcourir Formation $\to$ Expériences majeures en \textbf{liant} les étapes (pas de liste à la Prévert). « Ce stage m'a permis de... », « C'est pour cela que j'ai choisi cette option... ».
                \item \textbf{La valeur ajoutée (45s) :} « Ma force aujourd'hui, c'est [Compétence technique pointue / Soft skill prouvé + Exemple concret]. Je suis aussi [Qualité humaine]. »
                \item \textbf{La projection (30s) :} « Rejoindre votre entreprise [Nom] me permettrait de [mettre en œuvre / développer]... sur [projet/type de chantier]. Je suis disponible dès le [Date]. »
            \end{enumerate}
        \item \textbf{Attitude :} Regard franc, voix posée, sourire, pas de lecture (notes autorisées sur carte bristol mots-clés uniquement).
        \item \textbf{Gestion des questions pièges sur le CV :}
            \begin{itemize}
                \item \emph{Trou dans le parcours :} Honnêteté + Valorisation (formation en ligne, bénévolat, projet perso, recherche active).
                \item \emph{Échec / Note basse :} « J'ai eu des difficultés en [Matière], j'ai réagi en [Action], aujourd'hui je suis à niveau / j'ai compensé par [Autre force]. »
                \item \emph{Expérience sans lien :} « Cela m'a appris la rigueur / le relationnel / la gestion du temps, transférable à ce poste. »
            \end{itemize}
    \end{enumerate}
    \textbf{Réflexe :} S'entraîner à chronométrer (chronomètre) et s'enregistrer (smartphone).
\end{methode}

\begin{exoresolu}[Simulation de pitch oral (2 min 30)]
\textbf{Énoncé :} À partir du CV de M. Ngono Jean (exercice précédent), rédigez le script exact de sa présentation orale de 2 minutes 30 devant un jury de recrutement (Chef de projet BTP + RH). Respectez la structure en 4 temps de la méthode.
\tcblower
\textbf{Solution détaillée :}

\textbf{SCRIPT DE PRÉSENTATION -- NGONO Jean (2 min 30)}

\textbf{1. ACCROCHE (15 secondes)}
« Bonjour Monsieur, Madame. Je m'appelle \textbf{Jean Ngono}, j'ai 22 ans. Je suis actuellement en \textbf{Terminale Génie Civil (BTP)} au Lycée Technique de [Ville]. Je me présente aujourd'hui pour postuler à un \textbf{stage de fin d'études de 6 mois} (ou un premier emploi) en tant que \textbf{Technicien Chantier / Géomètre Topographe} au sein de votre structure. »

\textbf{2. PARCOURS \& FIL CONDUCTEUR (1 minute)}
« Mon parcours est construit sur deux piliers : \textbf{la technique terrain} et \textbf{la rigueur du suivi}. \\
Après un BEP BTP obtenu brillamment, j'ai validé mon Probatoire Série F4 avec \textbf{Mention Assez Bien (13,2)}. Ce qui m'a tout de suite passionné, c'est la \textbf{réalité du chantier}. \\
\textbf{Mon stage majeur chez RAZEL (3 mois sur la RN3)} a été le déclic : j'ai quitté la théorie pour le concret. J'ai \textbf{suivi quotidiennement la production béton} (contrôles Abrams, éprouvettes), j'ai \textbf{manipulé la station totale} pour l'implantation de bordures, et j'ai \textbf{lu des plans d'exécution complexes} (profils en long, détails ferraillage) sous AutoCAD. \\
\textbf{Mon stage à la Mairie} m'a apporté la vision \textbf{maîtrise d'ouvrage} : relevés de dégradés, contrôle sous-traitants. \\
Aujourd'hui, mon Projet de Fin d'Études sur \textbf{l'étude d'un pont en béton armé} me permet d'approfondir le calcul de structure et le DAO (Covadis). »

\textbf{3. VALEUR AJOUTÉE / FORCES (45 secondes)}
« Ce qui me différencie, c'est ma \textbf{double compétence "Terrain + DAO"}. Je ne suis pas seulement capable de lire un plan, je sais \textbf{comment il se construit} (dosage, coffrage, sécurité). \\
Je cite deux preuves : à RAZEL, \textbf{zéro non-conformité béton} sur mes 3 mois de suivi. Et ma capacité à \textbf{manager une équipe} : en tant que \textbf{capitaine de mon équipe de foot} depuis 3 ans, j'ai appris à fédérer, communiquer sous pression et tenir les délais -- des soft skills indispensables sur un chantier. \\
Je suis \textbf{opérationnel immédiatement} sur le suivi béton, l'implantation simple, la lecture plans et la tenue de journal de chantier. »

\textbf{4. PROJECTION / DISPONIBILITÉ (30 secondes)}
« Intégrer \textbf{[Nom de l'Entreprise]}, reconnue pour [citer un projet récent ou valeur de l'entreprise : ex: « vos ouvrages d'art sur l'autoroute Yaoundé-Douala »], serait pour moi l'opportunité idéale de \textbf{mettre ma rigueur au service de vos chantiers} tout en progressant aux côtés de vos chefs de projet expérimentés. \\
Je suis \textbf{disponible dès le 1er juillet} (fin des épreuves écrites), mobile sur l'ensemble du territoire, et mon Permis B sera validé cet été. \\
Je vous remercie de votre attention et reste à votre disposition pour toute question technique sur mes stages ou mon PFE. »

\textbf{Analyse :} Temps respecté ($\approx$ 150 mots/min $\to$ 375 mots = 2 min 30). Structure claire. Vocabulaire pro. Preuves chiffrées. Lien avec l'entreprise. Call to action final.
\end{exoresolu}

\begin{attention}[Les 5 erreurs qui coûtent des points au Bac (Épreuve de Production Écrite / Intégration)]
    \begin{enumerate}
        \item \textbf{CV « Fourre-tout » non ciblé :} Envoyer le même CV pour tous les postes. \textbf{Sanction :} Hors-sujet partiel (non-adaptation à la situation d'embauche). \textbf{Remède :} Analyser l'offre $\to$ Mots-clés $\to$ Réordonnancer/Reformuler.
        \item \textbf{Expériences « Muettes » (Sans verbes d'action ni chiffres) :} Ex: « Stage : Maçonnerie ». \textbf{Sanction :} Preuve de compétence absente, note faible en « Structure interne / Contenu ». \textbf{Remède :} Méthode STAR (Verbe + Contexte + Chiffre + Résultat).
        \item \textbf{Fautes d'orthographe / grammaire / syntaxe :} Notamment confusion \textbf{Conditionnel / Futur} (\emph{je souhaiterai} au lieu de \emph{je souhaiterais}), accords participes passés, homophones (\emph{a/à, et/est, ou/où}). \textbf{Sanction :} Pénalité lourde en « Langue française » (souvent 2 à 4 points sur 20). \textbf{Remède :} Relecture à voix haute + Correcteur + Tierce personne.
        \item \textbf{Incohérence dates / Trous inexpliqués :} Chevauchements impossibles, année blanche non justifiée. \textbf{Sanction :} Perte de crédibilité, note « Structure externe / Cohérence » basse. \textbf{Remède :} Vérifier la chronologie, expliquer les trous (formation, projet perso, recherche active) en une ligne.
        \item \textbf{Présentation soignée mais contenu vide / Mensonger :} Belle mise en page (Canva) mais compétences inventées, stages fictifs. \textbf{Sanction :} Élimination morale + note 0 si fraude avérée ; note basse « Adéquation profil/poste » si décalage flagrant. \textbf{Remède :} Honnêteté, valoriser le peu (bénévolat, projets scolaires, soft skills) plutôt qu'inventer.
    \end{enumerate}
    \textbf{Rappel ultime :} Au Bac, l'épreuve d'intégration demande souvent \textbf{CV + Lettre de motivation}. Ne négligez pas la lettre : elle doit \textbf{argumenter} ce que le CV \textbf{liste}.
\end{attention}

\newpage

\coursec{Exercices}

\courssub{Je vérifie mes connaissances}

\begin{exercice}[QCM : les rubriques du CV]
Pour chaque question, une seule réponse est correcte. Entoure la lettre correspondante.

\textbf{1.} La rubrique « État civil » d'un CV contient généralement :
\begin{itemize}
\item[a)] les diplômes obtenus ;
\item[b)] le nom, le prénom, la date de naissance, les coordonnées ;
\item[c)] les loisirs du candidat.
\end{itemize}

\textbf{2.} Dans la rubrique « Expériences professionnelles », l'ordre recommandé est :
\begin{itemize}
\item[a)] chronologique (du plus ancien au plus récent) ;
\item[b)] alphabétique ;
\item[c)] antichronologique (du plus récent au plus ancien).
\end{itemize}

\textbf{3.} La phrase « Je souhaiterais obtenir ce poste » emploie :
\begin{itemize}
\item[a)] le conditionnel de politesse ;
\item[b)] le subjonctif de souhait ;
\item[c)] l'impératif.
\end{itemize}

\textbf{4.} Un texte iconique est :
\begin{itemize}
\item[a)] un texte composé uniquement de phrases ;
\item[b)] un document où l'image (photo, pictogramme, mise en page) participe au sens ;
\item[c)] un texte rédigé en vers.
\end{itemize}
\end{exercice}

\begin{exercice}[Vrai ou faux]
Indique si chaque affirmation est vraie ou fausse, puis justifie ta réponse en une phrase.

\begin{enumerate}
\item Un CV peut comporter trois pages pour un jeune diplômé.
\item La rubrique « Formation » présente les diplômes du plus récent au plus ancien.
\item « Je veux absolument ce poste ! » est une formulation adaptée à un dossier de candidature.
\item La photo est obligatoire sur tout CV déposé au Cameroun.
\item Le subjonctif peut exprimer le souhait après « je désire que ».
\end{enumerate}
\end{exercice}

\begin{exercice}[Définitions]
Définis en deux ou trois phrases chacun des termes suivants, tels qu'ils s'appliquent au CV :
\begin{enumerate}
\item structure externe ;
\item structure interne ;
\item rubrique ;
\item ordre antichronologique ;
\item texte fonctionnel.
\end{enumerate}
\end{exercice}

\begin{exercice}[Modes verbaux]
Pour chaque phrase, indique le mode du verbe souligné (indicatif, subjonctif, conditionnel, impératif) et précise sa valeur (politesse, souhait, hypothèse, ordre, fait réel).

\begin{enumerate}
\item Je \emph{voudrais} mettre mes compétences au service de votre entreprise.
\item Il faut que je \emph{sois} disponible dès le mois de septembre.
\item J'\emph{ai effectué} un stage de trois mois à la SABC de Douala.
\item \emph{Contactez}-moi pour tout renseignement complémentaire.
\item Si j'\emph{obtenais} ce poste, je m'investirais pleinement.
\end{enumerate}
\end{exercice}

\courssub{Je m'entraîne}

\begin{exercice}[Remettre en ordre un CV]
Voici les rubriques éparpillées du CV de Manga Solange, jeune diplômée d'un Baccalauréat technique en secrétariat, candidate à un poste d'assistante de direction à Yaoundé :

\begin{center}
\begin{tabular}{|c|l|}
\hline
Lettre & Rubrique \\
\hline
A & Centres d'intérêt : lecture, chant choral \\
B & Formation : Baccalauréat G2 (2024), BEPC (2020) \\
C & État civil : nom, prénom, née le 12/03/2005 à Bafoussam, téléphone, courriel \\
D & Expérience professionnelle : stage de 2 mois, mairie de Yaoundé 5 (2024) \\
E & Compétences : dactylographie, traitement de texte, français et anglais \\
F & Références : M. Nkoulou, chef du personnel, mairie de Yaoundé 5 \\
\hline
\end{tabular}
\end{center}

\begin{enumerate}
\item Remets ces rubriques dans l'ordre convenable pour le poste visé.
\item Justifie ton choix en deux ou trois phrases : pourquoi placer telle rubrique avant telle autre ?
\end{enumerate}
\end{exercice}

\begin{exercice}[Réécrire des phrases maladroites]
Les phrases suivantes, extraites de lettres et de CV de candidats, sont trop directes. Réécris-les en employant le conditionnel de politesse ou le subjonctif du souhait.

\begin{enumerate}
\item Je veux ce poste de comptable.
\item Envoyez-moi une convocation pour l'entretien.
\item Je suis sûr que vous allez me recruter.
\item Donnez-moi une chance de prouver mes capacités.
\item Je demande que tu m'appelles vite. (adressé au directeur)
\end{enumerate}
\end{exercice}

\begin{exercice}[Corriger un CV fautif]
Le CV ci-dessous contient \textbf{dix erreurs} (orthographe, dates incohérentes, rubrique manquante, mode verbal fautif, information inutile). Releve chaque erreur et propose la correction.

\begin{center}
\begin{tabularx}{\linewidth}{|l|X|}
\hline
\multicolumn{2}{|c|}{\textbf{CURRICULUM VITAE}} \\
\hline
État civil & Fotso Arnaud, né le 30 février 2004 à Douala, célibataire, taille : 1 m 75 \\
\hline
Formation & BEPC en 2023 ; Probatoire F3 en 2022 ; Baccalauréat F3 en 2021 \\
\hline
Expérience & Stage à la CAMTEL, du 1\ier{} juin 2024 au 30 mai 2024 \\
\hline
Compétences & Je veut maîtriser l'électricité bâtiment ; je exige un bon salaire \\
\hline
Centres d'intérêt & Football, musique, je déteste lire \\
\hline
\end{tabularx}
\end{center}
\end{exercice}

\begin{exercice}[Analyser deux CV iconiques]
On présente à la classe deux CV fictifs pour un même poste de technicien en maintenance :
\begin{itemize}
\item \textbf{CV 1} : une page, sans photo, rubriques séparées par des filets, titres en gras, pictogrammes (téléphone, enveloppe) devant les coordonnées ;
\item \textbf{CV 2} : deux pages, photo d'identité, texte compact sans titres visibles, plusieurs polices de caractères mélangées.
\end{itemize}

\begin{enumerate}
\item Releve trois éléments iconiques (photo, pictogrammes, couleurs, mise en page) présents dans chaque CV.
\item Rédige un paragraphe argumenté (8 à 12 lignes) dans lequel tu montres lequel des deux CV est le plus efficace et pourquoi.
\end{enumerate}
\end{exercice}

\courssub{Je me prépare au Bac}

\begin{exercice}[Production de texte fonctionnel : rédiger un CV adapté à une offre]
\textbf{Situation.} Tu lis dans un quotidien camerounais l'offre d'emploi suivante :

\begin{center}
\begin{tabularx}{\linewidth}{|X|}
\hline
\textbf{AVIS DE RECRUTEMENT} \\
La société AGROCAM S.A., située à Bafoussam, recrute \textbf{un(e) technicien(ne) en agroalimentaire}. \\
Profil : Baccalauréat technique (série F ou équivalent) ; sens de l'organisation ; maîtrise de l'outil informatique souhaitée ; une première expérience (stage accepté) serait un atout. \\
Dossier : CV et lettre de motivation à déposer au plus tard le 30 septembre. \\
\hline
\end{tabularx}
\end{center}

\textbf{Travail à faire.}
\begin{enumerate}
\item Relève dans l'offre les trois qualités ou compétences exigées que ton CV devra mettre en valeur.
\item Rédige ton CV complet (une page) en respectant la structure externe (titre, rubriques repérables, mise en page soignée) et la structure interne (état civil, formation et expériences en ordre antichronologique, compétences, centres d'intérêt, références).
\item Soigne particulièrement la rubrique « Compétences » : rédige-la en deux ou trois phrases employant au moins une fois le conditionnel de politesse et une fois le subjonctif du souhait.
\item Justifie en cinq lignes l'ordre dans lequel tu as placé tes rubriques, en fonction du poste visé.
\end{enumerate}
\end{exercice}

\begin{exercice}[Correction et adaptation d'un CV]
\textbf{Situation.} Ton ami Etoundi Jean, titulaire d'un Baccalauréat F4 (Génie civil), te montre le CV qu'il s'apprête à envoyer pour un poste de conducteur de travaux dans une entreprise de BTP à Douala. Voici des extraits de son document :

\begin{itemize}
\item « Formation : Baccalauréat F4, 2020 ; Probatoire F4, 2023. »
\item « Expérience : chantier de l'immeuble Émeraude, Douala, du 10 janvier 2024 au 5 janvier 2024. »
\item « Compétences : je veux que vous me donniez ce poste parce que je suis le meilleur. »
\item Rubrique « État civil » absente ; aucune coordonnée téléphonique.
\end{itemize}

\textbf{Travail à faire.}
\begin{enumerate}
\item Releve et corrige les erreurs de dates et de cohérence.
\item Réécris la phrase de la rubrique « Compétences » en deux phrases correctes, l'une au conditionnel de politesse, l'autre avec un subjonctif de souhait.
\item Reconstitue la rubrique « État civil » complète qu'Etoundi doit ajouter.
\item Rédige la version corrigée et complète du CV d'Etoundi, adaptée au poste visé.
\end{enumerate}
\end{exercice}

\begin{exercice}[Activité d'intégration : rédiger et présenter son CV]
\textbf{Situation d'intégration.} À la fin de l'année scolaire, le club « Avenir » de ton lycée organise un forum de l'emploi. Plusieurs entreprises de ta localité (menuiserie, garage, cybercafé, microfinance) proposent des stages et des emplois de vacances. Tu décides de postuler à l'une de ces offres.

\textbf{Travail à faire.}
\begin{enumerate}
\item Choisis l'une des quatre entreprises et précise le poste visé.
\item Rédige ton CV personnel adapté à ce poste : structure externe soignée (une page, rubriques claires, titre), structure interne complète, ordre antichronologique respecté, au moins un conditionnel de politesse et un subjonctif de souhait correctement employés.
\item Prépare une présentation orale de deux minutes de ton CV devant la classe (le jury) : annonce ton identité, ton diplôme, ta compétence principale et ta motivation.
\item Le jury te posera trois questions (par exemple : « Pourquoi avez-vous placé cette rubrique en premier ? », « Pourquoi n'avez-vous pas mis de photo ? », « Quelle expérience vous rend apte à ce poste ? »). Prépare par écrit, en trois à cinq lignes chacune, tes réponses justifiant tes choix de rédaction.
\end{enumerate}
\end{exercice}

\newpage

\coursec{Corrigés des exercices}

\coursec{Corrigés des activités — Leçon VI : Rédiger un Curriculum Vitae}

\begin{corrige}[Exercice 1 — QCM : les rubriques du CV]
\textbf{1. Réponse b).} La rubrique « État civil » regroupe les informations d'identification du candidat : nom, prénom, date et lieu de naissance, situation familiale éventuelle, adresse, téléphone et courriel. Les diplômes relèvent de la rubrique « Formation » et les loisirs de la rubrique « Centres d'intérêt ».

\textbf{2. Réponse c).} Dans un CV, les expériences professionnelles (comme les diplômes) se présentent en \cle{ordre antichronologique}, c'est-à-dire du plus récent au plus ancien : le recruteur s'intéresse d'abord à ce que le candidat fait ou vient de faire.

\textbf{3. Réponse a).} « Je souhaiterais » est un \cle{conditionnel de politesse} (conditionnel présent du verbe \emph{souhaiter}) : il atténue la demande et marque le respect envers le destinataire. Ce n'est pas un subjonctif (le verbe n'est pas précédé de « que ») ni un impératif.

\textbf{4. Réponse b).} Un \cle{texte iconique} est un document où l'image — photo, pictogramme, couleur, mise en page — participe à la construction du sens, à côté du texte écrit. Le CV est un exemple typique de texte fonctionnel à dimension iconique.
\end{corrige}

\begin{corrige}[Exercice 2 — Vrai ou faux]
\begin{enumerate}
\item \textbf{Faux.} Un CV de jeune diplômé doit tenir sur \textbf{une seule page} : le recruteur consacre peu de temps à chaque dossier et une information diluée perd en efficacité.
\item \textbf{Vrai.} La rubrique « Formation » suit l'ordre antichronologique : le diplôme le plus récent (souvent le plus élevé) est placé en premier car c'est celui qui intéresse le plus le recruteur.
\item \textbf{Faux.} « Je veux absolument ce poste ! » est une formulation trop directe, impérative et familière ; dans un dossier de candidature, on emploie le conditionnel de politesse : « Je souhaiterais obtenir ce poste ».
\item \textbf{Faux.} La photo n'est pas obligatoire sur un CV au Cameroun ; elle est facultative et, si on la met, il faut choisir une photo d'identité sérieuse et récente.
\item \textbf{Vrai.} Après un verbe de volonté ou de souhait comme « je désire que », « je souhaite que », « je veux que », la subordonnée exige le \cle{subjonctif} : « Je désire que vous \emph{soyez} présent ».
\end{enumerate}
\end{corrige}

\begin{corrige}[Exercice 3 — Définitions]
\begin{enumerate}
\item \textbf{Structure externe} : c'est l'organisation matérielle et visuelle du CV : le titre (« Curriculum Vitae »), les rubriques clairement séparées et repérables, les titres en gras, les filets ou encadrés, l'aération de la page. Elle permet au recruteur de trouver l'information en un coup d'œil.
\item \textbf{Structure interne} : c'est le contenu logique du CV, c'est-à-dire l'ensemble des informations réparties dans les rubriques : état civil, formation, expériences professionnelles, compétences, centres d'intérêt, références. Elle obéit à des règles d'ordre (notamment l'ordre antichronologique).
\item \textbf{Rubrique} : chacune des parties titrées du CV, regroupant des informations de même nature (Formation, Expérience professionnelle, Compétences, etc.). Chaque rubrique porte un titre visible qui la signale.
\item \textbf{Ordre antichronologique} : mode de classement qui présente les éléments du plus récent au plus ancien. Dans un CV, il s'applique aux rubriques « Formation » et « Expériences professionnelles » afin de mettre en avant la situation actuelle du candidat.
\item \textbf{Texte fonctionnel} : texte qui vise une utilité pratique immédiate dans la vie sociale ou professionnelle (lettre, CV, facture, formulaire, compte rendu). Il se reconnaît à sa présentation codifiée, à son vocabulaire précis et à son objectif d'information ou d'action.
\end{enumerate}
\end{corrige}

\begin{corrige}[Exercice 4 — Modes verbaux]
\begin{enumerate}
\item « Je \emph{voudrais} mettre mes compétences\ldots » : \textbf{conditionnel présent}, valeur de \textbf{politesse} (demande atténuée adressée au recruteur).
\item « Il faut que je \emph{sois} disponible\ldots » : \textbf{subjonctif présent} (après « il faut que »), valeur d'\textbf{obligation} ; on retiendra que le subjonctif exprime ici ce qui est exigé, non un fait réel.
\item « J'\emph{ai effectué} un stage de trois mois\ldots » : \textbf{indicatif passé composé}, valeur de \textbf{fait réel} et accompli : le stage a réellement eu lieu.
\item « \emph{Contactez}-moi\ldots » : \textbf{impératif présent}, valeur d'\textbf{ordre} ou d'invitation ; c'est une formule de conclusion courante et acceptable dans un dossier de candidature.
\item « Si j'\emph{obtenais} ce poste, je m'investirais pleinement » : \textbf{indicatif imparfait} après « si », valeur d'\textbf{hypothèse} (condition), en corrélation avec le conditionnel « je m'investirais » dans la proposition principale.
\end{enumerate}
\end{corrige}

\begin{attention}[Point de vigilance — Exercice 4]
Après « si » exprimant une hypothèse, on n'emploie jamais le conditionnel : on dit « si j'obtenais » (imparfait) et non « si j'obtiendrais ». Le conditionnel se place dans l'autre proposition.
\end{attention}

\begin{corrige}[Exercice 5 — Remettre en ordre un CV]
\textbf{1. Ordre convenable : C — B — D — E — A — F.}

\begin{center}
\begin{tabular}{|c|l|}
\hline
Ordre & Rubrique \\
\hline
1 & C — État civil (identification du candidat) \\
2 & B — Formation (Baccalauréat G2, 2024 ; BEPC, 2020) \\
3 & D — Expérience professionnelle (stage, mairie de Yaoundé 5, 2024) \\
4 & E — Compétences (dactylographie, traitement de texte, langues) \\
5 & A — Centres d'intérêt (lecture, chant choral) \\
6 & F — Références (M. Nkoulou, chef du personnel) \\
\hline
\end{tabular}
\end{center}

\textbf{2. Justification.} L'état civil ouvre le CV car le recruteur doit d'abord savoir qui est le candidat et comment le joindre. Pour une jeune diplômée sans longue expérience, la rubrique « Formation » vient juste après : le Baccalauréat G2 est son principal atout pour un poste d'assistante de direction. L'expérience (stage) et les compétences (dactylographie, bilinguisme), directement liées au poste, suivent immédiatement. Les centres d'intérêt et les références, informations secondaires et complémentaires, ferment le document.
\end{corrige}

\begin{corrige}[Exercice 6 — Réécrire des phrases maladroites]
\begin{enumerate}
\item « Je veux ce poste de comptable. » $\rightarrow$ « Je \emph{souhaiterais} obtenir le poste de comptable offert par votre entreprise. » (conditionnel de politesse)
\item « Envoyez-moi une convocation pour l'entretien. » $\rightarrow$ « Je vous \emph{serais} reconnaissant de bien vouloir m'adresser une convocation pour l'entretien. » (conditionnel de politesse)
\item « Je suis sûr que vous allez me recruter. » $\rightarrow$ « Je \emph{souhaite} vivement que ma candidature \emph{retienne} votre attention. » (subjonctif du souhait)
\item « Donnez-moi une chance de prouver mes capacités. » $\rightarrow$ « Je \emph{désirerais} qu'on me \emph{donne} l'occasion de prouver mes capacités. » (conditionnel + subjonctif)
\item « Je demande que tu m'appelles vite. » $\rightarrow$ « Je \emph{souhaiterais} que vous \emph{m'appeliez} dès que possible, Monsieur le Directeur. » (conditionnel de politesse, subjonctif, vouvoiement obligatoire)
\end{enumerate}
\end{corrige}

\begin{attention}[Point de vigilance — Exercice 6]
Deux corrections sont exigées dans la phrase 5 : le passage du tutoiement au vouvoiement (registre soutenu obligatoire envers un directeur) et le remplacement de l'indicatif direct par le conditionnel de politesse suivi du subjonctif.
\end{attention}

\begin{corrige}[Exercice 7 — Corriger un CV fautif]
\begin{center}
\begin{tabularx}{\linewidth}{|c|X|X|}
\hline
\rowcolor{popblueL} N° & Erreur relevée & Correction \\
\hline
1 & « né le 30 février 2004 » : date impossible & « né le 28 février 2004 » (février n'a jamais 30 jours) \\
\hline
2 & « taille : 1 m 75 » : information inutile & Supprimer (la taille n'intéresse pas le recruteur) \\
\hline
3 & « BEPC en 2023 ; Probatoire F3 en 2022 ; Baccalauréat F3 en 2021 » : ordre et dates incohérents & Ordre antichronologique et dates logiques : « Baccalauréat F3 (2023) ; Probatoire F3 (2022) ; BEPC (2020) » \\
\hline
4 & « du 1\ier{} juin 2024 au 30 mai 2024 » : la fin précède le début & « du 1\ier{} mai 2024 au 30 juin 2024 » (par exemple) \\
\hline
5 & « Je veut » : faute de conjugaison & « Je veux » \\
\hline
6 & « je exige » : élision manquante et ton arrogant & « j'exige » — mieux : supprimer cette phrase inadaptée \\
\hline
7 & « Je veux maîtriser » : mode inadapté (on n'exige pas dans un CV) & « Je souhaiterais perfectionner ma maîtrise de l'électricité bâtiment » (conditionnel de politesse) \\
\hline
8 & « je exige un bon salaire » : prétention déplacée dans un CV & Supprimer ; le salaire se discute à l'entretien, pas dans le CV \\
\hline
9 & « je déteste lire » : information négative et rédaction personnelle & « Football, musique » (on ne mentionne jamais ce qu'on déteste) \\
\hline
10 & Rubrique « Références » absente & Ajouter : « Références : M. X, chef de service, CAMTEL Douala » \\
\hline
\end{tabularx}
\end{center}
\end{corrige}

\begin{attention}[Point de vigilance — Exercice 7]
Dans un CV, on vérifie toujours la cohérence des dates (un diplôme ne peut précéder le diplôme inférieur ; une période ne peut se terminer avant d'avoir commencé) et on bannit toute information négative, inutile ou prétentieuse.
\end{attention}

\begin{corrige}[Exercice 8 — Analyser deux CV iconiques]
\textbf{1. Éléments iconiques relevés.}

\begin{center}
\begin{tabularx}{\linewidth}{|l|X|X|}
\hline
\rowcolor{popblueL} Élément & CV 1 & CV 2 \\
\hline
Photo & absente & photo d'identité présente \\
\hline
Pictogrammes & téléphone, enveloppe devant les coordonnées & aucun \\
\hline
Mise en page & rubriques séparées par des filets, titres en gras, une seule page, une seule police & texte compact sans titres visibles, deux pages, plusieurs polices mélangées \\
\hline
\end{tabularx}
\end{center}

\textbf{2. Paragraphe argumenté (corrigé type).}

Le CV 1 est nettement le plus efficace des deux. D'abord, sa structure externe est soignée : les rubriques, séparées par des filets et annoncées par des titres en gras, permettent au recruteur de repérer instantanément l'information cherchée, alors que le texte compact du CV 2, dépourvu de titres visibles, oblige à une lecture fastidieuse. Ensuite, les pictogrammes du CV 1 (téléphone, enveloppe) guident l'œil vers les coordonnées : l'image y participe pleinement au sens, ce qui est la marque d'un texte iconique réussi. En outre, le CV 1 tient sur une seule page, ce qui convient parfaitement à un jeune candidat, tandis que les deux pages du CV 2 diluent l'information. Enfin, l'unité de police du CV 1 traduit le sérieux et la rigueur du candidat, qualités attendues d'un technicien en maintenance, alors que le mélange de polices du CV 2 donne une impression de désordre. Certes, le CV 2 comporte une photo, mais cet avantage ne compense pas ses défauts de lisibilité. En conclusion, c'est la clarté, la concision et la cohérence graphique du CV 1 qui retiendront l'attention du recruteur.
\end{corrige}

\begin{corrige}[Exercice 9 — Production de texte fonctionnel : rédiger un CV adapté à une offre]
\textbf{1. Qualités et compétences exigées par l'offre :}
\begin{itemize}
\item le diplôme exigé : Baccalauréat technique (série F ou équivalent) ;
\item le sens de l'organisation ;
\item la maîtrise de l'outil informatique (et, comme atout, une première expérience, même un stage).
\end{itemize}

\textbf{2. Corrigé type de CV (une page).}

\begin{center}
\begin{tabularx}{\linewidth}{|l|X|}
\hline
\multicolumn{2}{|c|}{\textbf{CURRICULUM VITAE}} \\
\hline
État civil & Kamga Yvan, né le 15 avril 2005 à Bafoussam, célibataire. Tél. : 6 90 00 00 00 ; courriel : yvan.kamga@mail.com ; B.P. 123 Bafoussam \\
\hline
Formation & 2024 : Baccalauréat technique F5 (agroalimentaire), Lycée technique de Bafoussam ; 2023 : Probatoire F5 ; 2020 : BEPC \\
\hline
Expérience professionnelle & Juillet–septembre 2024 : stage d'ouvrier de production, laiterie de Bafoussam (contrôle qualité, conditionnement) \\
\hline
Compétences & Techniques de transformation et de conservation des produits agroalimentaires ; maîtrise de l'outil informatique (traitement de texte, tableur) ; sens de l'organisation acquis lors du suivi des chaînes de production. Je souhaiterais mettre ces compétences au service de votre société, et je désire que mon sens du travail en équipe soit pleinement utile à votre production. \\
\hline
Centres d'intérêt & Football, lecture, jardinage \\
\hline
Références & M. Tchoupo, responsable de production, laiterie de Bafoussam \\
\hline
\end{tabularx}
\end{center}

\textbf{3. Modes employés dans la rubrique « Compétences » :} « Je \emph{souhaiterais} mettre\ldots » (conditionnel de politesse) ; « je désire que mon sens du travail en équipe \emph{soit} pleinement utile » (subjonctif du souhait après un verbe de volonté).

\textbf{4. Justification de l'ordre des rubriques (5 lignes).} L'état civil ouvre le CV pour identifier le candidat et permettre à l'entreprise de le contacter. La formation vient ensuite car le diplôme exigé (Baccalauréat technique série F) est le premier critère de l'offre : il faut le montrer immédiatement. L'expérience, même brève, suit car l'offre la présente comme un atout. Les compétences répondent point par point au profil demandé (organisation, informatique). Centres d'intérêt et références, informations complémentaires, ferment logiquement le document.

\textbf{Barème indicatif (sur 20).} Repérage des trois exigences de l'offre : 3 pts ; structure externe (titre, rubriques repérables, une page) : 4 pts ; structure interne complète et ordre antichronologique : 5 pts ; conditionnel de politesse et subjonctif correctement employés : 3 pts ; correction de la langue (orthographe, conjugaison) : 3 pts ; justification de l'ordre des rubriques : 2 pts.
\end{corrige}

\begin{attention}[Points de vigilance — Exercice 9]
Le correcteur vérifiera : la présence de toutes les rubriques obligatoires ; l'ordre antichronologique dans « Formation » et « Expérience » ; l'adéquation du contenu au profil de l'offre (série F, informatique, organisation) ; l'emploi correct du subjonctif après « je désire que » et du conditionnel de politesse ; l'absence de fautes dans les dates et les noms propres.
\end{attention}

\begin{corrige}[Exercice 10 — Correction et adaptation d'un CV]
\textbf{1. Erreurs de dates et de cohérence.}
\begin{itemize}
\item « Baccalauréat F4, 2020 ; Probatoire F4, 2023 » : double erreur. Le Probatoire se passe \emph{avant} le Baccalauréat, et l'ordre antichronologique impose de citer le diplôme le plus récent d'abord. Correction : « Baccalauréat F4 (2023) ; Probatoire F4 (2022) ».
\item « du 10 janvier 2024 au 5 janvier 2024 » : la date de fin précède la date de début, ce qui est impossible. Correction : « du 10 janvier 2024 au 5 juillet 2024 » (ou toute date postérieure au début).
\end{itemize}

\textbf{2. Réécriture de la rubrique « Compétences ».}

Phrase au conditionnel de politesse : « Je \emph{souhaiterais} mettre mes compétences en conduite de chantier au service de votre entreprise. »

Phrase avec subjonctif de souhait : « Je \emph{désire} que ma rigueur et mon sens de l'organisation \emph{soient} reconnus comme des atouts pour vos chantiers. »

\textbf{3. Rubrique « État civil » reconstituée.}

« Etoundi Jean, né le [date] à [lieu], célibataire. Téléphone : 6 XX XX XX XX ; courriel : etoundi.jean@mail.com ; adresse : B.P. XXXX Douala. » (Le téléphone est indispensable : sans lui, l'entreprise ne peut pas convoquer le candidat.)

\textbf{4. Version corrigée et complète du CV.}

\begin{center}
\begin{tabularx}{\linewidth}{|l|X|}
\hline
\multicolumn{2}{|c|}{\textbf{CURRICULUM VITAE}} \\
\hline
État civil & Etoundi Jean, né le 20 mai 2004 à Yaoundé, célibataire. Tél. : 6 95 00 00 00 ; courriel : etoundi.jean@mail.com ; B.P. 456 Douala \\
\hline
Formation & 2023 : Baccalauréat F4 (Génie civil), Lycée technique de Douala ; 2022 : Probatoire F4 ; 2019 : BEPC \\
\hline
Expérience professionnelle & Du 10 janvier 2024 au 5 juillet 2024 : stage sur le chantier de l'immeuble Émeraude, Douala (suivi des travaux, lecture de plans) \\
\hline
Compétences & Lecture de plans, métré, organisation d'équipe. Je souhaiterais mettre mes compétences en conduite de chantier au service de votre entreprise ; je désire que ma rigueur soit un atout pour vos travaux. \\
\hline
Centres d'intérêt & Football, dessin technique \\
\hline
Références & M. Abega, chef de chantier, immeuble Émeraude, Douala \\
\hline
\end{tabularx}
\end{center}

\textbf{Barème indicatif (sur 20).} Correction des deux incohérences de dates : 4 pts ; réécriture des deux phrases (conditionnel + subjonctif) : 4 pts ; état civil complet avec coordonnées : 4 pts ; CV final complet, structuré et adapté au poste : 6 pts ; correction de la langue : 2 pts.
\end{corrige}

\begin{attention}[Points de vigilance — Exercice 10]
La justification des corrections de dates doit être explicite (le Probatoire précède toujours le Baccalauréat ; une période ne peut finir avant son début). Le ton du CV doit rester modeste : « je suis le meilleur » est une prétention à bannir absolument.
\end{attention}

\begin{corrige}[Exercice 11 — Activité d'intégration : rédiger et présenter son CV]
\textbf{1. Choix de l'offre (exemple).} Entreprise choisie : le garage moderne de la localité ; poste visé : aide-mécanicien pour les vacances (stage rémunéré).

\textbf{2. Corrigé type de CV personnel.}

\begin{center}
\begin{tabularx}{\linewidth}{|l|X|}
\hline
\multicolumn{2}{|c|}{\textbf{CURRICULUM VITAE}} \\
\hline
État civil & Ngo Bell Franck, né le 3 novembre 2006 à Ebolowa, célibataire. Tél. : 6 99 00 00 00 ; courriel : franck.ngobell@mail.com \\
\hline
Formation & 2025 : classe de Terminale technique F2 (mécanique), Lycée technique d'Ebolowa ; 2024 : Probatoire F2 ; 2021 : BEPC \\
\hline
Expérience professionnelle & Juillet–août 2024 : stage d'aide-mécanicien, garage Central, Ebolowa (vidange, freinage, démontage) \\
\hline
Compétences & Entretien courant des véhicules, lecture de schémas mécaniques, ponctualité. Je souhaiterais approfondir ma pratique au sein de votre atelier, et je désire que mon sérieux soit utile à votre équipe. \\
\hline
Centres d'intérêt & Mécanique, football, musique \\
\hline
Références & M. Ondoa, patron du garage Central, Ebolowa \\
\hline
\end{tabularx}
\end{center}

\textbf{3. Présentation orale (plan de deux minutes).}
\begin{itemize}
\item \textbf{Identité (20 s)} : « Monsieur le Président, Mesdames et Messieurs les membres du jury, je me nomme Ngo Bell Franck. »
\item \textbf{Diplôme et formation (30 s)} : « Je prépare le Baccalauréat technique F2, spécialité mécanique, au Lycée technique d'Ebolowa, après avoir obtenu le Probatoire F2 en 2024. »
\item \textbf{Compétence principale et expérience (40 s)} : « J'ai effectué un stage d'aide-mécanicien au garage Central, où j'ai pratiqué la vidange, le freinage et le démontage. »
\item \textbf{Motivation (30 s)} : « Je souhaiterais mettre mon sérieux et ma ponctualité au service de votre atelier pendant les vacances. Je vous remercie de votre attention. »
\end{itemize}

\textbf{4. Réponses préparées aux questions du jury (corrigé type).}

\emph{« Pourquoi avez-vous placé cette rubrique en premier ? »} — J'ai placé l'état civil en premier parce que le recruteur doit d'abord identifier le candidat et pouvoir le contacter ; j'ai ensuite mis la formation avant l'expérience car, étant élève, mon diplôme préparé est mon principal atout.

\emph{« Pourquoi n'avez-vous pas mis de photo ? »} — La photo est facultative sur un CV ; je n'ai pas jugé utile d'en ajouter une, car c'est le contenu des rubriques qui doit convaincre. Si l'entreprise l'exige, je fournirai une photo d'identité récente et sérieuse.

\emph{« Quelle expérience vous rend apte à ce poste ? »} — Mon stage au garage Central m'a familiarisé avec les gestes de base de la mécanique (vidange, freinage, démontage) et avec la discipline d'un atelier ; cette première expérience, même brève, me rend immédiatement opérationnel comme aide-mécanicien.

\textbf{Barème indicatif (sur 20).} Choix clair du poste : 1 pt ; structure externe soignée (une page, titre, rubriques) : 4 pts ; structure interne complète et ordre antichronologique : 4 pts ; conditionnel de politesse et subjonctif corrects : 3 pts ; présentation orale (clarté, tenue, durée respectée) : 4 pts ; qualité des réponses justifiant les choix : 4 pts.
\end{corrige}

\begin{attention}[Points de vigilance — Exercice 11]
À l'oral : se lever, saluer le jury, parler distinctement sans lire le CV mot à mot, respecter les deux minutes. À l'écrit : vérifier chaque date, accorder les participes, employer le vouvoiement et éviter toute formulation familière ou prétentieuse. Toutes les réponses au jury doivent \emph{justifier} un choix de rédaction, jamais s'excuser.
\end{attention}

\newpage

\coursec{Fiche bilan}

\courssub{Carte mentale de la leçon}
\begin{popfigure}
\begin{tikzpicture}[
  mindmap,
  grow cyclic,
  every node/.style={concept, execute at begin node=\hskip0pt},
  concept color=popblue,
  level 1/.style={level distance=4.5cm, sibling angle=360/6, font=\large\bfseries},
  level 2/.style={level distance=3.2cm, sibling angle=360/6, font=\small},
  CV/.style={concept color=popblue},
  Ext/.style={concept color=poporange},
  Int/.style={concept color=popgreen},
  Lang/.style={concept color=poppurple},
  Icon/.style={concept color=poppink},
  Redac/.style={concept color=popgold}
]
\node[CV] {Le CV}
  child[Ext] { node[align=center]{Structure\\ Externe}
    child { node {Mise en page} }
    child { node {Police / Couleurs} }
    child { node {Format (PDF)} }
  }
  child[Int] { node[align=center]{Structure\\ Interne}
    child { node {Identité / Coord.} }
    child { node {Profil / Objectif} }
    child { node {Formation} }
    child { node {Expériences} }
    child { node {Compétences} }
    child { node {Centres d'intérêt} }
  }
  child[Lang] { node {Langue : Modes}
    child { node {Subjonctif : souhait, doute, nécessité} }
    child { node {Conditionnel : politesse, hypothèse, atténuation} }
  }
  child[Icon] { node[align=center]{Textes\\ Iconiques}
    child { node {Analyse offre emploi} }
    child { node {Mots-clés ciblés} }
    child { node {Adaptation CV} }
  }
  child[Redac] { node[align=center]{Rédaction\\ \& Présentation}
    child { node {Verbes d'action} }
    child { node {Chronologie inversée} }
    child { node {Relecture / Oral} }
  }
  child[Int] { node {Intégration}
    child { node {Cas pratique} }
    child { node {Simulation orale} }
  };
\end{tikzpicture}
\legende{Synthèse visuelle de la leçon VI : Le Curriculum Vitae (structure, langue, adaptation, rédaction).}
\end{popfigure}

\courssub{Bilan des savoirs et savoir-faire}
\begin{bilan}

\textbf{Définitions clés}
\begin{itemize}
  \item \cle{Curriculum Vitae (CV)} : Document synthétique présentant le parcours (formation, expériences, compétences) d'un candidat pour une candidature (emploi, stage, formation).
  \item \cle{Structure externe} : Aspect visuel (mise en page aérée, police lisible 11-12 pt, couleurs sobres, 1 page max pour junior, format PDF, photo pro optionnelle).
  \item \cle{Structure interne} : Rubriques obligatoires : \textbf{En-tête} (identité, coordonnées, titre visé), \textbf{Profil/Objectif} (3 lignes ciblées), \textbf{Formation} (diplômes, mentions, dates, anti-chronologique), \textbf{Expériences} (stages, jobs, missions, résultats chiffrés), \textbf{Compétences} (techniques, langues, informatiques, transversales), \textbf{Centres d'intérêt} (cohérence avec le poste).
  \item \cle{Adaptation} : Personnalisation du CV (mots-clés, ordre des rubriques, profil) en fonction de l'offre d'emploi analysée (texte iconique).
\end{itemize}

\textbf{Langue française : Valeurs des modes (rappel Bac)}
\begin{center}
\begin{tabularx}{\linewidth}{|l|X|X|}\hline
\rowcolor{popblueL} \textbf{Mode} & \textbf{Valeurs principales (contexte CV / Lettre)} & \textbf{Exemples CV / Lettre} \\ \hline
\textbf{Subjonctif} & \begin{itemize}\item Volonté, souhait, désir : \emph{Je souhaite que vous \textbf{receviez} ma candidature.}\item Nécessité, obligation : \emph{Il faut que je \textbf{maîtrise} l'anglais.}\item Doute, émotion : \emph{Je doute qu'il \textbf{soit} disponible.}\end{itemize} & \textit{Qu'il me soit permis de...}, \textit{Il est essentiel que vous \textbf{sachiez}...} \\ \hline
\textbf{Conditionnel} & \begin{itemize}\item Politesse, atténuation : \emph{Je \textbf{serais} honoré de vous rencontrer.}\item Hypothèse, condition : \emph{Si j'intègre votre équipe, je \textbf{m'investirais} pleinement.}\item Conseil : \emph{Vous \textbf{devriez} valoriser vos stages.}\end{itemize} & \textit{Je \textbf{voudrais} mettre mes compétences à profit...}, \textit{Cela \textbf{permettrait} de...} \\ \hline
\end{tabularx}
\end{center}

\textbf{Méthode express : Rédiger un CV efficace (5 étapes)}
\begin{enumerate}
  \item \textbf{Analyser l'offre} : Identifier mots-clés, compétences requises, valeurs de l'entreprise.
  \item \textbf{Choisir le titre} : Intitulé exact du poste visé (ex: \emph{Technicien en maintenance industrielle}).
  \item \textbf{Remplir les rubriques} (anti-chronologique) : Verbes d'action au passé composé ou infinitif (\emph{Géré, Conçu, Programmé, Vendu}). Chiffrer les résultats (\emph{+15\% CA, 50 clients/jour}).
  \item \textbf{Soigner la forme} : Alignements parfaits, puces régulières, pas de fautes, PDF nommé \texttt{Prenom\_Nom\_CV\_Poste.pdf}.
  \item \textbf{Relire et tester} : Lecture à voix haute, test ATS (pas de colonnes complexes, polices standards), avis tiers.
\end{enumerate}

\textbf{Pistes d'adaptation (Textes iconiques / Offre d'emploi)}
\begin{itemize}
  \item \cle{Mots-clés} : Reprendre termes exacts de l'offre (logiciels, certifications, soft skills).
  \item \cle{Ordre des rubriques} : Expérience avant Formation si profil senior ; Formation avant Expérience si junior.
  \item \cle{Profil} : Réécrire la phrase d'accroche pour matcher la fiche de poste.
  \item \cle{Compétences} : Mettre en avant celles demandées (ex: \emph{Autocad, Gestion de stock, Allemand B2}).
\end{itemize}

\end{bilan}

\courssub{Checklist avant le Bac / Probatoire}
\begin{aretenir}[Checklist CV — À cocher $\square$]
$\square$ Le CV tient sur \textbf{1 page} (format A4, PDF, nom de fichier pro).\par
$\square$ \textbf{En-tête complet} : Prénom NOM, Téléphone, Email pro, Ville, LinkedIn/Portfolio (liens cliquables), Titre visé.\par
$\square$ \textbf{Photo} : Professionnelle (fond neutre, tenue adaptée) ou absente (selon usage pays/secteur).\par
$\square$ \textbf{Profil / Objectif} : 2-3 lignes, personnalisé pour l'offre, verbes d'action, valeur ajoutée claire.\par
$\square$ \textbf{Formation} : Anti-chronologique, diplôme exact, établissement, année, mention, options pertinentes.\par
$\square$ \textbf{Expériences} : Anti-chronologique, Structure / Dates / Lieu / \textbf{3-5 puces} (Verbe action + Contexte + Résultat chiffré).\par
$\square$ \textbf{Compétences} : Classées (Techniques, Informatiques, Langues \textbf{niveau CECRL}, Permis, Soft skills), \textbf{pas de barre de progression}.\par
$\square$ \textbf{Centres d'intérêt} : Spécifiques (pas "lecture, sport"), montrent personnalité / soft skills (bénévolat, compétition, pratique artistique).\par
$\square$ \textbf{Zéro faute} : Relecture à l'envers, correcteur, impression test, accord participes passés, homophones.\par
$\square$ \textbf{Mots-clés de l'offre} : Présents naturellement dans Profil, Compétences, Expériences (passage ATS).\par
$\square$ \textbf{Présentation orale} : Pitch 1 min (Parcours $\to$ Motivation $\to$ Projet) prêt, posture, regard, voix.\par
$\square$ \textbf{Lettre de motivation} : Assortie (même en-tête, police, ton), structure \textbf{Vous / Moi / Nous}, subjonctif/conditionnel maîtrisés.
\end{aretenir}

\findecours

\end{document}
